\documentclass[12pt,twoside, openany]{book}

% IMPORTANT: Compile with XeLaTeX on Overleaf

% For book dimensions
\usepackage[
  paperwidth=6in,
  paperheight=9in,
  inner=0.8in,      % Margin near the spine
  outer=0.6in,      % Margin near the edge
  top=0.8in,
  bottom=0.8in,
  includeheadfoot   % Ensures headers/footers stay inside the margins
]{geometry}

% Prevent text overflow
\sloppy

\usepackage[strings]{underscore}

% Dotted leaders for chapters
\usepackage{tocloft}
\renewcommand{\cftdotsep}{1.5}
\renewcommand{\cftchapleader}{\cftdotfill{\cftdotsep}}

% \renewcommand{\cftchapfont}{\normalfont}
\renewcommand{\cftchappagefont}{\normalfont}

% Indentation settings
\cftsetindents{chapter}{0em}{2em}

% For no Κεφάλαιο Ν before the title
\usepackage{titlesec}

\titleformat{\chapter}
  {\normalfont\LARGE\filcenter} 
  {\thechapter}
  {1em}
  {}

\titlespacing*{\chapter}
  {0pt}   % left margin
  {1cm}   % space before chapter title
  {4.5cm}   % space after chapter title

% Language support
\usepackage{polyglossia}
\setdefaultlanguage{greek}
\setotherlanguage{english}

% Font support 
\setmainfont{Georgia}
\newfontfamily\greekfont{Georgia}
\newfontfamily\englishfont{Georgia}

% Typography improvements
\usepackage{microtype}
\usepackage{setspace}
\usepackage{csquotes}
\usepackage{emptypage}

\begin{document}

\frontmatter

\title{
\Huge Πλευρές του Ίδιου Νομίζματος\\[0.5em]
}
\author{Γιώργος Τόλης}
\date{}
\mainmatter


\maketitle

\cleardoublepage

\tableofcontents

\cleardoublepage

\chapter{}

Επτά, έξι, έντεκα και πέντε χιλιόμετρα σήμερα. Τέσσερα, έντεκα και δεκαεπτά την προηγούμενη μέρα. Μπότες να κινούνται πάνω κάτω ξανά και ξανά. Δεν υπάρχει απαλλαγή στον πόλεμο. Δεν κοιτάς τι έχεις μπροστά σου γιατί είναι πάλι μπότες, μπότες που κινούνται πάνω κάτω πάλι. Άντες τρελαίνονται τρελαίνονται από την επανάληψη της κίνησης όλη την μέρα. Τους έχει κατακτήσει ο φόβος, ο φόβος ότι κάποια μέρα δεν θα έχουν στο πλάι τους συναδέλφους τους. Όσο και να προσπαθήσεις να σκεφτείς κάτι διαφορετικό δεν μπορείς. Δεν υπάρχει απαλλαγή στον πόλεμο. Κάνεις ό,τι σου ζητήσουν οι από πάνω σου και ελπίζεις πως είναι το σωστό, όσο διαβολικό και να είναι.

Εμείς χωριστήκαμε από τους υπόλοιπους πριν τέσσερις μέρες. Από τότε ο ρυθμός με τον οποίο περπατάμε έχει αυξηθεί αρκετά καθώς έχουμε μείνει μόνο εγώ, ο Ράινερ και ο Γιόχαν. Ο αντισυνταγματάρχης και ο συνταγματάρχης δηλαδή, απλά μου επέτρεψαν όσο δουλεύουμε στο “Σχέδιο” να τους απευθύνομαι με τα κανονικά τους ονόματα. Γνώριζα πως η ευγένεια τους αυτή δεν θα κρατούσε πολύ μετά την ολοκλήρωση της αποστολής μας, οπότε έπρεπε να την ευχαριστηθώ όσο μπορούσα.

«Ράινερ, μπορείς να μου πιάσεις μία το νερό από το σακίδιο σου;» ζήτησε ο Γιόχαν και πήρε άλλη μία δαγκονιά από το ψωμί που είχε επιλέξει να φάει σήμερα για βραδινό.

Ο Ράινερ του το έδωσε χωρίς να πει κάτι παραπάνω και έγειρε ξανά το πλευρό του στο δέντρο που στερεονόταν όρθιος πριν την παράκληση του Γιόχαν.

«Μα καλά εσείς δεν θέλετε να φάτε τίποτα;» μας ρώτησε ο Γιόχαν. «Ειδικά εσύ Μάρκο. Κανόνισε να φτάσουμε στο αρχοντικό και να είσαι πιο λεπτός από το πιστόλι μου»

«Δεν προλαβαίνω να χάσω τόσα κιλά Γιόχαν, χαλάρωσε» του απάντησα. «Εξάλλου, δεν νομίζω να τους πειράξει κιόλας. Ίσα ίσα, μπορεί να φανώ πιο πειστικός»

«Καλά, ό,τι νομίζεις. Πρέπει να τρώτε τρία γεύματα την ημέρα. Μας δώσανε αρκετό φαγητό για το ταξίδι μας, δεν νομίζω να μας τελειώσει πριν φτάσουμε. Εξάλλου πόσο ακόμη θέλουμε. Μία; Δύο μέρες;»

«Δύο. Με τον ρυθμό που περπατάμε τουλάχιστον.» του είπε ο Ράινερ.

«Γιατί, περπατάμε αργά;»

«Όχι, ίσα ίσα. Απλά αν κάνουμε τριάντα χιλιόμετρα τη μέρα θα φτάσουμε το πρωί της έβδομης μέρας, όπως ακριβώς είχε υπολογίσει ο στρατηγός»

«Και να σου πω ρε Ράινερ, εσύ είσαι βέβαιος ότι θα καταφέρεις μόνος σου;»

Ο Ράινερ του έριξε μία ειρωνική στραβή ματιά λες και δεν ήταν στις ειδικές δυνάμεις του στρατού της χώρας πριν τον χωρισμό που προκάλεσε ο πόλεμος.

«Πιστεύεται ότι θα τα καταφέρουμε;» ρώτησα για να εμπλακώ και γω στη συζήτηση.

«Εμείς ή η επανάσταση συνολικά;» είπε ο Ράινερ

«Ε… κυρίως η μεγάλη εικόνα με προβληματίζει. Τι και αν οι υπολογισμοί του στρατηγού είναι λάθος. Τι και αν δεν είμαστε τυχεροί; Τι και αν…»

«Αν η γιαγιά μου είχε ροδάκια θα ήταν κάρο» με διέκοψε ο Γιόχαν και ξαναπήρε άλλη μία δαγκονιά από  το ψωμί του. «Αν θέλει ο Θεός να νικήσουμε, θα μας βοηθήσει»

«Άσε τα θεία τώρα γέρο.» του είπε εκνευρισμένος ο Ράινερ. «Το μόνο στο οποίο εξαρτόμαστε αυτή τη στιγμή είναι οι υποκριτικές σας ικανότητες, στην διακριτικότητα μου και στο στρατό μας στις πρώτες γραμμές του πολέμου. Όλα τα άλλα είναι τελείως άσχετα με τη μοίρα μας»

Άσχετα μπορεί να μην ήταν, αλλά ο Ράινερ είχε δίκιο. Ο μόνος που θα έπρεπε να βασιζόμαστε από όταν αρχίσει η πρώτη φάση του “Σχεδίου” θα ήταν οι εαυτοί μας και οι ικανότητες μας, χωρίς εξωτερικούς παράγοντας να μας αποσπούν. 

«Γρήγορα στα φυτά»

«Τι;» είπε σαν να μην τον άκουσε ο Γιόχαν.

«Μη ρωτάς, κάναμε βλακεία» είπε ο Ράινερ και πήδηξε για να κρυφτεί πίσω από μερικούς θάμνους, πετώντας δύο από τις τσάντες με τις προμήθειες πιο πέρα, μακριά από τη φωτιά. 

Εμείς ακολουθήσαμε συντώμος, κρύβοντας και την τρίτη τσάντα μας στον θάμνο που σκύψαμε όλοι μας για να καληφθούμε. Η φωτιά έμεινε πίσω αναμένη.

Ο Ράινερ άκουσε ήχους που έμοιαζαν με φωνές σε συνδυασμό με ρόδες να τρίβονται στο έδαφος. Το μέρος που είχαμε κατασκηνώσει εκείνο το απόγευμα έτυχε και ήταν δέκα μέτρα από έναν δρόμο και κανείς από τους τρεις μας δεν το αντιλήφθηκε. Πάλι καλά ο Ράινερ μας είπε να κρυφτούμε, γιατί αλλιώς το “Σχέδιο” θα είχε αποτύχει πριν καν ξεκινήσει η πρώτη φάση του. Η φωτιά όμως; Την ξεχάσαμε αναμμένη. Από την άλλη, δεν μπορούσαμε να κάνουμε και τίποτα. Οι περαστικοί θα παρατηρούσαν είτε την φωτιά, είτε εμάς. Πιστεύω κάναμε τη σωστή επιλογή.

Δεν μπόρεσα να δω το αυτοκίνητο που πέρασε, ούτε μπόρεσα να διακρίνω κάποια φωνή καθώς αυτό μας προσπερνούσε. Είχαμε πέσει και οι τρεις μπρούμυτα ώστε να μας καλύπτουν τα κλαδιά των θάμνων.

«Τώρα σε τι βασιστήκαμε;» είπε ειρωνικά ο Γιόχαν.

«Στην ικανότητα μας να προσαρμοζόμαστε» του απάντησε ο Ράινερ και σηκώθηκε, ξεσκονίζοντας τη ζακέτα του και αναθεωρώντας  την χαζή επιλογή του να ανάψουν φωτιά πριν νυχτώσει επειδή οι τρεις τους είχαν κουραστεί. «Σηκωθείτε τώρα, δεν είμαστε σε καλό σημείο, πρέπει να μετακινηθούμε»


\chapter{}

«Έξι, επτά, έντεκα ή και όσες θέλετε. Σας ξαναλέω, όσες φορές και να έρθετε στο γραφείο μου στις πρώτες γραμμές, που σας ξαναλέω είναι επικίνδυνο, αλλά εσείς δεν το καταλαβαίνετε κυρία μου, δεν υπάρχει περίπτωση ο γιος σας να γίνει στρατηγός» είπε ο στρατηγός μας στη μητέρα μου εκνευρισμένος με την επανεμφάνιση μας στο στρατόπεδο για πολλοστή φορά. 

Ακόμη και εγώ έχω χάσει το μέτρημα πλέον. Κάθε λίγες μέρες τον τελευταίο μήνα η μητέρα μου με βάζει να πηγαίνω στον στρατηγό και να του δείχνω τις στρατηγικές πολέμου που ζωγραφίζω στο τετράδιο μου πάνω από τις καρικατούρες για χάρτη που σχεδιάζω για να είναι ρεαλιστικά τα σενάρια. Άρχισα να σκέφτομαι διάφορους τρόπους οργάνωσης του στρατού ώστε να μεγιστοποιείται την επίδοση της επίθεσης και η άμυνα μας να είναι όσο το δυνατόν πιο κλειστή όχι γιατί είχα σκοπό να γίνω αρχηγός ή κάτι, αλλά για να περάσει ο χρόνος. Από τότε που άρχισε ο πόλεμος δεν έχω και τίποτα να κάνω στο σπίτι, μόνο αυτό και αργότερα η αναπαράσταση των στρατηγικών με τα ξύλινα στρατιωτάκια είχα για να περάσω τον χρόνο μου. Κυρίως μου τράβηξε το ενδιαφέρον να προσπαθώ να οργανώσω τρόπους με τους οποίους θα μπορούσαμε να σκορπίσουμε τον στρατό μας επειδή η τρέχουσα στρατηγική με την οποία λειτουργεί ο τρέχων στρατηγός έχει ένα τεράστιο μειονέκτημα που αν δεν ήταν προσεκτικός θα του κόστιζε τον πόλεμο.

«Μα, οι στρατηγικές… δεν… δεν σας αποδεικνύουν ότι ο γιος μου έχει την ικανότητα να…»

«Ο γιος σας δεν είναι ικανός να ηγηθεί στρατό» διέκοψε τη μητέρα μου ο στρατηγός. «Είναι ακόμη νέος. Η μόνη ελπίδα του να γίνει υψηλόβαθμος στον στρατό θα ήταν η πτώση όλης της πρώτης γραμμής. Και τότε, όχι μόνο ο γιος σας, αλλά μάλλον και όλοι οι άλλοι νέοι από τα υπόλοιπα αρχοντικά θα αγωνίζονταν για την κατάκτηση βαθμών στον στρατό»

«Δεν μπορείς να είσαι σοβαρός. Ο Μάρκο Μπερνούλι είναι γιος του προηγούμενου στρατηγού. Πως μπορείς να τον υποτιμάς έτσι;»

«Μπορώ και θα το κάνω κυρία μου. Είναι δουλειά μου να διακρίνω τις ικανότητες των αντρών μου και στο γιο σας δεν βλέπω τις απαραίτητες ικανότητες ηγεσίας.»

«Καλά, μη μας ξαναμιλήσεις ποτέ εσύ και οι άντρες σου, ακούς;» του φώναξε εκνευρισμένη η μητέρα μου. «Εμείς φεύγουμε τώρα και δεν θα ξαναγυρίσουμε»

«Χαρά μου» της είπε ειρωνικά ο στρατηγός αποχεραιτώντας την και στρέφοντας κατευθείαν την προσοχή του στον βοηθό του που κρατούσε μερικά έγγραφα και περίμενε αρκετή ώρα για να τελειώσει τη συζήτηση με τη μητέρα μου.

«Αντίο» είπε η μητέρα μου και με τράβηξε από το χέρι. «Πάμε Μάρκο, Άλμπερτ»

«Όπως επιθυμείτε κυρία μου» είπε ο Άλμπερτ και μας συνόδεψε μέχρι το αυτοκίνητο μας όπου μας άνοιξε και την πίσω πόρτα ώστε να κάτσουμε. Αυτός κάθισε μπροστά από το τιμόνι και άρχισε να οδηγεί. 

«Ο γιος σας δεν έχει ικανότητες ηγεσίας…» μιμήθηκε η μητέρα μου τον στρατηγό και έβαλε τα κλάματα. «Δεν καταλαβαίνω γιατί δεν υπάρχει σεβασμός μεταξύ αριστοκρατών. Το σύστημα πάντα ήταν “ο γιος του κάτι είναι ο επόμενος κάτι”, αλλά τώρα που δόθηκε σε  αυτόν τον στρατηγό η αρχηγία προσωρινά… δηλαδή… δέκα χρόνια, αλλά δεν έχει σημασία αυτό. Από τότε που πέθανε ο πατέρας σου Μάρκο, επειδή ήσουν μικρός, έχει αρχίσει να μη μας ευνοεί»

«Δεν πειράζει κυρία» προσπάθησα να την ανακουφίσει ο Άλμπερτ, όμως χωρίς αποτέλεσμα.

Το παραλήρημα της μητέρας μου συνέχισε καθόλη τη διάρκεια του πεντάωρου ταξιδιού μας. Ακόμη και στο διάλειμμα που έκανε ο Άλμπερτ για να γεμίσει το αυτοκίνητο πάλι καύσιμα, η μητέρα μου συνέχισε να συζητάει με τον εαυτό της για το “σύστημα” και το πως ήμουν μία… όχι απογοήτευση απαραίτητα, αλλά συνώνυμες λέξεις ακούστηκαν να βγαίνουν από το στόμα της. 

Η μόνη στιγμή που κάποιος άλλος εκτός από την μητέρα μου μίλησε κατά τη διάρκεια του ταξιδιού πίσω στο αρχοντικό ήταν όταν παρατήρησα μία αναμμένη φωτιά στο δάσος, περίπου δέκα μέτρα μακριά από τον δρόμο. Προσπάθησα να σηκωθώ λίγο από το κάθισμα μου για να παρατηρήσω πιο προσεκτικά πόσοι άνθρωποι βρίσκονταν κοντά στη φωτιά, όμως μπόρεσα να δω μόνο δύο βαθουλώματα στο έδαφος. Λες να είχαν φύγει οι περιπλανόμενοι; Τι δουλειά είχαν μέσα στο δάσος; Ήταν μαζί μας; Ήταν αντιστασιακοί; Γιατί υπάρχει καν κάποια φωτιά αναμένη ενώ δεν έχει νυχτώσει καν; Όλα τα σενάρια ήταν ανοιχτά, δεν μπορείς να βγάλεις συμπέρασμα μόνο με μία φωτιά.

Ακόμη και να μην είναι απειλή για εμάς, ποιος ο λόγος να περιπλανούνται δύο άνθρωποι σε περίοδο πολέμου; Είναι χαμένοι; Αν δεν είναι, γνωρίζουν προς τα πιες μεριές είναι τα στρατόπεδα και οι εμπόλεμες ζώνες; 

«Άλμπερτ, πόσο ακόμη έχουμε για να φτάσουμε σπίτι;» τον ρώτησα από περιέργεια να δω πόσο κοντά βρισκόμασταν στο αρχοντικό.

«40 λεπτά» μου απάντησε και έριξε μια πεταχτή ματιά προς το μέρος μου χωρίς να γυρίσει το κεφάλι του. «Γιατί ρωτάτε κύριε;»

«Τίποτα… κάτι σκέφτομαι»


\chapter{}

«Μα δεν θέλω να φύγω μαμά» παραπονέθηκε η αδελφή μου. «Θέλω να μείνω με τον μπαμπά»

«Δεν μπορείς, είναι επικίνδυνο να μείνεις στην πρωτεύουσα Γκρέτα μου» της απάντησε η μητέρα.

Το σπίτι μας σε λίγες μέρες θα γινόταν στρατόπεδο και αργότερα ίσως πεδίο μάχης αν υπήρξε υποχώρηση από τους αντιστασιακούς. Ο Πατέρας μας ως στρατηγός θα έμενε πίσω προφανώς για να ηγηθεί του πολέμου. Η αλλαγή αυτή φέρνει, όχι μόνο εμάς, αλλά όλες τις οικογένειες σε δύσκολη θέση καθώς πρέπει να μετακινηθούμε μακριά από τον τόπο στον οποίο ζούμε τα τελευταία χρόνια, μερικοί ίσως όλοι τους τη ζωή, και να αφήσουμε παράλληλα πίσω μέλη της οικογένειας μας με την ελπίδα ότι κάποια μέρα θα τους ξαναδούμε.

Η πιθανότητα όμως οριστικού αποχαιρετισμού ήταν αυτό που ενοχλούσε όλους εκείνη τη μέρα. Η πράξη του να αφήσεις τον γιο σου, τον πατέρα σου, τον αδελφό σου να πεθάνει χωρίς τη θέληση του για έναν πόλεμο που ποιος ξέρει καν γιατί εξελίσσεται σου προκαλεί περίεργα συναισθήματα. Άλλες οικογένειες περνάνε αυτή τη στιγμή πιο ήπια, μία απλή αγκαλιά αρκεί για να πουν αντίο στους συγγενείς τους. Άλλοι, σαν εμάς, χρειάζονται τον χρόνο τους.

«Μπορείτε να φύγετε μαμά, εγώ θα κάτσω εδώ με τον μπαμπά» επέμενε η Γκρέτα.

«Αγάπη μου…» την πλησίασε ο πατέρας μας και άπλωσε το χέρι του γύρω από τον ώμο της «…δεν μπορείς να μείνεις εδώ. Τις επόμενες μέρες θα εξελιχθεί ραγδαία ο πόλεμος και πολύ πιθανόν…»

«Πολύ πιθανόν τι;» τον έσπρωξε η αδελφή μου και τον κοίταξε με τα δύο παιδικά της μάτια να αφήνουν ένα κύμα δακρύων περιμένοντας την απάντηση του.

«Αν μείνεις εδώ δεν θα μπορούμε να σε φροντήσουμε» της απάντησε προσπαθώντας να μην ξεστομίσει κάτι που θα μετάνιωνε.

«Δε με πειράζει αυτό. Εγώ θέλω να είμαι δίπλα στον πατερούλη μου. Δίπλα…»

«Αρκετά Γκρέτα» της φώναξε η μητέρα και την τράβηξε προς το μέρος μας. «Ανέβα στο τρένο και περίμενε εμένα και τον Μάρκο στις θέσεις μας»

Η Γκρέτα την κοίταξε με μίσος. Ύστερα από λίγα δευτερόλεπτα σιωπής  μετά την ένταση μεταξύ τους, η Γκρέτα ανέβηκε τις σκάλες και έστριψε δεξιά διστακτικά για να βρει τις θέσεις μας. Πριν την χάσουμε από το οπτικό μας πεδίο, αυτή σταμάτησε για λίγο.

«Θα πάω… τουαλέτα» ήταν το τελευταίο πράγμα που μας είπε και εξαφανίστηκε.

Η μητέρα μου περίμενε λίγο να απομακρυνθεί η Γκρέτα από την πόρτα του τρένου. Όταν έκρινε πως ήταν σε αρκετά μεγάλη απόσταση, έπεσε στην αγκαλιά του άντρα της και με την σειρά της άφησε τα δικά της δάκρυα να κυλήσουν πάνω στη στολή του στρατηγού. Εκείνος το μόνο που μπορούσε να κάνει ήταν να την αγκαλιάσει μία τελευταία φορά και να την φιλήσει στο πάνω μέρος του κεφαλιού καθώς αυτή είχε σωριαστεί στο στήθος του.

«Μάρκο, έλα εδώ και συ» μου είπε ο πατέρας μου κάνοντας νεύμα να πλησιάσω.

Και αυτό έκανα. Άφησα τις βαλίτσες που κουβαλούσα απαλά στο πάτωμα και στάθηκα δίπλα στον πατέρα μου. Αυτός μου χαμογέλασε και με τράβηξε σφιχτά πάνω του. Δεν πρόλαβα στην ηλικία που ήμουν να αντιδράσω κατευθείαν, ούτε ως αντανακλαστικό. Ένιωσα μόνο τον πατέρα μου να σφίγγει και τους δύο μας πάνω του. Τότε ήταν που κατάλαβα ότι ο πατέρας μας δεν ήταν κάποιος ακραία θαρραλέος χωρίς κανένα άλλο συναίσθημα που θα πήγαινε στις πρώτες γραμμές του πολέμου και θα οδηγούσε την αντίσταση στη νίκη με την ηγεσία του. Όχι. Η είκονα που είχα για τον πατέρα μου κατέρρευσε εκείνη τη στιγμή. Ξαφνικά ο κλειστός άντρας που ήξερα τόσα χρόνια εξέφραζε συναισθήματα. Συναισθήματα φόβου. Φόβου αυτού που σκεφτόμουν προηγούμενος; Ότι δεν θα γυρίσει; Ότι μπορεί να μην μας ξαναδεί πότε; Ποιος ξέρει. 

«Όλα θα πάνε καλά» ψιθύρισε στο αυτί της μητέρας μου καθώς μας έκλεινε πάνω του πιο σφιχτά για να απολαύσει την ζεστασιά της οικογένειας του μία τελευταία φορά πριν σηκώσει το κεφάλι και μας αφήσει να ανεβούμε στο τρένο.

«Συγγνώμη για την Γκρέτα» του απολογήθηκε η μητέρα.

«Δεν χρειάζεται να συγχωρείς. Είναι παιδί ακόμη. Πόσο μάλλον κοριτσάκι. Τα συναισθήματα που μπορεί νιώθει αυτή τη στιγμή δεν μπορούν να περιγραφούν από κανέναν. Το μόνο που μπορούμε να κάνουμε αυτή τη στιγμή είναι υπομονή» της απάντησε ο πατέρας μου με το χέρι στον ώμο της και ύστερα έκανε μισό βήμα πίσω για να μας αφήσει να προχωρήσουμε προς τις αποσκευές.

«Όλα θα πάνε καλά…» ψιθύρισε η μητέρα μου καθώς σήκωνε την πρώτη βαλίτσα για να μου τη δώσει. «Αντίο λοιπόν»

«Καλό ταξίδι» ήταν οι τελευταίες λέξεις που άκουσα τον πατέρα μου να ξεστωμίζει.

«Πάμε Μάρκο, ανέβα» μου είπε η μητέρα μου και με άφησε να ανέβω πρώτος στο τρένο.

Στρέφοντας το βλέμμα μου δεξιά για να προχωρήσω στον διάδρομο προς τις θέσεις μας είδα άλλες οικογένειες όπως έξω από το τρένο. Άλλες οικογένειες που όπως και στον έξω κόσμο περνούσαν όμοια κατάσταση με εμάς. Όλοι τους βέβαια είχαν την προσοχή τους στραμμένη πάνω μας, στην οικογένεια του στρατηγού, οι οποίοι θα ταξίδευαν μαζί τους. Όλοι τους παράλληλα ήξεραν, μαζί και εμείς, ότι δεν μας ξεχώριζε τίποτα από το πλήθος. Τύχαινε απλά ο πατέρας μου να είναι “ο στρατηγός”.  

Κανείς δεν σχολίασε την κατάσταση, όχι από κακία, από σεβασμό. Από σεβασμό που όλες οι οικογένειες αντιμετωπίζονται ισότιμα. Που όλοι έχουν την ευκαιρία να φύγουν από την πρωτεύουσα με τα πρώτα τρένα, χωρίς διακρίσεις. Πως κάποιοι δεν ευνοούνται περισσότερο από την θέση κάποιου άλλου στην οικογένεια τους για να σώσουν το τομάρι τους. Όλοι μας εκείνη τη στιγμή βρισκόμασταν στο ίδιο επίπεδο.

Όταν βρήκαμε τελικά τις θέσεις μας, άφησα την βαλίτσα μου στο πάτωμα και η μητέρα μου ακούμπησε προβληματισμένη τις δύο δικές της στο κάθισμα της στο ένα από τα τέσσερα καθίσματα που μας είχαν ανατεθεί και συνέχισε να περπατάει στον διάδρομο. Μόνος μου τώρα, δεν είχα πλέον κάποιον να μου μιλήσει, να γκρινιάζει ή να σωριάζεται επειδή φεύγουμε. Ο μόνος ήχος που έφτανε στο αυτί μου πλέον ήταν η βαβούρα από όλους τους υπόλοιπους στο τρένο που συζητούσαν. 

Χαλάρωσα, γιατί τι άλλο μπορούσα να κάνω, και κοίταξα έξω από το παράθυρο. Εκεί είδα τον πατέρα μου ακίνητο μέσα σε ένα κινούμενο πλήθος να μου χαμογελάει. Εγώ παρέμεινα σκυθρωπός, χωρίς να εκφράσω κάποιο συναίσθημα. Σήκωσα μόνο το χέρι μου για να στηρίξω το κεφάλι μου το οποίο ακούμπησα εν τέλη στη τζαμαρία και κοίταξα προς τις ράγες ανάμεσα από την αποβάθρα και το τρένο. 

Λίγα δευτερόλεπτα μετά, ένα να χέρι μου χτύπησε το τζάμι. Ήταν πάλι ο πατέρας μου. Τον είδα πάλι να χαμογελάει και ξεφύσηξα. Αυτός, βλέποντας με σε απόγνωση μου ξεστόμησε  μία λέξη με έντονη άρθρωση: χαμογέλα. Εκείνη τη στιγμή όμως δεν μπορούσα να χαμογελάσω ούτε να προσπαθούσα. Ο πατέρας μου το πήρε ως πρόκληση μάλλον γιατί άρχισε να κάνει αστείες γκριμάτσες τη μία μετά την άλλη, και για να πω την αλήθεια είχα ξεκαρδιστεί όντας παιδί. Στη συνέχεια άρχισε να προσποιείται πως είναι πεζός που διαβάζει εφημερίδα και σκοντάφτει σε μπανανόφλουδα. Αυτό το αστείο πάντα έκανε τον μικρό Μάρκο να γελάσει, όσες φορές και να το είχε δει να εξελίσσεται, γιατί κάθε φορά ο πατέρας μου έβαζε και μια νέα πινελιά.

Το καμπανάκι του τρένου όμως τώρα χτυπούσε. Οι πόρτες έκλεισαν. Το πλήθος από στρατιώτες μειωνόταν. Ο πατέρας μου; Πλησίασε το τρένο και άρχισε να σπρώχνει το τρένο από το πλάι. Εγώ παραξενεύτηκα στην αρχή γιατί αυτό που έκανε ήταν ανούσιο, όμως αυτός το έκανε για μένα, για μία τελευταία παντομίμα. Σταμάτησε να σπρώχνει μετά από λίγη ώρα γιατί έβλεπε ότι το τρένο δεν ξεκινούσε και έκανε πως βαριανασαίνει. Όταν ήχησε ξανά το καμπανάνι του τρένου ο πατέρας μου σήκωσε τα μανίκια της στολής για να φανούν τα μπράτσα του και προσποιήθηκε πως τα ξεσκονίζει πριν λυγίσει τα χέρια του λες και κάνει επίδειξη των μυών του. Έπειτα ήρθε πάλι δίπλα στο τρένο και μιμήθηκε πως το σπρώχνει από το πλάι πάλι, όμως τώρα ο οδηγός του τρένου θα βοήθαγε τον πατέρα μου και το μυαλό του οκτάχρονου Μάρκου είχε τρελαθεί. Ο πατέρας μου συνέχισε να σπρώχνει το τρένο όσο αυτό επιτάχυνε, μέχρις ότου δεν μπορούσε να συμβαδίσει με αυτό. Πριν αφήσει τα χέρια του από το τρένο έκανε πως το εκσφενδονίζει λες και είναι ακόντιο και έπεσε κάτω. Ένα τελευταίο αστείο για τον γιο του ή μία αστούμπαλι πτώση από τον στρατηγό, ποιος ξέρει. Εγώ εκείνη τη στιγμή γέλασα και είδα τον πατέρα μου να σηκώνει το κεφάλι του προς τα πάνω για να δει αν πέτυχε τον σκοπό του. Η τελευταία εικόνα που έχω από τον πατέρα μου είναι αυτό το χαμόγελο που είχε ενώ σηκωνόταν. Μετά βγήκαμε από τον σταθμό και δεν μπορούσα να δω πλέον τους ανθρώπους στην αποβάθρα. \newline\newline

«Έτσι ήταν λοιπόν που λέτε η τελευταία συνάντηση με τον πατέρα μου. Όμως γιατί ρωτάς Γιόχαν;» παραξενεύτηκα με την περίεργη ερώτηση που μου είχε κάνει από το πουθενά ο αντισυνταγματάρχης όταν έμαθε ότι είμαι γιος του πρώην στρατηγού.

«Για τον πατέρα σου δυστυχώς ξέρεις ότι δεν είναι πλέον μαζί μας» μου απευθύνθηκε και κάθισε αντιδιαμετρικά μου, ρίχνοντας ένα μικρό ξυλαράκι στη φωτιά. «Την αδελφή σου την βρήκατε ποτέ;»

«Όχι…» του απάντησα και παραξενεύτηκα. «Όλο το ταξίδι την ψάχναμε στο τρένο για να βρούμε που έχει κρυφτεί αλλά δεν καταφέραμε να την βρούμε. Η μητέρα μου υπέθεσε πως η Γκρέτα κατέβηκε από το τρένο και πήγε στον πατέρα μου και αυτός θα την βοηθούσε να μπει σε κάποιο άλλο τρένο με κάποιον συνεργάτη ίσως για να αποχωρήσει. Όμως ούτε αυτό έγινε ποτέ και μετά από έναν μήνα παρακολούθησης των επιβατών που έβγαιναν από τα τρένα στην αποβάθρα κοντά στο νέο μας σπίτι τα παρατήσαμε και θεωρήσαμε πως η αδελφή μου πήγε σε κάποιο άλλο μέρος για να τη φροντήσουν. Μπορεί ο πατέρας μου να μην μπορούσε να την στείλει σε μας, μπορεί να μην ήθελε αυτή να επιστρέψει σε μας, μπορεί…»

«Έχεις δίκιο πως η αδελφή σου είχε κατέβει από το τρένο» διέκοψε τον συλλογισμό μου ο Γιόχαν «…αλλά δεν επιβιβάστηκε ποτέ σε άλλο τρένο»


\chapter{}


«Μα δεν θέλω να φύγεις πατέρα» παραπονέθηκε η αδελφή μου. «Γιατί δεν μπορείς να μείνεις μαζί μας;»

«Είναι πολύ ριψοκίνδυνο να μην είμαι στις πρώτες γραμμές του πολέμου Κλάρα μου» της απάντησε η μητέρα.

Ο πατέρας μας είχε μαζέψει στην είσοδο του αρχοντικού για να μας ανακοινώσει ότι έπρεπε να αποχωρήσει με προορισμό ένα εγκαταλελειμμένο χωριό κοντά στην πρωτεύουσα, το οποίο είχε γίνει στρατόπεδο και από όπου θα μπορούσε να ηγηθεί του στρατού των αριστοκρατών χωρίς να χρειάζεται να στέλνει αγγελιοφόρους. Ως στρατηγός, έπρεπε να πάρει αυτή τη δύσκολη απόφαση. Δύσκολη πάντα ήθελε να βρίσκεται εντός μίας ώρας ταξίδι με το αυτοκίνητο από το αρχοντικό της οικογένειας του. Όμως θα έκανε τη διαχείριση του στρατού του πολύ πιο άμεση και αυτό θα ήταν μεγάλη αναβάθμιση σε σχέση με την τρέχουσα κατάσταση που επικρατεί.

Η αλλαγή αυτή βέβαια θα είχε το προφανές μειονέκτημα ότι ο πατέρας θα βρισκόταν μακριά από εμάς, που δεν θα επηρέαζε κανέναν άλλον παρά εμάς. Το γεγονός ότι υπάρχει πιθανότητα να μην τον βλέπαμε ποτέ ξανά είχε περάσει από τη σκέψη ολονών μας, ιδιαίτερα της αδελφής μου. Ήταν η πρώτη φορά που ο πατέρας θα πήγαινε τόσο βαθιά στην καρδιά του πολέμου, πέντε ώρες ταξίδι από το σπίτι του, όμως ήθελε να είναι η πρώτη και η τελευταία φορά που θα το έκανε. Ήθελε να βάλει τέλος στον πόλεμο που ποιος ξέρει καν γιατί οι δύο πλευρές χύνουν το αίμα των συνανθρώπων τους. Ίσως φταίει η απληστία. Ίσως υπάρχει κάποια αδικία που δεν βλέπει η αντίθετη μεριά. Μόνο οι αρχηγοί ξέρουν.

«Μπορείς να μου φτιάξεις και μένα μια βαλίτσα Άλμπερτ, θα έρθω με τον μπαμπά» επέμενε η Κλάρα.

«Αγάπη μου…» την πλησίασε ο πατέρας μας και άπλωσε το χέρι του γύρω από τον ώμο της «…δεν μπορείς να έρθεις μαζί. Τις επόμενες μέρες θα εξελιχθεί ραγδαία ο πόλεμος και πολύ πιθανόν…»

«Πολύ πιθανόν τι;» τον έσπρωξε η αδελφή μου και τον κοίταξε με τα δύο παιδικά της μάτια να αφήνουν ένα κύμα δακρύων περιμένοντας την απάντηση του.

«Αν έρθεις μαζί μου δεν θα μπορούμε να σε φροντήσουμε» της απάντησε προσπαθώντας να μην ξεστομίσει κάτι που θα μετάνιωνε.

«Δε με πειράζει αυτό. Εγώ θέλω να είμαι δίπλα στον πατερούλη μου. Δίπλα…»

«Αρκετά Κλάρα» της φώναξε η μητέρα μας την τράβηξε προς το μέρος μας. «Ανέβα στο δωμάτιο σου και περίμενε εμένα και τον Μάρκο να έρθουμε σε λίγο»

Η Κλάρα την κοίταξε με μίσος. Ύστερα από λίγα δευτερόλεπτα σιωπής μετά την ένταση μεταξύ τους, η Κλάρα ανέβηκε τις σκάλες και έστριψε δεξιά για να πάει προς το δωμάτιο της, ρίχνοντας μία κλεφτή ματιά πίσω για να δει αν την κοιτούσαμε. Πριν τη χάσουμε από το οπτικό μας πεδίο, αυτή σταμάτησε για λίγο και στράφηκε προς το μέρος μας λες και ήθελε να πει κάτι τελευταίο, αλλά σταμάτησε πριν βγει κάποια λέξη από το στόμα της και έσπευσε προς το δωμάτιο της.

Η μητέρα περίμενε λίγο να απομακρυνθεί η Κλάρα από τις σκάλες. Όταν έκρινε πως ήταν σε αρκετά μεγάλη απόσταση, έπεσε στην αγκαλιά του άντρα της και με την σειρά της άφησε τα δικά της δάκρυα να κυλήσουν πάνω στη στολή του στρατηγού. Εκείνος το μόνο που μπορούσε να κάνει για να την παρηγορήσει ήταν να την αγκαλιάσει μία τελευταία φορά και να την φιλήσει στο πάνω μέρος του κεφαλιού της καθώς αυτή είχε σωριαστεί στο στήθος του.

«Μάρκο, έλα δω και συ» μου είπε ο πατέρας κάνοντας μου νεύμα να πλησιάσω.

Και αυτό έκανα. Προχώρησα προς το μέρος του σιγά σιγά και αυτός μου χαμογέλασε και με τράβηξε σφιχτά πάνω του. Δεν πρόλαβα στην ηλικία που ήμουν να αντιδράσω κατευθείαν, ούτε ως αντανακλαστικό. Ένιωσα μόνο τον πατέρα μου να σφίγγει και τους δύο μας πάνω του. Τότε ήταν που κατάλαβα πως ο πατέρας μας δεν ήταν… όχι μη θαρραλέος, αυτό θα ήταν προσβλητικό. Απλά μου χάλασε την παραίσθηση που είχα για αυτόν πως είναι ένας βράχος χωρίς συναισθήματα που ηγείται τον στρατό του με πειθαρχία. Όμως ο πατέρας μου είχε δείξει την ευάλωτη του πλευρά για πρώτη φορά. Ξαφνικά ο κλειστός άντρας που γνώριζα τόσα χρόνια εξέφραζε συναισθήματα. Συναισθήματα φόβου. Φόβου ότι εκθέτει τη ζωή του για μία μεγαλύτερη πιθανότητα νίκης, η οποία θα σήμανε τέλος. Τέλος σε ένα μαρτύριο.

«Όλα θα πάνε καλά» ψιθύρισε στο αυτί της μητέρας καθώς μας έκλεινε πάνω του πιο σφιχτά για να απολαύσει την ζεστασιά της οικογένειας του μία τελευταία φορά πριν σηκώσει το κεφάλι του και χαλαρώσει την αγγαλιά του.

«Συγγνώμη για την Κλάρα» του απολογήθηκε ο μητέρα.

«Δεν χρειάζεται να απολογείσαι. Είναι παιδί ακόμη. Πόσο μάλλον κοριτσάκι. Πρέπει να υπερχειλίζει από συναισθήματα αυτή τη στιγμή που εμείς δεν μπορούμε να συλλάβουμε. Το μόνο που μπορούμε να κάνουμε είναι υπομονή» της απάντησε ο πατέρας μας με το χέρι στον ώμο της και ύστερα έκανε μισό βήμα πίσω για να κοιτάξει να δει που ήταν ο Άλμπερτ… εσύ, που ήσουν εσύ. 

«Όλα θα πάνε καλά…» άκουσα τη μητέρα να ψιθυρίζει καθώς έκανε και αυτή μερικά βήματα πίσω για να αποδεχτεί την κατάσταση. «Καλό ταξίδι λοιπόν»

«Αντίο…» μας χαμογέλασε ο πατέρας για να να σιγουρέψει πως η τελευταία εικόνα που είχα από αυτόν θα ήταν κάτι ευχάριστο, όμως εγώ δεν μπορούσα να τον κοιτάξω.

«Το αυτοκίνητο είναι έτοιμο κύριε» εμφανίστηκες τότε εσύ από την είσοδο του αρχοντικού καλώντας τον πατέρα. 

Όμως αυτός δεν είχε σκοπό να αφήσει και τα δύο του παιδιά θλιμμένα πίσω. Τον είδα να με πλησιάζει και αγχώθηκα. Αυτός, με το χαρακτηριστικό του χαμόγελο έσκηψε μπροστά μου και μου είπε μία λέξη: «χαμογέλα». Εκείνη τη στιγμή όμως δεν μπορούσα να χαμογελάσω ούτε να προσπαθούσα. 

Την υπόλοιπη ιστορία την ξέρεις αφού ήσουν εκεί. Ο πατέρας μου το πήρε ως πρόκληση και πήγε ενάμιση αιώνα στο παρελθόν και έγινε γελοτοποιός, κάνοντας μου γκριμάτσες και λέγοντας μου λογοπαίγνια που ήξερε ότι μου άρεσαν, και ο μικρός Μάρκο προφανώς δεν μπορούσε να αντισταθεί. Εσύ απ’ ότι θυμάμαι παρατηρούσες από την είσοδο και έριξες και ένα πεταχτό χαμογελάκι βλέποντας μας να περνάμε μία τελευταία ωραία στιγμή μαζί.

Αυτό που μου άρεσε περισσότερο όμως εκείνη τη μέρα ήταν όταν ανέβηκα στο δωμάτιο μου και κοίταξα έξω προς τον δρόμο της πόλης για να σας δω να αποχωρείται. Ο πατέρας μου κοιτούσε στο παράθυρο μου για να δει πότε θα εμφανιστεί ο γιος του. Μάλλον ήλπιζε να εμφανιστεί και η κόρη του γιατί φάνηκε σαν να περίμενε λίγο για να αρχίσει το σόου, αλλά από ένα απογοητευμένη εκπνοή, αυτός σήκωσε τα μανίκια της στολής για να φανούν τα μπράτσα του. Έπειτα αυτός προσπάθησε να σπρώξει το αμάξι μας αλλά δεν τα κατάφερε, που προφανώς δεν θα τα κατάφερνε. Στη συνέχεια όμως έκανε μία κίνηση σαν να ξεσκονίζει τους δικέφαλους του και ξαναπροσπάθησε να μετακινήσει το αμάξι από το πλάι. Ο μικρός Μάρκο είχε πάθει πλάκα με το γεγονός ότι το αμάξι κουνήθηκε όντως τη δεύτερη φορά. Μου πήρε αρκετό καιρό να καταλάβω ότι εσύ απλά το προχώρησες λίγο πιο μπροστά και ξενέρωσα πολύ όταν μου το επιβεβαίωσες. Πλάκα είχε όμως και μετά που ο πατέρας έκανε κίνηση σαν να πετάει ακόντιο και άφησε το αυτοκίνητο να κυλήσει μπροστά του και μάλλον από απροσεξία σκόνταψε και πήγες να τον βοηθήσεις. Χάλασε και αυτό λίγο την ψευδαίσθηση που θέλατε να προκαλέσετε, αλλά το μυαλό του μικρού Μάρκου είχε διασκεδάσει το θέαμα. Ο σκοπός του να αφήσει μια ευχάριστη τελευταία εντύπωση στον γιο του είχε επιτευχθεί… ή μπορεί απλά εκείνη τη στιγμή να ήταν ατσούμπαλος. Εγώ θέλω να πιστεύω πως είναι το πρώτο, και μη πεις τίποτα να μου χαλάσεις και αυτή τη ψευδαίσθηση.\newline\newline

«Έτσι ήταν λοιπόν η τελευταία συνάντηση με τον πατέρα από την οπτική μου. Όμως γιατί ρωτάς Άλμπερτ;» παραξενεύτηκα με την περίεργη ερώτηση που μου είχε κάνει από το πουθενά ο μπάτλερ μας όταν φτάσαμε σπίτι από το ταξίδι μας στο στρατόπεδο.

«Για τον πατέρα σου δυστυχώς γνωρίζεις πως δεν είναι πλέον μαζί μας» μου απευθύνθηκε ο Άλμπερτ και μου τράβηξε μία καρέκλα από το τραπέζι για να κάτσω. «Όμως, για την αδελφή σου δεν είχατε ποτέ κάποιο νέο; Έτσι δεν είναι;»

«Όχι…» του απάντησα και παραξενεύτηκα. «Μία εβδομάδα την ψάχναμε σε όλη την πόλη, κανένας όμως δεν είχε παρατηρήσει την Κλάρα ούτε να φεύγει, ούτε να περπατάει στους δρόμους. Θεωρήσαμε μάταιο να συνεχίσουμε να την ψάχνουμε μετά από τόσο καιρό και απλά η μητέρα μου έστειλε ένα γράμμα στον στρατηγό άντρα της να φτιάξει μία ομάδα αναζήτησης της κόρης της, αλλά αυτό δεν έγινε ποτέ. Τι, θες να μου πεις απλά ότι δεν έφτασε το γράμμα στο στρατόπεδο; Είχαμε φτάσει σε αυτό το συμπέρασμα με τη μητέρα εδώ και πολύ καιρό. Ότι μπορεί…»

«Το γράμμα έφτασε στο στρατόπεδο και ο πατέρας σου πρόλαβε να το διαβάσει…» με έβαλε ο Άλμπερτ να καθίσω και ύστερα κάθισε και αυτός δίπλα μου «…αλλά ο πατέρας σου ήξερε πολύ καλά που βρισκόταν η κόρη του»


\chapter{}


«Η αδελφή μου δεν επιβιβάστηκε σε άλλο τρένο;» παραξενεύτηκα με την ξαφνική αλλαγή πλοκής που προκάλεσαν τα λόγια του Γιόχαν. «Τι εννοείς;»

«Αυτό που κατάλαβες. Η αδελφή σου δεν έφυγε ποτέ από τις πρώτες γραμμές του πολέμου» μου απάντησε και με άφησε να επεξεργαστώ τα νέα.

«Ώστε… για αυτό δεν με έχουν αφήσει να πάω ποτέ στις πρώτες γραμμές του πολέμου. Γιατί μπορεί να μου προκληθούν συναισθήματα ή κάτι;»

«Όχι» με γείωσε ο αντισυνταγματάρχης κατευθείαν «αυτό δεν έχει καμία σχέση. Απλά είσαι μικρός ακόμη και όλοι οι στρατιώτες είκοσι δύο ετών και νεότεροι οφείλουν να μένουν στις πίσω γραμμές, εκτός αν ο στρατηγός κρίνει διαφορετικά, όπως τώρα που είμαστε βαθιά σε αντίπαλο έδαφος»

«Τότε… γιατί μου αναφέρεις ότι η αδελφή μου βρίσκεται στις πρώτες γραμμές της αντίστασης;»

«Ο πατέρας σου ήταν καλός άνθρωπος Μάρκο, το ξέρεις;» φάνηκε να προσπαθεί να αλλάξει το θέμα ο Γιόχαν.\newline\newline

Δύο μέρες μετά τις αποχωρήσεις των τρένων από την πρωτεύουσα, όλοι που είχαμε μείνει στήναμε το στρατόπεδο όπου θα μέναμε για ποιος ήξερε τότε πόσο χρόνο. Ως κοντινός συνεργάτης και φίλος του πατέρα σου, είχα την τιμή να μπορώ να στήσω τη σκηνή μου κοντά στου στρατηγού. Όχι δίπλα, όμως αρκετά κοντά όσο προετοιμαζομουν να σηκώσω τους κεντρικούς στύλους για να πάρει το μάτι μου ένα μικρό κοριτσάκι να τρέχει μέσα στο στρατόπεδο, ψάχνοντας κάποιον. Γιατί βρισκόταν στο στρατόπεδο; Άφησα ό,τι έκανα εκείνη τη στιγμή και έτρεξα να την προλάβω πριν προκαλέσει κάποια αναστάτωση. Αυτή όμως ήξερε ποιον ήθελε να βρει και ήθελε να τον βρει γρήγορα. Κοιτώντας με πεταχτές ματιές μέσα στις ήδη σηκωμένες σκηνές, προσέχοντας μη την παρατηρήσει κάποιος από μέσα, το κοριτσάκι κινούνταν λες και ήταν ερπετό στο χώμα και ο μόνος που την είχε παρατηρήσει μέσα στο χάος του στρατοπέδου ήμουν εγώ.

Φτάνοντας στη σκηνή του στρατηγού αυτή έριξε την πεταχτή πλάγια ματιά μέσα στο υφασμάτινο κάλυμμα και χαρούμενη μάλλον που βρήκε αυτόν που έψαχνε, έτρεξε βιαστικά μέσα.

«Μπαμπά» ακούστηκε η φωνούλα της μέσα από τη σκηνή.

«Τι κάνεις εσύ εδώ» της είπε ο στρατηγός χωρίς κανένα μίσος στη φωνή του για το λάθος που είχε κάνει η κόρη του. «Που είναι η μαμά και ο Μάρκο;» την ρώτησε.

«Έφυγαν με το τρένο πριν δύο μέρες» της άκουσα να του λέει πριν γυρίσω να κοιτάξω την κατάσταση μέσα στη σκηνή.

«Και συ γιατί είσαι εδώ; Τι έκανες δύο μέρες;»

«Κρυβόμουν» που απάντησε λες και αυτή ήταν η αυτονόητη απάντηση.

Ο στρατηγός με είδε να μπαίνω στη σκηνή του και μου έριξε μια πεταχτή ματιά λες και είχα εμπλακεί και γω στην εισβολή της κόρης του στο στρατόπεδο αλλά δεν μου μίλησε. Προσπάθησε να λύσει τον προβληματισμό του εσωτερικά.

«Θεέ μου, δεν έχεις φάει τίποτα δύο μέρες» σκέφτηκε ο στρατηγός και σηκώθηκε γρήγορα από την καρέκλα του.

«Δεν πειράζει μπαμπά, μπορώ να περιμένω» του είπε η κόρη του και αγκάλιασε το αριστερό του πόδι.

«Μα…» πήγε να της πει κάτι ο στρατηγός, αλλά δεν ήθελε να της χαλάσει τη στιγμή.

Περίμενε λίγο να χαλαρώσει η λαβή της κόρης του και έσκηψε για να της μιλήσει πρόσωπο με πρόσωπο.

«Αυτό που έκανες δεν είναι σωστό, αλλά δεν είναι και αδιόρθωτο» της εξήγησε προσπαθώντας να της μιλήσει όσο πιο επιεικώς γίνεται. «Γιόχαν

«Μάλιστα!»

«Μπορείς να μας φέρεις δύο τεμάχια από ό,τι έχει μαγειρέψει ο μάγειρας μέχρι στιγμής πριν επιστρέψεις στη δουλειά σου; Αν είναι πολύ νωρίς και δεν έχει μαγειρέψει τίποτα, φέρε μας μία φρατζόλα ψωμί αν μπορείς»

«Μάλιστα» του απάντησα και έσπευσα να εκπληρώσω τη διαταγή του.

«Σε ευχαριστώ φίλε μου» ήταν οι τελευταίες λέξεις που άκουσα από τον πατέρα σου πριν βγω από τη σκηνή του.

Ο μάγειρας προφανώς δεν είχε προλάβει να φτιάξει γεύματα για όλους στο στρατόπεδο εκείνο το πρωί, καθώς ήταν νωρίς ακόμη, μάλλον δέκα ή έντεκα το πρωί. Μου έδωσε δύο πιάτα τα οποία περιείχαν μία δόση από το χθεσινό φαγητό και λίγα λαχανικά από το σημερινό.

«Για ποιον τόσο νωρίς;» με ρώτησε.

«Για τον στρατηγό και την κόρη του» του απάντησα.

«Την κόρη του; Μα δεν έχει αποχωρήσει όλος ο άμαχος πληθυσμός;»

«Όπως φαίνεται όχι» του είπα και πήρα τα πιάτα από τον πάγκο του. «Σε ευχαριστώ»

«Δώσε στον στρατηγό τις καλημέρες μου» μου είπε ο μάγειρας και έστρεψε την προσοχή του πίσω στην κατσαρόλα του με το σημερινό φαγητό.

Δεν θυμάμαι καν πόσο μακριά ήμουν, αλλά καθώς επέστρεφα στη σκηνή του στρατηγού, είδα μία χειροβομβίδα να εμφανίζεται από το πουθενά, να σκίζει το ύφασμα της οροφής της σκηνής και στη συνέχεια μία έκρηξη από το εσωτερικό της σκηνής να την ρίχνει στο έδαφος και να αρχίζει φωτιά και πανικό στο στρατόπεδο. Τα πιάτα που κρατούσα τώρα ήταν η μικρότερη ανησυχία μου πλέον, καθώς τα άφησα στο πάτωμα λες και δεν με ενδιέφερε και έτρεξα προς τη φωτιά. 

Ανάμεσα στα θρύψαλα της σκηνής και της θερμότητας που προκαλούσαν οι φλόγες που έκαιγαν το ύφασμα είδα μία εικόνα που με τραυμάτισε για αρκετό καιρό. Είδα έναν άντρα ο οποίος είχε χάσει κομμάτια δέρματος από την πλάτη του και όλο το υπόλοιπο του σώμα που είχε επιβιώσει ήταν σκούρο από τη μαύρη πυρίτιδα της χειροβομβίδας. Αυτός ο άντρας ήταν ο στρατηγός και βρισκόταν πάνω κυριολεκτικά από την κόρη του, την οποία είχε κάνει μία τελευταία προσπάθεια να προστατεύσει από την αιφνιδιαστική επίθεση που είχε γίνει. 

Είχε πετύχει, η κόρη του βέβαια είχε τραυματιστεί στη δεξιά πλευρά του σώματος της και δεν είχε τη δύναμη να κουνήσει κανένα της άκρο από εκείνη την πλευρά για να σηκώσει τον πατέρα της.

«Μπαμπά, γιατί… δεν κουνιέσαι;» την άκουσα χαρακτηριστικά να λέει, έχοντας τη δυνατότητα να δει μόνο το πρόσωπο του πατέρα της το οποίο είχε χάσει όλη τη ζωντάνια από πάνω του.

«Όχι, όχι, όχι» αγχώθηκα και γύρισα τον στρατηγό ανάσκελα ώστε να πάρω την κόρη του από το πάτωμα και να τρέξω έξω από τη σκηνή. 

Έλεγξα τους παλμούς του και δεν μπόρεσα να νίωσω αίμα να διαπερνά τις φλέβες του λαιμού του. Σηκώθηκα και κοίταξα άλλη μία φορά τα μάτια του, μήπως ο στρατηγός είχε κάποια τελευταία λέξη να πει, όμως μετά από τέτοιο τραύμα που υπέστει είχε χάσει τη ζωή του μέσα σε κλάσματα του δευτερολέπτου. Ένα δάκρυ κύλησε από το μάτι μου βλέποντας τον αρχηγό μας νεκρό τόσο απλά, προσπαθώντας να σώσει την κόρη του. Ήξερα, όμως, πως η επιθυμία του ακόμη και ζωντανός εκείνη τη στιγμή, θα ήταν να σώσει την οικογένεια του.

«Μη φοβάσαι, όλα θα πάνε καλά…» της ψιθύρισα στο αυτί όσο αυτή έβγαζε βογγητα πόνου και προσπαθούσε να με κρατήσει σφιχτά.

Με την κόρη του σφιχτά στην αγκαλιά μου, έτρεξα έξω από τη σκηνή, ανάμεσα από τους άντρες που προσπαθούσαν να σβήσουν τη φωτιά. 

«Γιατρός!» φώναξα έχοντας στο μυαλό μου τη φροντίδα της μικρής. 

Εκείνη τη στιγμή άκουσα ένα πυροβολισμό από τη δεξιά μου πλευρά και πριν προλάβω να αντιδράσω, ένιωσα το κεφάλι του κοριτσιού να βαραίνει και αίμα να κυλάει στον αριστερό μου ώμο. Σταμάτησα να τρέχω και κοίταξα την κόρη του στρατηγού.

«Σκοπευτής δυτικά» φώναξε ένας στρατιώτης και όλοι έτρεξαν προς εκείνη την κατεύθυνση.

Εγώ δεν μπόρεσα. Είχα πέσει στα γόνατα και από την δεύτερη τραυματική εικόνα της ημέρας.  Το κοριτσάκι του στρατηγού με μία τρύπα στο κρανίο της. Αίμα να κυλάει και από τις δύο μεριές του κεφαλιού της στο χώμα του στρατοπέδου. Τα χέρια μου να τρέμουν, όχι από φόβο, ο στρατηγός ήταν και αυτός νεκρός οπότε δεν χρειαζόταν να τον ενημερώσω πως η κόρη του είχε πεθάνει. Εκείνη τη στιγμή με ένοιαζε λιγότερο ο στρατηγός και περισσότερο η αθώα παιδική ζωή την οποία είχε στερήσει ο όποιος σκοπευτής αντίπαλος μας.

Ευτυχώς για μένα οι άντρες μας ήταν γρήγοροι στον εντοπισμό του και στην αφόπλιση και επαναφορά του στο στρατόπεδο. Δυστυχώς για αυτόν, θα τον πήγαιναν στον στρατηγό για να αποφασίσουν τη μοίρα του. Αυτοί οι στρατιώτες δεν ήξεραν όμως πως ο στρατηγός δεν ήταν πλέον ζωντανός. 

Πλησιάζοντας την σκηνή του στρατηγού, οι άντρες μας με είδαν να στέκομαι γονατιστός με ένα νεκρό κοριτσάκι στην αγκαλιά μου. Αυτοί προσπάθησαν να με προσπεράσουν αγνοώντας όλο το σκηνικό και προσπάθησαν να σύρουν τον αντίπαλο μας προς τον στρατηγό.

«Ήταν ένα αθώο παιδί…» ξεφώνισα αλλά κανείς δεν μου έδωσε προσοχή.

«Σήκω Γιόχαν» κοίταξα πάνω και είδα τον αντιστράτηγο ο οποίος δεν είχε ίχνος συναισθήματος στο πρόσωπο του.

«Την σκότωσε…»

«Σήκω…» προσπάθησε να μου ξαναπεί ο αντιστράτηγος, αλλά εγώ σηκώθηκα και έτρεξα να προλάβω τους στρατιώτες πριν μπουν στη σκηνή του στρατηγού. «Ήταν απλά ένα αθώο παιδί» φώναξα στον αντίπαλο μας καθώς πλησίαζα τον όχλο.

«Ακολουθώ οδηγίες» μου απάντησε.

«Ήταν παιδί…» του είπα, σαν να μην το είχα επαναλάβει ήδη τρεις φορές.

«Έκανα αυτό που ήταν απαραίτητο για τον σκοπό μας»

«Ήταν παιδί» τον άρπαξα από τον γιακά του και του φώναξα στο πρόσωπο του. «Ένα αθώο παιδί που έτυχε να βρίσκεται κατά λάθος στο στρατόπεδο»

«Ηρέμησε Γιόχαν» άκουσα τον αντιστράτηγο να μου λέει από πίσω μου.

«Τη σκότωσες εν ψυχρώ» τον αγνόησα και συνέχισα να μιλάω στον άντρα. «Τι μπορεί να είναι ο σκοπός σας; Αυτό είναι έγκλημα πολέμου. Και ας εξαιρέσουμε και το γεγονός ότι μας επιτίθεσαι εκτός ώρας μάχης»

«Γιατί ρε, τι θα γίνει;» μου απάντησε θρασύτατα αυτός με τα χέρια δεμένα πίσω από την πλάτη του.

«Τι θα γίνει;» έβγαλα ένα στιλέτο που είχα στη ζώνη μου και το κάρφωσα στο αριστερό του νεφρό.

«Αρκετά» μας χώρισε ο αντιστράτηγος ρίχνοντας μας και τους δύο στο έδαφος, όμως μόνο εγώ είχα τη δύναμη να στηριχτώ στα πόδια μου ενώ ο άλλος άντρας αργοπέθαινε. «Ταξίαρχε Γιόχαν Βάγκνερ, αυτή δεν είναι συμπεριφορά υψηλόβαθμου στο στρατό μας» μου είπε και στάθηκε μπροστά μου για να μου επιβάλλει την κυριαρχία του. «Από σήμερα είστε και πάλι επιλοχίας» μου είπε και στράφηκε προς στη σκηνή του στρατηγού για να τον φωνάξει έξω, λες και αν ο στρατηγός είχε ακούσε όλη αυτή τη φασαρία και ήταν ζωντανός δεν θα είχε βγει έξω από μόνος του.

«Ο στρατηγός είναι νεκρός» του είπα και σηκώθηκα. «Αν θες πάρε τη θέση του» του είπα προσπαθώντας να μη φανώ είρωνας και μπήκα μέσα στη σκηνή με το πτώμα του κοριτσιού σηκωμένο στον ώμο μου για σηκώσω και το νεκρό σώμα του πατέρα της.\newline\newline

«Στη συνέχεια ζήτησα από δύο άντρες να με βοηθήσουν να τους θάψουμε κοντά στο στρατόπεδο και κάπως έτσι βρέθηκα από το να είμαι ένας αξιότιμος ταξίαρχος στον στρατό, στο να είμαι ένας ακόμη επιλοχίας» μου είπε ο Γιόχαν και δάγκωσε ένα κομμάτι από τη φρατζόλα που είχε αρπάξει ενώ αφυγούνταν την ιστορία.

«Μάλιστα…» του απάντησα σκεπτόμενος. «Όλα αυτά πριν δεκατρία χρόνια και δεν υπήρχε καμία ενημέρωση πέραν του γεγονότος ότι ο στρατηγός είχε πεθάνει;» τον ρώτησα.

«Θα σας προκαλούσε πολύ μεγαλύτερη αναστάτωση αν μαθαίνατε εκείνη τη στιγμή την αλήθεια» προσπάθησε να μου απαντήσει διπλοματικά.

«Έχεις δίκιο» έκοψα την συζήτηση πριν μου προκληθούν παραπάνω συναισθήματα καθώς βρισκόμασταν κοντά στο αρχοντικό.


\chapter{}


«Ήξερε καλά που βρισκόταν η κόρη του;» αιφνιδιάστηκα με τα λόγια του Άλμπερτ. «Τι εννοείς;»

«Αυτό που κατάλαβες. Ο πατέρας σου ήξερε καλά που βρισκόταν η κόρη του, για αυτό δεν έστειλε ποτέ ομάδα αναζήτησης για να την ψάξουν» μου απάντησε αυτός με περίεργο τρόπο ώστε να μη μου αποκαλύψει ακριβώς τι εννοούσε ακόμη.

«Γιατί τέτοια λόγια τέτοια ώρα Άλμπερτ; Τι σε κάνει να αναρωτιέσαι για την αδελφή μου τώρα και όχι πριν δεκατρία χρόνια που γύρισες στο αρχοντικό από το στρατόπεδο όταν ο πατέρας και η προηγούμενη μπάτλερ μας απεβίωσαν την ίδια περίοδο;»

«Επισκέφτηκα τον τάφο του πατέρα σου κοντά στο στρατόπεδο όσο εσύ και η μητέρα σου μιλούσατε με τον στρατηγό… και νομίζω ήρθε η στιγμή να σου πω την αλήθεια» είπε και στήριξε το ένα χέρι του στον αριστερό μου ώμο. 

«Ποια αλήθεια; Που βρισκόταν η αδελφή μου ενώ εμείς την αναζητούσαμε Άλμπερτ;»

«Ο πατέρας σου ήταν καλός άνθρωπος Μάρκο, το ξέρεις;» φάνηκε να προσπαθεί να αλλάξει το θέμα ο Άμπερτ.\newline\newline

Τότε, όπως ξέρεις, ήμουν λοχαγός του και κοντινός φίλος του πατέρα σου. Με είχε θέσει υπεύθυνο για τις μετακινήσεις του και εγώ φυσικά είχα δεχτεί γιατί ήταν μία τίμια και εύκολη σχετικά θέση. Ενώ οδηγούσα προς το στρατόπεδο εκείνη το μέρα, μία ώρα περίπου πριν φτάσουμε, άκουσα μία φωνή από το χώρο των αποσκευών, αλλά το αγνόησα γιατί νόμισα ότι ήταν η κούραση μετά από 5 ώρες οδήγηση. Όμως αυτός ο ήχος ακούστηκε ξανά πιο έντονα και το παρατήρησε και ο πατέρας σου αυτή τη φορά.

Σταματήσαμε το αυτοκίνητο λοιπόν και ανοίξαμε την πίσω πόρτα για να βρούμε την πηγή του ήχου. Μείναμε και οι δύο άφωνοι όταν είδαμε την αδελφή σου να κοιμάται πάνω στις αποσκευές μας και να ομιλεί στον ύπνο της. Ο πατέρας σου πετάχτηκε προς τα πίσω και ακούμπησε τις παλάμες του στο πρόσωπο του ξαφνιασμένος, αλλά όχι θυμωμένος. Εγώ προσπάθησα να σκεφτώ κάτι να κάνω εκείνη τη στιγμή αλλά δεν μπορούσα. Δεν μπορούσαμε να γυρίσουμε πίσω γιατί δεν είχαμε τα καύσιμα και τον χρόνο να κάνουμε τέτοια αλλαγή στα σχέδια. Από την άλλη είχαμε τονίσει τόσες φορές ότι θα ήταν επικίνδυνο και αφελές να πάρουμε ένα μικρό κοριτσάκι στο στρατόπεδο μαζί μας, οπότε ήμασταν σε αδιέξοδο.

«Μπαμπά» ακούσαμε τη φωνούλα της ενώ ήμασταν και οι δύο σε σύγχυση. 

«Τι κάνεις εσύ εδώ;» της είπε ο πατέρας της χωρίς κανένα μίσος στη φωνή του για το λάθος που είχε κάνει η κόρη του. «Γιατί δεν είσαι με τη μαμά και το Μάρκο;»

«Επειδή κρυβόμουν» απάντησε αυτή και απλά του χαμογέλασε γιατί ήξερε ότι όλες οι άλλες απαντήσεις που μπορούσε να δώσει θα ήταν ακόμη πιο λάθος.

Ο πατέρας σου τότε μου έριξε ένα βλέμμα λες και είχα εμπλακεί και εγώ στην εισβολή της κόρης του στο αυτοκίνητο και με είδε που τον κοίταζα περιμένοντας την εντολή του για την επόμενη κίνηση. Κατάλαβε από το άγχος μου πως ούτε εγώ είχα ιδέα για την ξαφνική συντροφιά που μας βρήκε πέντε ώρες μέσα στο ταξίδι και χαλάρωσε για να σκεφτεί. Εκείνος έκανε ένα βήμα πίσω από το αυτοκίνητο και η κόρη του εκτινάχθηκε από τις αποσκευές και προσγειώθηκε στο πόδι του για να τον αγκαλιάσει.

«Μάλλον θα πεινάς» της είπε και χάιδεψε το κεφαλάκι της.

«Πολύ» του απάντησε αυτή καταφατικά.

«Άλμπερτ, πάμε στο στρατόπεδο και θα δούμε τι μπορούμε να κάνουμε με αυτήν την ταραχοποιό που μας βρήκε» μου είπε ο κύριος και πήρε στην αγκαλιά του την κόρη του.

«Μάλιστα!» του απάντησα απλά εγώ και ξεκίνησα πάλι το αυτοκίνητο.

Τις επόμενες μέρες προσπαθούσαμε να βρούμε κάποιον τρόπο να γυρίσουμε την αδελφή σου πίσω στο αρχοντικό. Τα τρένα τα είχαν ήδη καταλάβει οι αντιστασιακοί για τους δικούς τους. Από οχήματα είχαμε αρκετά, αλλά λόγω του μειωμένου στρατού που είχαμε σε σχέση με τους αντιπάλους μας, δεν είχαμε την πολυτέλεια να μπορέσουμε να ξοδέψουμε ούτε έναν στρατιώτη για αυτή την αποστολή. Τόσο σοβαρή ήταν και είναι ακόμη η διαφορά των αριθμών των δύο στρατών. Η μόνη διαφορά είναι πως εμείς είχαμε μέχρι στιγμής λίγο πιο έξυπνους αρχηγούς που ηγήθηκαν συντηρητικά και στρατηγικά, σε σχέση με τους αντιστασιακούς.

Αυτό δεν σημαίνει βέβαια πως δεν είχαμε και εμείς τις απώλειες μας. Έναν μήνα μετά την άφιξη μας στο στρατόπεδο και μετά από πολλές μάχες, είχαμε αποδεχτεί και οι δυο μας πως η αδελφή σου θα έμενε μαζί μέχρι να υποχωρήσει κάποιος από τους δύο στρατούς. Ο πατέρας σου όμως, παρά τις αμέτρητες προειδοποιήσεις που του είχα κάνει να μην τριγυρνάει με την κόρη του στο στρατόπεδο, αυτός συνέχιζε να την εκθέτει σε κίνδυνο από πιθανό αιφνιδιασμό των αντιστασιακών. Μου έλεγε πως δεν πρόκειται οι αντιστασιακοι να επιτεθούν στο στρατόπεδο τους εκτός μάχης, πράγμα που αποδείχτηκε λάθος μία μοιραία μέρα, μία μέρα που με τραυμάτισε.

Εκείνη τη μέρα ο στρατηγός και η κόρη του είχαν βγει εκείνο το πρωί σχετικά νωρίς από την σκηνή τους για να συναντηθούν με δύο αξιωματικούς που είχαν νέα για αυτόν. Εγώ τους άφησα έξω από τη σκηνή και τους αποχαιρέτησα, συνέχισα όμως να κοιτάζω προς το μέρος τους καθώς αυτοί απομακρύνονταν. Το τραγικό είναι ότι δεν πρόλαβαν να φτάσουν πολύ μακριά. Περίπου 20 μέτρα μακριά μου πρέπει να βρίσκονταν όταν παρατήρησα μία χειροβομβίδα που από το πουθενά έσκασε στα πόδια του στρατηγού και αυτός αντανακλαστικά πήρε αγκαλιά την κόρη του και άρχισε να τρέχει… για μισό δευτερόλεπτο πριν σκάσει η χειροβομβίδα. Ο στρατηγός πετάχτηκε στον αέρα γυρνώντας την πλάτη του όσο μπορούσε προς το έδαφος ώστε να πέσει ανάσκελα και να σώσει την κόρη του από την πρόσκρουση στο δάπεδο.

Πλησίασα όσο πιο γρήγορα και επιφυλακτικά μπορούσα, ψάχνοντας για εχθρικό στρατιώτη από την αντίθετη μεριά της τροχιάς της χειροβομβίδας. Όταν έφτασα δίπλα στον κύριο, αυτός ήταν κατάμαυρος, με το δέρμα στο πίσω μέρος των ποδιών και του κάτω μέρος της πλάτης του να είχε ξεφλουδίσει. 

«Μπαμπά, γιατί… δεν κουνιέσαι;» είδα την μικρή να προσπαθήσει να ξυπνήσει μάταια τον πατέρα της εστιάζοντας το βλέμμα της μόνο στο πρόσωπο του, το οποίο είχε χάσει όλη τη ζωντάνια από πάνω του.

«Όχι, όχι, όχι» αγχώθηκα και κάθισα δίπλα στον κύριο, ακουμπώντας το χέρι μου στον λαιμό του για να δω αν είχε παλμό και παράλληλα ψάχνοντας για τον δράστη. «Μη φοβάσαι» προσπάθησα να καθησυχάσω την αδελφή σου «όλα θα πάνε καλά» 

Δυστυχώς για μένα, ο πατέρας σου ήταν νεκρός. Ευτυχώς για την κατάσταση, είδα μία λάμψη ανάμεσα σε δύο θάμνους και υπέθεσα πως είναι σκοπευτής. Δυστυχώς για όλους μας, η μικρή προσπάθησε να σηκώσει τον πατέρα της τη στιγμή που τράβηξα το όπλο μου για να πυροβολήσω προς το μέρος του σκοπευτεί και αυτός πρόλαβε να ρίξει τη χαριστική βολή. Είχε αστοχήσει τελείως το νεκρό πτώμα του στρατηγού, που μάλλον ήθελε να πετύχει για να σιγουρευτεί πως ήταν νεκρός, και είχε πετύχει την αδελφή σου στο πλάι του κρανίου της.

Ήξερα ότι είχα “αποτύχει” όταν άκουσα τον ήχο που έκανε η σφαίρα καθώς έκοβε τον άνεμο και ένιωσα ένα βάρος στα πόδια μου, συνοδευόμενο από αίμα να κυλάει στις μπότες μου.

«Γιατρός!» φώναξα πριν κοιτάξω κάτω, αλλά όλες οι προσπάθειες μου ήταν μάταιες.

Η κόρη του στρατηγού ήταν νεκρή. Η εικόνα του νεκρού πατέρα σου και της αθώας αδερφής σου στα πόδια μου με έριξε στα γόνατα. Σήκωσα το κεφάλι της για να δω το πρόσωπο της ακόμη μία φόρα αλλά αυτό έκανε την κατάσταση χειρότερη. Τα χέρια μου έτρεμαν. Όχι από φόβο, ο κύριος ήταν νεκρός, δεν χρειαζόταν να του παραδώσω κάποια κακά νέα. Εκείνη τη στιγμή ε ένοιαζε λιγότερο ο στρατηγός και περισσότερο η αθώα παιδική ζωή την οποία είχε στερήσει ο όποιος σκοπευτής αντίπαλος μας.

«Που είναι ο αντίπαλος;» με ρώτησε ένας στρατιώτης καθώς εγώ ήμουν πεσμένος στα γόνατα και προσπαθούσα να συγκρατήσω τα συναισθήματα μου.

Δεν συγκρατήθηκα όμως για πολύ. Με το που με ρώτησε σηκώθηκα και έτρεξα προς την περιοχή στην οποία είχα πυροβολήσει. Εκεί βρήκα έναν αντιστασιακό να προσπαθεί να κλείσει την πληγή στον αριστερό του ώμο που του είχα προκαλέσει. Όταν πλησίασα αυτός προσπάθησε να με ακινητοποιήσει σημαδεύοντας το πιστόλι του προς το μέρος μου αλλά εγώ του κλώτσησα το χέρι και έριξα το όπλο του κάτω αφοπλίζοντας τον.

«Ήταν απλά ένα αθώο παιδί…» του είπα.

«Ακολουθώ εντολές» μου απάντησε.

«Ήταν παιδί…» του ξαναείπα μήπως δεν με άκουσε την πρώτη φορά.

«Έκανα αυτό που ήταν απαραίτητο για…»

«Ήταν παιδί» του φώναξα στο πρόσωπο καθώς τον άρπαξα από τον γιακά του για να τον φέρω πιο κοντά. «Τι δεν καταλαβαίνεις;»

«Λοχαγέ» άκουσα μερικούς άντρες μας να πλησιάζουν.

«Τη σκότωσες εν ψυχρώ» τους αγνόησα και συνέχισα να μιλάω στον άντρα. «Τι μπορεί να είναι αυτό που ήταν τόσο υπεράνω ώστε να σου δίνει το δικαίωμα να αφαιρέσεις τη ζωή ενός παιδιού; Και ας εξαιρέσουμε το γεγονός ότι μας επιτίθεσαι εκτός ώρας μάχης»

«Γιατί εσείς πως πήρατε τη ζωή του στρατηγού μας;» μου απάντησε ο άντρας προκλητικά. «Τώρα που υπήρξε ανταπόδοση…»

Δεν πρόλαβε να τελειώσει την πρόταση του πριν του καρφώσω το στιλέτο μου στο δεξί του νεφρό.

«Λοχαγέ τι κάνεις;» μου είπε ένας στρατιώτης από αυτούς που με ακολούθησαν. «Το ξέρεις ότι οι ανώτεροι θα θυμώσουν που πήρες τη ζωή εχθρού εκτός ώρας μάχης αντί να τον αιχμαλωτίσεις για να εξάγουμε πληροφορίες»

«Παραιτούμαι» του απάντησα μονολεκτικά.

«Τι;»

«Δεν με ενδιαφέρει άλλο ο πόλεμος»

«Και ποιος θα είναι υπεύθυνος για τις μετακινήσεις του…»

«Ο στρατηγός είναι νεκρός. Δεν είδες το πτώμα του;» του είπα και σηκώθηκα τραβόντας το στιλέτο μου από το νεφρό του άντρα. «Θα επιστρέψω στο αρχοντικό του στρατηγού. Αν θέλεις τη θέση μου πάρτη» του είπα και έσκισα το διακριτικό βαθμού μου και του έδωσα.\newline\newline

«Στη συνέχεια γυρίσαμε στο στρατόπεδο για να ενημερώσουμε τους υπόλοιπους για την αιφνιδιαστική επίθεση της οποίας ήμασταν θύματα και ως τελευταία διαταγή φρόντισα να θαφτεί ο πατέρας και η αδελφή σου σε ένα σημείο λίγο πιο έξω από το στρατόπεδο. Ήξερα πως κάποια στιγμή, με τον έναν τρόπο ή με τον άλλον θα ερχόσουν και εσύ στο στρατόπεδο. Ήθελα να το κρατήσω μυστικό όσο μπορούσα από εσένα και τη μητέρα σου»

«Γιατί;» τον ρώτησα.

«Επειδή θα δημιουργούσε αναστάτωση στην οικογένεια και ίσως πυροδοτήσει τον θυμό που μπορεί να είχες για τους αντιστασιακούς και να σε οδηγούσε τη ζωή σου σε κακό μονοπάτι. Περίμενα μία κατάλληλη στιγμή ηρεμίας για να σου το ανακοινώσω όσο πιο ομαλά μπορούσα»

«Μάλιστα…» του απάντησα σκεπτόμενος. «Δεκατρία χρόνια λοιπόν το κράταγες μέσα σου;»

Αυτός μου έγνεψε θετικά με το κεφάλι του.

«Είναι κάτι άλλο που θα ήθελες να γνωρίζεις;»

«Γιατί μερικές φορές μου μιλάς στον ενικό και άλλες φορές στον πληθυντικό;»

Αυτός μου χαμογέλασε. 

«Χαίρομαι που τα νέα δεν σε επηρέασαν τόσο όσο νόμιζα» είπε και σηκώθηκε, αφήνοντας με μόνο μου στην τραπεζαρία.

Είχα όμως άδικο. Τα νέα… τα παλιά; Η ανακοίνωση που μου έκανε όχι μόνο με είχε επηρεάσει, μου είχε όντως προκαλέσει μία δόση θυμού προς τους αντιστασιακούς που δύσκολα θα έσβηνε. Τα συναισθήματα μου δεν ήταν τόσο έντονα που θα με έκαναν να ενταχθώ στον στρατό, όμως ποιος ξέρει τι θα έκανα σε έναν αντιστασιακό αν τον είχα μπροστά μου.


\chapter{}


«Είστε έτοιμοι;» μας είπε ο Ράινερ καθώς τράβηξε το σακίδιο του προς το μέρος του.

«Εμένα γιατί με ρωτάς;» του απάντησα. «Στον Γιόχαν θα έπρεπε να απευθύνεσαι»

«Είμαι έτοιμος φίλε μου» είπε και τέντωσε το πόδι του προς το μέρος του Ράινερ, ισιώνοντας το ώστε η πατούσα και η γάμπα το να βρίσκονται στην ίδια ευθεία.

«Ωραία λοιπόν» είπε ο Ράινερ βγάζοντας το μαχαίρι του από το σακίδιο του και πλησιάζοντας τον. «Τρία, δύο…»

Και με το ένα έχωσε το μαχαίρι του με δύναμη στην πατούσα του ευάλωτου Γιόχαν. 

Πριν βγάλει το μαχαίρι, προσπάθησα να τυλίξω σφιχτά επιδέσμους και από τις δύο μεριές της πληγής ώστε να με την εκχώρηση του μαχαιριού να μην προλάβει να χάσει πολύ αίμα ο Γιόχαν και να αποφύγουμε την αποτυχία πριν καν ξεκινήσουμε την πρώτη φάση του “Σχεδίου”. Ο Ράινερ στη συνέχεια αφαίρεσε προσεκτικά το μαχαίρι του, δίνοντας χρόνο στον Γιόχαν να προσαρμοστεί, ο οποίος δεν είχε βγάλει άχνα όλη αυτή την ώρα. Βγάζοντας τελείως το μαχαίρι του, ο Ράινερ έκλεισε την πληγή βιαστικά με μερικές στρώσεις παραπανήσιου επιδέσμου που τράβηξε από το ρολό μου όσο εγώ τύλιγα και με άφησε να ολοκληρώσω τις στρώσεις καλύμματος γύρω από την πληγή και άφησε το χέρι του μόνο όταν ξεκίνησα να τυλίγω την περιοχή πάνω από την πληγή.

«Ωραία» είπε ο αντισυνταγματάρχης και σηκώθηκε, αφήνοντας με να ολοκληρώσω την εφαρμογή των λίγων πρώτων βοηθειών που ήξερα πάνω στον Γιόχαν. «Να σας υπενθυμίσω τα τρία σινιάλα ή τα θυμάστε;»

«Χαζοί νομίζεις ότι είμαστε;» του απάντησε ο Γιόχαν, με μία δόση πόνου στη φωνή του.

«Συγγνώμη συνταγματάρχη, δεν ήθελα να σε προσβάλω» του απάντησε και σήκωσε το σακίδιο του για να το βάλει στους ώμους του. «Κοιμηθείτε εδώ το βράδυ. Είναι αρκετά κοντά στο αρχοντικό για να μπορεί να σε κουβαλήσει ο Μάρκο και αρκετά μακριά ώστε να μπορεί να αφήσετε τα πράγματα σας στους θάμνους και κανείς από την πόλη να μην τα βρει. Φροντίστε να κοιμηθείτε με τα καλά σας ρούχα στο χώμα ώστε να είστε πιο πειστικοί»

«Μάλιστα» του απάντησα ολοκληρώνοντας το δέσιμο του ποδιού του συνταγματάρχη.

«Ελπίζω να σας ξαναδώ» ήταν οι τελευταίες λέξεις του Ράινερ πριν μας γυρίσει την πλάτη του για να προχωρήσει προς την πόλη.

Καθισμένοι γύρω από τη φωτιά, μόνο ένα πράγμα μπορούσα να σκεφτώ εκείνη τη στιγμή και δεν είχε καμία σχέση με το “Σχέδιο”. Πώς γίνεται ο Γιόχαν να μην έβγαλε ούτε μία κραυγή πόνου όταν τον τραυμάτισε ο Ράινερ, παρόλο που προφανώς ένιωσε πόνο; Στη θέση του οριακά θα είχα πεθάνει από τον πόνο και η αποστολή μας θα είχε τελειώσει πριν καν αρχίσει. Πρέπει να σταματήσω να το λέω αυτό βέβαια. Όσο σκέφτομαι αρνητικά θα μου μένει στο μυαλό το πεσιμιστικό αποτέλεσμα και θα αποτύχουμε όντως. Δεν θα έκανε κακό όμως να τον ρωτήσω από περιέργεια.

«Γιόχαν;»

«Ναι Μάρκο;» γύρισε το βλέμμα του προς το μέρος μου.

«Πώς γίνεται να μην έβγαλες άχνα όταν σε μαχαίρωσε ο Ράινερ;»

«Αυτό είναι θέμα εμπειρίας μικρέ. Είναι και ο λόγος που είμαι εδώ μαζί σου για αυτή την αποστολή» είπε και ξάπλωσε. «Έχω τραυματιστεί θανάσιμα αρκετές φορές τα τελευταία χρόνια, αυτή τη πληγούλα δεν μπορεί να συγκριθεί με τις υπόλοιπες. Αν ο στρατηγός διάλεγε κάποιον σαν εσένα να πάρει τη θέση μου θα ήταν αλλιώς τα πράγματα»

«Σαν εμένα;» ξαφνιάστηκα με τα λόγια του Γιόχαν και μου ήρθε στο μυαλό η μέρα που ήμασταν στοιχισμένοι στο στρατόπεδο που στεγαζόμουν και ο στρατηγός φώναξε το όνομα μου και μου είπε να ακολουθήσω τον συνταγματάρχη και τον αντισυνταγματάρχη. Αφού ο στρατηγός με είχε επιλέξει για την αποστολή, για άλλο ρόλο μεν, αλλά το γεγονός ότι βρισκόμουν εκεί σήμαινε ότι είχα την εμπιστοσύνη του.

«Συγγνώμη, δεν ήθελα να σε προσβάλλω» είπε ο Γιόχαν λες και διάβαζε τις σκέψεις μου «ήθελα απλά να πω πως κάποιος νεότερος στη θέση μου ίσως τα έκανε μαντάρα»

«Συμφωνώ» επιβεβαίωσα τα λόγια του και έπεσα προς τα πίσω για να ξαπλώσω. «Τα λέμε το πρωί λοιπόν;»

«Αν μπορέσω να κοιμηθώ» μου απάντησε ο συνταγματάρχης και έκλεισε τα μάτια του.

Η νύχτα πέρασε πολύ εύκολα παραδόξως. Παρόλο που κοιμόμασταν στο χώμα του δάσους, κανείς από τους δύο μας δεν παραπονέθηκε καθώς ήταν πιο εύκολο από ότι περιμέναμε. Προφανώς είχαμε φορέσει τα καλά μας ρούχα, όπως μας είπε ο Ράινερ,  ώστε να μπορέσουν να λερωθούν το βράδυ καθώς θα κυλιόμασταν ασυναίσθητα στο χώμα. Ήταν μία από τις έξυπνες ιδέες που είχα κατά τη διάρκεια του ταξιδιού για να φανούμε πιο πειστικοί κατά την άφιξη μας στο αρχοντικό. Είχα χαρεί όταν είδα τις εκφράσεις στα πρόσωπα του συνταγματάρχη και του αντισυνταγματάρχη όταν τους πρότεινα την ιδέα μου και αυτοί αναρωτήθηκαν πως δεν το είχαν σκεφτεί οι ίδιοι νωρίτερα.

Το πρωί ξύπνησα δεύτερος. Ο Γιόχαν είχε ξυπνήσει πρώτος, αλλά με το τραυματισμένο πόδι του δεν μπορούσε να κάνει και πολλά. Όχι ότι χρειαζόταν κιόλας, αφού είχα σβήσει τη φωτιά πριν κοιμηθούμε για να μην κινήσουμε υποψίες.

«Καλημέρα» μου είπε και έσπρωξε το σακίδιο του κάτω από τον θάμβο πίσω του. «Πήρες τον χρόνο σου χωρίς το κοκόρι τον Ράινερ, ο ήλιος είναι σχεδόν ενενήντα μοίρες» σχολίασε.

«Συγγνώμη» του είπα και σηκώθηκα να μαζέψω τα πράγματα μου πριν αποχωρήσουμε.

«Δεν πειράζει, έχουμε χρόνο. Μη βιάζεσαι» μου είπε αυτός χαλαρός για να με καθησυχάσει.

Του χαμογέλασα για να καταλάβει ότι τον άκουσα και επιβράδινα τις κινήσεις μου. Δεν είχα όρεξη να φάω πρωινό, το άγχος με είχε καταβάλει. Είχα πει στον εαυτό μου πως δεν θα έκανα άλλες αρνητικές σκέψεις συνειδητά, όμως ασυνείδητα δεν μπορούσα να το ελέγξω μέχρι όλος μου ο εαυτός να πιστέψει στο “Σχέδιο” και να σταματήσει να βρίσκει μικρές τρύπες που μπορεί να μην επηρεάζουν καν το τελικό αποτέλεσμα.

«Όποτε είσαι έτοιμος πες μου» είπε ο Γιόχαν και τέντωσε τα χέρια του προς τα πίσω για να πάρει μια πιο ανακουφιστική στάση.

«Είμαι έτοιμος» του είπα και πέταξα τα όπλα μας και το δικό μου σακίδιο στον θάμνο που είχε κρύψει και αυτός τις δικές του προμήθειες.

«Ωραία λοιπόν» μου απάντησε και μου έκανε νεύμα να τον πλησιάσω για να τον σηκώσω. 

Όταν φύγαμε, φροντίσαμε να πάρουμε μαζί μας μόνο τα απαραίτητα για να μη μας βαραίνει τίποτα. Δεν περπατήσαμε πάνω από είκοσι λεπτά για να φτάσουμε στον στόχο μας. Με το που τελείωσε η πρασινάδα του δάσους στο μονοπάτι που είχαμε ακολουθήσει, βρεθήκαμε στην πίσω μεριά ενός τεράστιου κτιρίου περιτρυγιρισμένο από έναν κήπο ο οποίος έμοιαζε πάρα πολύ με λαβύρινθο. Το κτίριο αυτό ταίριαζε την περιγραφή που μας είχε δώσει ο στρατηγός, οπότε ήμασταν σίγουροι πως βρισκόμασταν στο σωστό μέρος.

Ο συνταγματάρχης βάσταξε βλέποντας την πόρτα, καθώς ήταν ανυψωμένη και θα έπρεπε όταν φτάσουμε να ανεβούμε πέντε σκαλιά. Ευτυχώς για μας, ο δρόμος για να φτάσουμε στην πίσω πόρτα του αρχοντικού ήταν ένα απλό πέτρινο μονοπάτι, χωρίς φυτά του κήπου να μας μπλοκάρουν ενώ περπατούσαμε σε αυτό. Ακολουθώντας αυτό το μονοπάτι, παρατήρησα από μακριά μας έναν νεαρό άντρα, ίσως συνομήλικος με εμένα, να σκουπίζει πεσμένα φύλλα στην δεξιά μεριά του κήπου και να ρίχνει πεταχτές ματιές προς ένα παράθυρο στο πλάι του αρχοντικού. Πριν προλάβω όμως να του φωνάξω κάτι, τα βλέμματα μας διασταυρώθηκαν. Αυτός, μάλλον από φόβο, έριξε κάτω τη σκούπα του και έτρεξε προς την αντίθετη μεριά.

«Μη τρέχεις» προσπάθησα να του φωνάξω, αλλά μάλλον ο δεν με άκουσε ο συνομήλικος μου.

«Παράτα τον, δεν πειράζει» μου είπε ο συνταγματάρχης και συνεχίσαμε να περπατάμε προς την πόρτα.

Φτάνοντας στις σκάλες, προσάρμοσα τον συνάδελφο μου καλύτερα στον ώμο μου και αρχίσαμε να ανεβαίνουμε τα σκαλοπάτια αργά και σταθερά. ‘Ενα σκαλοπάτι ο συνταγματάρχης με το καλό του πόδι, ένα εγώ για στήριγμα.

«Ακίνητοι» άκουσα μία αντρική φωνή από πίσω μου καθώς προετοιμαζόμασταν να ανεβούμε το προτελευταίο σκαλί.

Εμείς δεν είπαμε τίποτα και σηκώσαμε τα χέρια μας ψηλά αντανακλαστικά όταν ακούσαμε τον άντρα πίσω μας. Προσπαθώντας να κρατάω και τον συνταγματάρχη, γύρισα προς το μέρος του άντρα πίσω μας και είδα τον νεαρό κηπουρό να μας σημαδεύει με ένα πιστόλι.

«Τι κάνετε εδώ;» μας ρώτησε αυτός χωρίς να κατεβάσει το όπλο του.

«Είμαστε θύματα του πολέμου» του απάντησα. «Χάσαμε το σπίτι μας στην πολιορκία που έγινε στην πόλη δίπλα από τη δική σας και προσπάθησα να κουβαλήσω τον γέρο μου μέχρι εδώ για να μπορέσει να τον βοηθήσει κάποιος»

«Ποια πόλη δίπλα από τη δική μας;» με ρώτησε αυτός λες και δε γνώριζε πως υπάρχει μόνο μία πόλη με μεγαλοαστούς κοντά στη δική του που είχε πολιορκηθεί πρόσφατα από τους αντιστασιακούς.

«Από την…»

«Τι γίνεται εδώ;» άνοιξε ένας καλοντυμένος άντρας την πόρτα και είδε την κατάσταση στην οποία βρισκόμασταν. «Κυρία, μη πλησιάσετε» είπε ο άντρας και έκανε νεύμα με το χέρι του σε μία μεσήλικη γυναίκα η οποία βρισκόταν μερικά μέτρα πιο πίσω του.

«Είναι εχθροί Άλμπερτ» είπε ο νέος άντρας σφίγγοντας το όπλο του για να φτιάξει τον στόχο του.

«Όχι» είπε ο δεύτερος άντρας και πλησίασε τον συνταγματάρχη. «Είναι τραυματισμένοι» είπε και με βοήθησε και ξαπλώσουμε τον συνάδελφο μου κάτω. «Από που είστε;» τους ρώτησε. 

«Από την διπλανή πόλη. Έφερα τον γέρο μου ως εδώ μήπως υπάρχει περίπτωση να μας βοηθήσει κάποιος» επανέλαβα το ψέμα που είχαμε προσυμφωνήσει.

«Μάλιστα» μου απάντησε αυτός και κοίταξε το πόδι του συνταγματάρχη καλύτερα, αλλά μάλλον δεν ήταν γιατρός γιατί δεν προσπάθησε να ξετυλίξει τον επίδεσμο για να δει την πληγή που υπήρχε από κάτω. Ακούμπησε απλά το δικό του χέρι πάνω στο πόδι του συναδέλφου μου, στο οποίο είχε δεμένο με μία άσπρη πετσέτα τον δείκτη του, η οποία είχε σημάδια αίματος.

«Λέει ψέματα» φώναξε ο νέος στον μεγαλύτερο του. «Προσπαθεί να σε πείσει να τους βοηθήσεις Άλμπερτ»

«Τι σε κάνει να το λες αυτό;» είπε ο Άλμπερτ κοιτώντας τον συνταγματάρχη στα μάτια βλέποντας ένα βλέμμα τρόμου. «Οι δύο άντρες αυτοί δεν φαίνονται για αντιστασιακοί» του απάντησε και σηκώθηκε.

«Τι λες Άλμπερτ;» είπε και με στόχευσε στο κεφάλι. «Αν… αν δεν ήταν, γιατί ήρθαν από το δάσος και όχι από την πόλη;»

«Γιατί, όπως μας είπαν, κατάγονται από την πόλη βορειοανατολικά της δικής μας»

«Και η πληγή του. Πως εξηγείς…»

«Αρκετά» ήρθε προς το μέρος μας η μεσήλικη γυναίκα και πλησίασε και αυτή τον συνταγματάρχη χαμογελώντας.

«Κυρία;» είπε ο Άλμπερτ.

«Μαμά;» παραξενεύτηκε ο νεότερος άντρας.

«Μου ήρθε μία ιδέα. Τι και αν παίρναμε αυτούς τους δύο άντρες υπό την προστασία μας» είπε και σηκώθηκε να κοιτάξει τους δύο συγκάτοικους της. «Θα βοήθαγε πολύ την εικόνα μας ως αριστοκρατική οικογένεια, δεν νομίζετε; “Η οικογένεια Μπερνούλι έσωσε έναν νέο άντρα και τον ηλικιωμένο παππού του”. Φαντάζομαι ήδη τα νέα να μεταδίδονται στην πόλη μας. Αν φτάσουν τα καλά λόγια στα αυτιά του στρατηγού μπορεί να αλλάξει η γνώμη του στρατηγού, ε; Τι πιστεύετε;»

«Δεν νομίζω πως δουλεύουν έτσι οι βαθμίδες του στρατού…» της είπε ο Άλμπερτ «…όμως θα μπορούσαμε να τους βοηθήσουμε μέχρις ότου να μπορέσουν να επανενταχθούν στην κοινωνία»

«Ναι… να τους βοηθήσουμε» του είπε η “κυρία” και του έκλεισε το μάτι προσπαθώντας να είναι διακριτική αλλά το υπόλοιπο σώμα της την πρόδωσε.

«Τι;» παραξενεύτηκε ο Άλμπερτ.

«Εσύ τι πιστεύεις;» απευθύνθηκε η “κυρία” στον νέο άντρα με το πιστόλι.

Αυτός έμεινε ακίνητος χωρίς το βλέμμα του να χάσει τον στόχο του.

«Τι σκέφτεσαι;» τον ρώτησε αυτή.

Ο άντρας δεν φάνηκε να αντιδρά στην φωνή της μητέρας του. Ο Άλμπερτ τον πλησίασε για να τον ηρεμήσει και άγγιξε το μπράτσο του νεότερου άντρα που. Φάνηκε σαν να του ψιθυρίζει κάτι για αρκετή ώρα, αλλά δεν μπόρεσα να τον ακούσω καθαρά. Μετά από δύο ή τρία δευτερόλεπτα, ο νέος χαμήλωσε το όπλο του.

«Καλά» είπε, έδωσε το πιστόλι του στον Άλμπερτ και μας προσπέρασε, μπαίνοντας μέσα στο αρχοντικό.

«Ωραία, λοιπόν» μας είπε η γυναίκα για να μας τραβήξει την προσοχή πάνω της «καλωσορίσατε στο αρχοντικό μας»


\chapter{}


«Δεν είσαι λίγο μεγάλος για να παίζεις με στρατιωτάκια;» μου είπε ο Άλμπερτ.

«Αυτό λες κάθε φορά» του είπα και συνέχισα να κινώ τα στρατιωτάκια μου στον χειροποίητο χάρτη που είχα φτιάξει ως ταμπλό για το παιχνίδι μου ώστε να είναι πιο ρεαλιστικό.

«Θέλεις κάτι να φας τουλάχιστον;» με ρώτησε πριν κλείσει την πόρτα του δωματίου μου.

«Ναι αμέ» του απάντησα και στράφηκα πίσω στο παιχνίδι μου για να τελειώσω το φανταστικό σενάριο που προσπαθούσα να εξελίξω.

«Θα σε φωνάξω κάτω όταν έχω κάτι έτοιμο» μου χαμογέλασε και έκλεισε την πόρτα.

Από χθες που μου ανακοίνωσε τον θάνατο της αδελφής μου, η συμπεριφορά του είχε αλλάξει. Ήταν πιο χαλαρός με εμένα, λες και είχε φύγει ένα βάρος από πάνω του. Άραγε η μητέρα γνωρίζει ή βρίσκεται στην ίδια κατάσταση που βρισκόταν και ο εαυτός μου πριν εικοσιτέσσερις ώρες; Είχε δίκιο που μας το έκρυψε ο Άλμπερτ, η διαταραχή ήταν αναπόφευκτη όσο ομαλά και να μου το έλεγε. 

«Αχ, δεν μπορώ να συγκεντρωθώ» είπα και σηκώθηκα για να κατέβω στην τραπεζαρία.

Κοιτώντας τον χάρτη μου πάλι, το βλέμμα μου έπεσε στο στρατόπεδο στις πρώτες γραμμές, στο οποίο είχα προσθέσει μία επιπλέον ταφόπλακα δίπλα από αυτή του πατέρα μου με το όνομα της αδελφής μου. Είχα τοποθετήσει έναν στρατιώτη κατά τύχη σε θέση που φαινόταν σαν να την φύλλαγε από τους αντιστασιακούς. Η ειρωνεία με έκανε να αφήσω ασυνείδητα ένα γέλιο πριν αποχωρήσω από το δωμάτιο μου.

Κατεβαίνοντας τις σκάλες, άκουσα την φωνή του Άλμπερτ από την κουζίνα.

«Α στο καλό» φώναξε και έτρεξα να δω τι συνέβει.

Μπαίνοντας στην κουζίνα είδα τον Άλμπερτ να λυγίζει το σώμα του πάνω από τον πάγκο της κουζίνας και να προσπαθεί να τυλίξει τον δείκτη του πρόχειρα με την πρώτη πετσέτα που βρήκε. Από το δάκτυλο του έσταζε αίμα, όχι πολύ, αλλά αρκετό για να λερώσει την σανίδα κοπής που χρησιμοποιούσε. Λίγο πιο δίπλα είχε ακουμπήσει τον ένοχο για τον τραυματισμό του, το μαχαίρι με το οποίο είχε ξεκινήσει να κόβει τα φρούτα του γεύματος μου.

«Α… Μάρκο» μου είπε και επιτάχυνε το τύλιγμα του δακτύλου του.

«Μη σε νοιάζει» του απάντησα και σήκωσα το μαχαίρι από τον πάγκο για να του ρίξω λίγο νερό ώστε να φύγει το αίμα από πάνω του.

«Να πάρει η ευχή» είπε ο Άλμπερτ τελειώνοντας το δέσιμο του χεριού του και έπιασε μία ακόμη πετσέτα για να καθαρίσει το χαμό που προκάλεσε.

«Μη σε νοιάζει το γεύμα» του είπα και ακούμπησα το μαχαίρι στο συρτάρι μαζί με τα υπόλοιπα.

«Δεν είναι μόνο αυτό» μου είπε αυτός «πρέπει να κάνω κι άλλες δουλειές σήμερα, όπως να καθαρίσω τον κήπο, να…»

«Θες να σε βοηθήσω;» τον διέκοψα και είδα ένα χαμόγελο να αστράφτει στο πρόσωπο του.

«Δεν μπορείς Μάρκο, δεν είναι δουλειά για σένα» μου απάντησε αυτός κρατώντας το δάχτυλο του.

«Έλα μωρέ, μία φορά στο τόσο έχω την ευκαιρία να κάνω κάτι στο σπίτι» του είπα. «Μας έχεις καλομαθημένους, η αλήθεια είναι»

«Και… τι θες να κάνεις;»

«Δεν είπες ότι θες να καθαρίσεις τον κήπο; Ε, δώσε μου την σκούπα να μαζέψω τα φύλλα να είναι όλα σε ένα σημείο και όταν κλείσει η πληγή σου σε λίγες ώρες θα συνεχίσεις εσύ με τις δουλειές σου. Εξάλλου τι άλλο έχω να κάνω όλη μέρα;»

Αυτός με κοίταξε με ένα περίεργο βλέμμα λες και δεν είχα ξαναπροσφερθεί να τον βοηθήσω. Αυτός πάντα χρησιμοποιούσε την ίδια δικαιολογία, πως η δουλειά του ως μπάτλερ μας ήταν να μας βοηθάει με τις εργασίες στο αρχοντικό ώστε να μπορούμε να ζούμε πιο άνετα, που είναι μεν ο λόγος που τον προσλάβαμε αρχικά, αλλά μετά από τα χθεσινά νομίζω υπάρχει και κάποιος βαθύτερος λόγος που δεν με αφήνει ποτέ να τον βοηθήσω, παρά την πρόθεση μου. Τώρα που είχε τραυματιστεί, όχι τρελά βέβαια, ήταν η ευκαιρία μου να τον πείσω.

«Να πάω λοιπόν;» τον ρώτησα καθώς αυτός δεν μου είχε απαντήσει.

«Ναι βεβαίως. Άλλαξε, βέβαια, ρούχα πρώτα. Δεν θα ήθελες να λερώσεις τα καλά σου» μου είπε και έφτιαξε τη στάση του σώματος του. «Η σκούπα είναι κοντά στην καλύβα στην αριστερή μεριά του κήπου, εκεί που φυλάω και τα υπόλοιπα όργανα κηπουρικής και τα κρυφά όπλα για περίπτωση ανάγκης» μου είπε και περπάτησε προς την είσοδο του σπιτιού. «Εγώ θα βρίσκομαι στο παράθυρο του διαδρόμου από εκείνη τη μεριά να σου δίνω οδηγίες, εντάξει;»

«Αριστερά όπως βγαίνω ή όπως κοιτάς το σπίτι;»

«Όπως βγαίνεις δεξιά» μου απάντησε.

«Και που θα βρω ρούχα για να αλλάξω;»

Αυτός ξεφύσηξε και μου έκανε νεύμα να τον ακολουθήσω στο δωμάτιο του, όπου μου έδωσε μία απλή μπλούζα και μία σαλοπέτα για να φορέσω όσο θα σκούπιζα τα φύλλα. Στην αρχή δυσκολεύτηκα να βρω πως ακριβώς θα φόραγα τα ρούχα που μου έδωσε, αλλά όταν ο Άλμπερτ με βοήθησε έβγαλε νόημα. 

Οι δυσκολίες μου συνέχισαν και όταν κατέβηκα στον κήπο και άρχισα να σκουπίζω. Τόσες φορές είχα δει τον Άλμπερτ να κάνει ακριβώς την ίδια δουλειά αλλά δεν μπορούσα να καταλάβω το πως. Πως έπρεπε να πιάσω την σκούπα, με το καλό μου χέρι ψηλά και το αντίθετο χαμηλά; Ο Άλμπερτ μου έλεγε το αντίθετο, αλλά αυτό μου φαίνεται περίεργο. Επίσης, προς τα που έπρεπε να σκουπίζω τα φύλλα ώστε να τα σύρω όλα σε ένα σημείο. Προσπάθησα να τα τραβάω προς το μέρος μου αλλά ο Άλμπερτ μου φώναζε πως θα λέρωνα τα πόδια μου έτσι και πως έπρεπε να στέκομαι στο πλάι και να σπρώχνω αντί να τραβάω με την σκούπα. Η εμπειρία χρόνων προσπαθούσε λεκτικά να εξηγήσει σε έναν οριακά άχρηστο στο να συγχρονίζει το σώμα του, αλλά δεν μπορούσαμε να συνεργαστούμε.

Η ασυνεννοησία συνέχισε μέχρις που είδα δύο άντρες, έναν νέο, μάλλον κοντά στην ηλικία μου, και έναν γέρο να πλησιάζουν την πίσω πόρτα του σπιτιού μας. Κοιτώντας προς το μέρος τους, τα μάτια μου και του συνομηλίκου μου διασταυρώθηκαν. Εκείνη τη στιγμή, το πρώτο πράγμα που σκέφτηκα ήταν πως αυτοί ήταν πιθανοί εισβολείς. Άφησα ό,τι έκανα εκείνη τη στιγμή και έτρεξα προς την καλύβα για να πάρω ένα όπλο. Στην προσπάθεια μου να προλάβω τους εισβολείς πριν φτάσουν στην πόρτα έχασα το βηματισμό μου και σκόνταψα στο ίσιωμα. Ρίχνοντας μία πεταχτή ματιά στο παράθυρο από το οποίο μου έδινε οδηγίες ο Άλμπερτ, αντίκρισα ένα ανοιχτό παράθυρο χωρίς κάποιον να στέκεται πλέον εκεί. 

«Μη τρέχεις» άκουσα μία αντρική φωνή πιθανότατα να μου απευθύνεται.

Εγώ συγκεντρωμένος στον στόχο μου τον αγνόησα και σηκώθηκα. Στην καλύβα βρήκα εύκαιρο μόνο ένα πιστόλι, το οποίο και άρπαξα και έτρεξα προς την κατεύθυνση των δύο αντρών. Φτάνοντας στο πίσω μέρος του σπιτιού, είδα τους δύο άντρες να ανεβαίνουν τις σκάλες της πίσω πόρτας.

«Ακίνητοι» τους είπα και στόχευσα ενδιάμεσα τους ώστε αν ένας από τους δύο προσπαθούσε να κάνει κάτι να είχα σχεδόν έτοιμο τον στόχο μου.

Αυτοί σήκωσαν τα χέρια τους και προσπάθησαν να γυρίσουν προς το μέρος μου σιγά για να μη τους πυροβολούσα. Δεν είχα ιδέα αν το πιστόλι ήταν γεμάτο ή όχι, αλλά και μόνο η απειλή που προκαλούσε ήταν αρκετή για να τρομάξει τους δύο άντρες.

«Τι κάνετε εδώ;» τους ρώτησα προσπαθώντας να μην χάσω την προσοχή μου από πάνω τους.

«Είμαστε θύματα του πολέμου» μου απάντησε ο νεότερος. «Χάσαμε το σπίτι μας στην πολιορκία που έγινε στην πόλη δίπλα από τη δική σας και προσπάθησα να κουβαλήσω τον γέρο μου μέχρι εδώ για να μπορέσει να τον βοηθήσει κάποιος»

«Ποια πόλη δίπλα από τη δική μας;» τον ρώτησα προσπαθώντας να τον σκαλώσω. Αν δίσταζε θα μου έδινε ένα σημάδι πως έλεγε ψέματα, καθώς υπάρχει μόνο μία πόλη κοντά στη δική μας που έχει πολιορκηθεί από αντιστασιακούς πρόσφατα.

«Από την…»

«Τι γίνεται εδώ;» άνοιξε την πόρτα ο Άλμπερτ πριν προλάβει να μου απαντήσει ο συνομήλικος μου. «Κυρία μη πλησιάσετε» έκανε νεύμα με το αριστερό του χέρι στη μητέρα μου που βρισκόταν μερικά μέτρα πιο πίσω του.

«Είναι εχθροί Άλμπερτ» πρόσθεσα και κράτησα το πιστόλι μου πιο σφιχτά, διορθώνοντας τον στόχο μου.

«Όχι» είπε αυτός και πλησίασε τον έναν άντρα ο οποίος φαινόταν τραυματισμένος πρόσφατα, καθώς ο επίδεσμος που είχε γύρω από το πόδι του είχε αρκετό αίμα το οποίο είχε ακόμη έντονο χρώμα. «Είναι τραυματισμένοι» μου εξήγησε ο Άλμπερτ και βοήθησε τον νεαρό να ξαπλώσουν τον παππού του σε ένα σκαλοπάτι. «Από που είστε;» τους ρώτησε.

«Από την διπλανή πόλη. Έφερα τον γέρο μου ως εδώ μήπως υπάρχει περίπτωση να μας βοηθήσει κάποιος» επανέλαβε τη δικαιολογία του ο νέος.

«Μάλιστα» του απάντησε ο Άλμπερτ και ακούμπησε το τραυματισμένο του χέρι στο πόδι του γέρου άντρα. Υπήρχε μία μικρή παύση στη συζήτηση, μάλλον ο Άλμπερτ κάτι σκεφτόταν.

«Λέει ψέματα» του φώναξα χωρίς να χαλαρώσω τη λαβή μου. «Προσπαθεί να σε πείσει να τους βοηθήσεις Άλμπερτ»

«Τι σε κάνει να το λες αυτό;» με ρώτησε χωρίς να με κοιτάξει. Το βλέμμα του ήταν κλειδωμένο πάνω στον φρεσκοτραυματισμένο ηλικιωμένο. «Οι δύο άντρες αυτοί δεν φαίνονται για αντιστασιακοί» είπε και σηκώθηκε.

«Τι λες Άλμπερτ;» του είπα και στόχευσα τον συνομήλικο μου στο κεφάλι. «Αν… αν δεν ήταν, γιατί ήρθαν από το δάσος και όχι από την πόλη;» προσπάθησα να το παίξω ντετέκτιβ, καταλαβαίνοντας πως μπορεί να είχα κάνει λάθος εκτίμηση για το αίμα στον επίδεσμο του ηλικιωμένου άντρα, καθώς ο πιο έμπειρος στον πόλεμο Άλμπερτ θα είχε πιάσει μία τέτοια λεπτομέρεια.

«Γιατί, όπως μας είπαν, κατάγονται από την πόλη βορειοανατολικά της δικής μας» μου απάντησε αυτός με σοβαρό ύφος.

«Και η πληγή του. Πως εξηγείς…» έκανα μία τελευταία προσπάθεια να τον πείσω.

«Αρκετά» πλησίασε η μητέρα μου και έσκυψε για να κοιτάξει τον γέρο άντρα που είχε καθίσει στα σκαλοπάτια της με ένα ύπουλο χαμόγελο.

«Κυρία;» της απευθύνθηκε ο Άλμπερτ.

«Μαμά;» παραξενεύτηκα.

«Μου ήρθε μία ιδέα. Τι και αν παίρναμε αυτούς τους δύο άντρες υπό την προστασία μας» είπα και σηκώθηκε για να κοιτάξει εμένα και τον Άλμπερτ. «Θα βοήθαγε πολύ την εικόνα μας ως αριστοκρατική οικογένεια, δεν νομίζετε; “Η οικογένεια Μπερνούλι έσωσε έναν νέο άντρα και τον ηλικιωμένο παππού του”. Φαντάζομαι ήδη τα νέα να μεταδίδονται στην πόλη μας. Αν φτάσουν τα καλά λόγια στα αυτιά του στρατηγού μπορεί να αλλάξει η γνώμη του στρατηγού, ε; Τι πιστεύετε;»

«Δεν νομίζω πως δουλεύουν έτσι οι βαθμίδες του στρατού…» της είπε ο Άλμπερτ προσπαθώντας να μη φανεί ειρωνικός «…όμως θα μπορούσαμε να τους βοηθήσουμε μέχρις ότου να μπορέσουν να επανενταχθούν στην κοινωνία»

«Ναι… να τους βοηθήσουμε» του είπε η μητέρα και του έκλεισε το μάτι προσπαθώντας να είναι διακριτική αλλά το υπόλοιπο σώμα της την πρόδωσε.

«Τι;» παραξενεύτηκε ο Άλμπερτ.

«Εσύ τι πιστεύεις;» μου απευθύνθηκε η μητέρα.

Έμεινα για μια στιγμή ακίνητος. Ήταν με τα καλά της η μητέρα; Γιατί να αφήσουμε δύο άγνωστους άντρες να μπουν στο αρχοντικό μας; Ποιος ξέρει τι μπορεί να μας κάνουν, ακόμη και μεγαλοαστοί να είναι; Η αφέλεια της μητέρας μου έχει φτάσει στα ύψη.

«Τι σκέφτεσαι;» με ρώτησε η μητέρα.

Δεν της απάντησα, ήμουν συγκεντρωμένος στον στόχο μου. Ο Άλμπερτ παρατήρησε ότι η τρέχουσα κατάσταση με προβλημάτιζε και αποφάσισε να με πλησίασε.

«Τι σκέφτεσαι Μάρκο;» μου είπε και τύλιξε το αριστερό του χέρι γύρω από τον ώμο μου. «Οι άντρες είναι άοπλοι και το μόνο που κουβαλάνε είναι αυτό το σακούλι στη μέση του νεαρού άντρα που μάλλον έχει μέσα κέρματα, τα τελευταία πιθανότατα χρήματα που έχουν στο όνομα τους. Σε παρακαλώ» μου είπε και στάθηκε περιμένοντας την απάντηση μου με απλωμένο το τραυματισμένο του χέρι για αφήσω το πιστόλι στην παλάμη του.

«Καλά» του απάντησα, του έδωσα το όπλο και άρχισα να περπατάω προς το σπίτι.

«Ωραία, λοιπόν» είπε η μητέρα μου καθώς την προσπέρασα «καλωσορίσατε στο αρχοντικό μας»


\chapter{}

«Ώστε…» προσπάθησα να σπάσω τον πάγο «…εσύ είσαι ο κηπουρός ή…»

«Όχι, είμαι μέλος της οικογένειας» μου απάντησε αυτός εκνευρισμένος.

«Και ο άλλος άντρας είναι;»

«Ο Άλμπερτ, ο άντρας, είναι ο άντρας μου και η γυναίκα είναι η μητέρα μου. Οι τρεις μας μένουμε…»

«Ο άντρας σου;» παραξενεύτηκα.

«Όχι, ο μπάτλερ μου» σταύρωσε τα χέρια του ο συνομήλικος και ακούμπησε την πλάτη του στην καρέκλα του και σταμάτησε να μου μιλάει.

Όσο περιμέναμε να μας οργανώσουν ένα δωμάτιο, εγώ, ο συνταγματάρχης και ο νέος, που όπως φαίνεται δεν ήταν κηπουρός αλλά ο γιος του αρχοντικού, καθόμασταν σε ένα τεράστιο δωμάτιο το οποίο είχε στη μέση ένα μικρό στρογγυλό τραπεζάκι με τέσσερις καρέκλες, δύο καναπέδες, έναν κοντά στην πίσω πόρτα του αρχοντικού και έναν στην απέναντι μεριά, κοντά σε μία άλλη πόρτα η οποία λογικά οδηγούσε στο χολ. Οι τοίχοι ήταν λείοι, κάταστρποι και υπερβολικά ψηλοί σε σχέση με το εμβαδόν του δωματίου. Στο σημείο που ενώνονται με την οροφή υπήρχε ένα ξύλινο γύψινο, το οποίο είχε σκαλισμένα μικρά σχεδιάκια λες και ήταν κίονας κορινθιακού ρυθμού. 

Το μόνο ενοχλητικό σε εκείνο το δωμάτιο ήταν ο νέος άντρας που δεν μπορούσε να σταματήσει να κοιτάζει εμένα και τον συνταγματάρχη επίμονα. Είχε και αυτός τις υποψίες του, που λογικό μεν, αλλά μόνο και μόνο που δύο από τους τρεις κατοίκους του αρχοντικού συμφώνησαν να μας βοηθήσουν. Ή τουλάχιστον νομίζω ότι είναι τρεις. Έτσι όπως πήγε να συνεχίσει την πρόταση του ο συνομήλικος μου, το μόνο λογικό θα ήταν αυτό. Το “Σχέδιο” ήταν αρχικά φτιαγμένο για την αντιμετώπιση περισσότερων μελών του αρχοντικού. Ο στρατηγός είχε υπολογίσει βέβαια την έλλειψη πατρικής φιγούρας σε αυτό το αρχοντικό, αλλά είχε επίσης προβλέψει πως θα είχαν πάνω από έναν βοηθό, ίσως δύο ή τρεις κατά προτίμηση. Δεν νομίζω όμως αυτό να μας χαλάσει τα πλάνα μας. Ίσα ίσα, εμείς στην τελική θα ήμασταν δύο προετοιμασμένοι και αυτοί τρεις απροετοίμαστοι. 

«Όλα είναι έτοιμα» άκουσα την γυναικεία φωνή της κυρίας από την δεύτερη πόρτα πίσω μου. «Θα σας οδηγήσουμε στο δωμάτιο σας» είπε και έκανε στην άκρη οπότε σήκωσα τον συνταγματάρχη στον ώμο μου και αρχίσαμε να περπατάμε.

Διασχίζοντας τη δεύτερη πόρτα, βρεθήκαμε σε ένα τεράστιο χολ, με μία τεράστια πόρτα, μάλλον της κύριας εισόδου, να μας αντικρίζει. Στα δεξιά της πόρτα βρισκόταν μία μαρμάρινη σκάλα που οδηγούσε στον πρώτο όροφο, η οποία είχε μαύρα κάγκελα και ξύλινη κουπαστή με χρυσαφένιο χρώμα. Λίγο πιο μπροστά, στις δύο η ώρα, βρισκόταν άλλος ένας καναπές, πόσοι καναπέδες πια, και μία ακόμη πόρτα που οδηγούσε σε ένα ακόμη δωμάτιο. Στις έντεκα η ώρα, αριστερά της πόρτας  ο χώρος ήταν σχετικά κενός, με το μόνο διακοσμητικό να είναι μερικά φυτά και δύο πίνακες κρεμασμένοι στον τοίχο πάνω από αυτά.

«Από εδώ κύριοι» μας είπε η κυρία και μας κατεύθυνε προς τα αριστερά μας όπου αρκετά μακριά μας, στην άκρη του χολ, βρισκόταν ο Άλμπερτ δίπλα σε ένα μικρό άνοιγμα στον τοίχο παράλληλο με την είσοδο, χωρίς κάποια πόρτα να το καλύπτει.

Πλησιάζοντας, είδαμε κι άλλες σκάλες, οι οποίες τώρα οδηγούσαν προς ένα υπόγειο. Κοιταχτήκαμε με τον γέρο παραξενεμένοι για λίγο. Ο Άλμπερτ και η “κυρία” άρχισαν να κατεβαίνουν αυτές τις σκάλες ανενόχλητοι, οπότε τους ακολουθήσαμε. Υποθέσαμε ότι το δωμάτιο μας βρίσκεται στον κάτω όροφο, πράγμα που δεν μας χάλαγε.

Φτάνοντας στο κάτω μέρος της σκάλας αργά και σταθερά με τον συνταγματάρχη στον ώμο μου, ο Άλμπερτ άνοιξε μία σιδερένια πόρτα και η “κυρία” μας έδειξε το νέο μας δωμάτιο, το οποίο ήταν ένα σκοτεινό κιβώτιο με δύο παπλώματα στο πάτωμα, συνοδευόμενα με δύο μαξιλάρια στις κορυφές τους.

«Μην ανησυχείτε, έχετε και λάμπα» είπε η “κυρία” και τράβηξε τον διακόπτη ενός λαμπτήρα που κρεμόταν από το καλώδιο του οριακά στη μέση του δωματίου και φώτισε ο χώρος, αποκαλύπτοντας λεκέδες στους τοίχους. «Εμείς φεύγουμε για να σας αφήσουμε στην ησυχία σας τώρα. Ο Άλμπερτ θα επιστρέψει αργότερα να σας φέρει το γεύμα σας. Αν θέλετε να αφοδεύσετε, απλά πείτε το μας και… μάλλον θα σας ακούσουμε» είπε η “κυρία” και έκλεισαν την πόρτα πίσω τους, κλειδώνοντας μας μέσα στο υπόγειο.

«Μην ανησυχείτε, έχετε και λάμπα» την ειρωνεύτηκε ο συνταγματάρχης και έσπρωξε την λάμπα σαν εκκρεμές και αυτή άρχισε να ταλαντώνεται. Ύστερα, αποχώρησε από την στήριξη μου του έδινα και κουτσός πλησίασε το ένα από τα δύο παπλώματα και ξάπλωσε ανάσκελα.

Εγώ, σκεπτόμενος τα τελευταία λόγια της μεσήλικης γυναίκας, προσπάθησα από περιέργεια να ανοίξω την πόρτα και, όπως είχα προβλέψει, αυτή ήταν κλειδωμένη. Κοίταξα προς το μέρος του συνταγματάρχη και ξεφύσηξα από απόγνωση.

«Ξάπλωσε και συ» με προσκάλεσε.

«Πρέπει να ξεκινήσουμε την δεύτερη φάση του “Σχεδίου” ε;» είπα ενώ έσκυψα για να καθίσω παράλληλα στον συνταγματάρχη, χωρίς ακόμη να ξαπλώσω τελείως.

«Μην αγχώνεσαι, έχουμε χρόνο» μου απάντησε εκείνος. «Θα περάσουν αρκετές μέρες μέχρι τα σινιάλα. Μέχρι τότε,  έχεις χρόνο να παίξεις τον ρόλο σου»

«Μα αυτό με αγχώνει. Το δεύτερο μέρος του “Σχεδίου”. Ο ρόλος σου στην αποστολή μας οριακά έχει τελειώσει…» είπα και τον κοίταξα γνωρίζοντας τι θα ερχόταν για αυτών μόλις τα δύο από τα τρία σινιάλα είχαν λάβει χώρα. «Αλήθεια… γιατί συμφώνησες να έχεις αυτό το ρόλο;»

Αυτός με κοίταξε και μου χαμογέλασε. 

«Όταν έχεις ζήσει τα τελευταία δεκαπέντε χρόνια μέσα στο πεδίο της μάχης, σταματάς πλέον να φοβάσαι το αναπόφευκτο. Θέλεις απλά να τελειώσει. Όταν έχεις χάσει σχεδόν όλη σου την οικογένεια, τους φίλους σου, τους συναδέλφους σου και η κατάσταση που επικρατεί σου λέει ουσιαστικά να μην νιώσεις κανένα συναίσθημα για αυτούς και να συνεχίσεις να προχωράς, θέλεις κάποια στιγμή αυτό το μαρτύριο να τελειώσει» είπε ο συνταγματάρχης και σκούπισε το πρόσωπο του πρόωρα για να μη δω τα δάκρυα του. «Έτσι, όταν μου πρόσφερε ο στρατηγός αυτή τη θέση, δεν θα πω πως δε δίστασα λιγάκι, αλλά ήταν η μόνη διαφυγή που μπορούσα να βρω»

«Μα, αν το “Σχέδιο” πετύχει, θα ήσουν ελεύθερος να φύγεις από τον στρατό»

«Δεν έχει σημασία. Έχω ξοδέψει όλη μου τη ζωή φορώντας τη στολή μου, πριν και αφότου διχαστεί ο κόσμος σε “αριστοκράτες” και “αντιστασιακούς”. Δεν είμαι χρήσιμος έξω από το πεδίο της μάχης»

«Γιατί το λες αυτό;»

Δεν μου απάντησε. Γύρισε το σώμα του πλάγια και έφερε το μαξιλάρι του σε καλύτερη θέση ώστε να πάρει μάλλον έναν γρήγορο ύπνο όσο περιμέναμε νέα από τους πάνω. Ξάπλωσα και εγώ ανάσκελα και κοίταξα το ταβάνι, με τη λάμπα στη μέση του δωματίου να μου τυφλώνει το ένα μου μάτι.

«Οπότε, τι κάνουμε τώρα;» τον ρώτησα πριν τον πάρει ο ύπνος.

«Υπομονή»


\chapter{}

«Ώστε…» μου απευθύνθηκε ο νεότερος άντρας «…εσύ είσαι ο κηπουρός ή…»

«Όχι» τον διέκοψα πριν με προσβάλλει περισσότερο «είμαι μέλος της οικογένειας»

«Και ο άλλος άντρας είναι;» συνέχισε να με ρωτάει.

«Ο Άλμπερτ, ο άντρας, είναι ο μπάτλερ μου και η γυναίκα είναι η μητέρα μου» του απάντησα συνοψίζοντας την σχέση όλων μας στο αρχοντικό, μπας και είχε περαιτέρω ερωτήσεις. «Οι τρεις μας μένουμε…»

«Ο άντρας σου;» με ρώτησε ο συνομήλικος μου παραξενεμένος. Δεν κατάλαβε καν αυτή την απλή πρόταση;

«Όχι, ο μπάτλερ μου» του απάντησα νευριασμένος, υποθέτοντας ότι μιλούσε για τον Άλμπερτ. Για ποιον άλλο μπορούσε να αναφέρετε;

Όσο περιμέναμε την μητέρα και τον Άλμπερτ να οργανώσουν ένα δωμάτιο στο σπίτι μας, ήμουν υποχρεωμένος να προσέχω τους δύο νέους μας “συγκάτοικους”. Όσο περισσότερο παρατηρούσα τις κινήσεις τους, τόσο πιο ύποπτοι φαίνονταν. Ο γέρος δεν έκανε και πολλά πράγματα, απλά καθόταν στη θέση του σκεπτόμενος κοιτάζοντας το πάτωμα. Ο νέος από την άλλη συμπεριφερόταν λες και δεν είχε ξαναβρεθεί σε σπίτι με τέσσερις τοίχους προηγουμένως. Το βλέμμα του στρεφόταν από τους τοίχους, στους καναπέδες, στην οροφή, στις πόρτες και μετά πάλι στους τοίχους. Τι το ενδιαφέρον έβρισκε στην αρχιτεκτονική του σπιτιού; Όλα σε αυτή την πόλη είναι όμοια. Καλά… το αρχοντικό είναι ίσως λίγο πιο τραβηγμένο σε μέγεθος, αλλά σίγουρα δεν είναι το μόνο διόροφο στην πόλη μας, οπότε και στην πόλη από την οποία ισχυρίζονται ότι κατάγονται θα είχε ψηλά σπίτια. 

Ώρες ώρες ο ύποπτος άντρας γυρνούσε και κοιτούσε και εμένα ενδιάμεσα της απόλαυσης του νέου του σπιτιού, βέβαια δεν ξέρω αν η μητέρα θα τους άφηνε για πολύ στο ισόγειο. Ο τρόπος που έκλεισε το μάτι της όταν ο Άλμπερτ ανέφερε να βοηθήσουμε τους δύο άντρες θα μπορούσε να σημαίνει μόνο ένα πράγμα.

«Όλα είναι έτοιμα» είδα τη μητέρα να ανοίγει την πόρτα που οδηγούσε στο χολ. «Θα σας οδηγήσουμε στο δωμάτιο σας» είπε και έκανε λίγο στην άκρη περιμένοντας τους δύο άντρες να σηκωθούν. 

Εγώ περίμενα για λίγο στο τραπέζι που καθόμασταν. Τους είδα να διασχίζουν την πόρτα και να σταματάνε μπροστά από τη μητέρα. Εξέταζαν και οι δυο τους το νέο δωμάτιο στο οποίο βρίσκονταν και δεν την άφησαν να περάσει μπροστά τους να τους καθοδηγήσει. Αυτή έκανε μία κίνηση για να τους προσπεράσει από τα δεξιά, αλλά οι δύο άντρες έκλειναν τον στενό διάδρομο και αυτή δυσανασχέτησε. Βλέποντας τους άντρες να κάνουν άλλο ένα βήμα μπροστά, η μητέρα χώθηκε από την αριστερή πλευρά τους αυτή τη φορά και συνέχισε την πορεία της αριστερά.

«Από εδώ κύριοι» τους είπε και αυτοί την ακολούθησαν.

Προς τα αριστερά, όμως, δε βρισκόταν κάποιο δωμάτιο. Σηκώθηκα για να τους ακολουθήσω από περιέργεια να μάθω που θα φιλοξενούσαμε τους νέους μας συγκάτοικους για τις επόμενες μέρες. Διασχίζοντας τον διάδρομο τους είδα να στέκονται διστακτικά μπροστά από τις σκάλες του υπογείου, με τον Άλμπερτ και την μητέρα πιθανότατα να έχουν ήδη κατέβει στο υπόγειο. Πλησίασα τους άντρες προσεκτικά, ώστε να μη με καταλάβουν, καθώς αυτοί ξεκίνησαν να ακολουθούν τις σκάλες προς το υπόγειο. Φτάνοντας στις σκάλες, έβγαλα το κεφάλι μου ίσα ίσα για να μπορέσω να κρυφοκοιτάξω. Οι άντρες κουτσά στραβά κατέβηκαν τις σκάλες, με τον νεότερο να κουβαλάει στον ώμο του τον γέρο του και η μητέρα τους περίμενε στην κάτω μεριά της σκάλας διατηρώντας το χαμόγελο της “οικοδεσπότριας” της. Προσπάθησα να ακούσω τον διάλογο τους όταν ο Άλμπερτ τους άνοιξε την ατσάλινη πόρτα του υπογείου και η μητέρα το φως, όμως δεν μπορούσα να διακρίνω τι συζητούσαν εκείνο το μικρό χρονικό διάστημα. Μετά από τη μικρή συζήτηση τους, ο Άλμπερτ άφησε την μητέρα να περάσει έξω από το υπόγειο πριν κλειδώσει την πόρτα και οι δυο τους ξεκινήσουν να οδεύουν προς το χολ.

Προσπαθώντας να κρύψω τον εαυτό μου και να αποφύγω μία στιγμή αμηχανίας όταν οι δυο τους επέστρεφαν πίσω στο χολ, το καλύτερο που πρόλαβα να κάνω ήταν να μετακινηθώ μερικά μέτρα μακριά από την είσοδο του υπογείου και να καθίσω στο πάτωμα. Η μητέρα δεν μου έδωσε σημασία όταν με προσπέρασε. Ο Άλμπερτ, από την άλλη, με πλησίασε.

«Είσαι καλά Μάρκο;» με ρώτησε γονατίζοντας μπροστά μου για να είναι στο ίδιο επίπεδο με εμένα.

«Όχι» του απάντησα.

«Γιατί;»

«Επειδή όλα τα σημάδια είναι μπροστά σας. Οι δύο άντρες αυτοί δεν είναι φιλικοί. Είναι κατάσκοποι του αντίπαλου μας»

«Τι σε κάνει να το λες αυτό;»

«Όσο σας περιμέναμε, ο συνομήλικος μου συμπεριφερόταν περίεργα» του απάντησα στην ερώτηση που μου είχε ξανακάνει με νέα σημάδια. «Επίσης, οι άντρες δέχτηκαν να μείνουν στο υπόγειο;»

«Δεν το αρνήθηκαν»

«Ποιος αριστοκράτης θα δεχόταν να μείνει σε υπόγειο;»

Η τελευταία πρόταση μου τον έβαλε σε σκέψεις.

«Δεν συμφωνείς;» 

«Οι δύο άντρες αυτοί προφανώς ψάχνανε ένα μέρος να μείνουν. Ποιος ξέρει τι ταλαιπωρίες πέρασαν τις τελευταίες μέρες. Είναι λογικό να δεχτούν την πρώτη προσφορά που έρθει προς το μέρος τους, ακόμη και αν αυτή είναι όσο περίεργη όσο η δική μας» προσπάθησε ο Άλμπερτ να συνδέσει όλα τα στοιχεία που γνωρίζαμε για τους άντρες.

«Τους ρωτήσατε τουλάχιστον πως λέγονται;»

«Η μητέρα σου δε με άφησε. Ισχυρίζεται πως αν αυτοί οι δύο άντρες είναι κάποιοι γνωστοί στο αντίπαλο στρατόπεδο και μαθαινόταν πως τους έχουμε υπό την προστασία μας θα εκτιθέμεθα σε κίνδυνο»

«Καλά» του απάντησα και σταύρωσα τα χέρια μου.

«Σήκω τώρα και σταμάτα να κάνεις σαν μικρό παιδί» μου πρότεινε καθώς σηκωνόταν και αυτός. «Πήγαινε να αλλάξεις, θα αναλάβω εγώ τον κήπο»

«Και το χέρι σου;»

«Σιγά το τραύμα, μπορώ να το διαχειριστώ. Δεν χρειάζεται καν να το δει γιατρός» μου απάντησε αυτός και άρχισε να αποχωρεί.

«Και μετά τι θα κάνουμε;»

«Ό,τι κάναμε και τις υπόλοιπες μέρες. Υπομονή»


\chapter{}

Το επόμενο… πρωί; Δεν ξέραμε καν τι ώρα ήταν σε αυτό το δωμάτιο που βρισκόμασταν. Υπέθεσα πως ήταν πρωί καθώς ήταν η τρίτη φορά που ο μπάτλερ είχε έρθει στο δωμάτιο μας για να μας αφήσει δύο πιάτα με φαγητό. Η ήχος της κλειδαριάς μας ξυπνούσε κάθε φορά, ήμασταν εκπαιδευμένοι στρατιώτες εξάλλου.

«Ορίστε το πρωινό σας κύριοι» μας απευθύνθηκε ο Άλμπερτ και άφησε στα πόδια μας. 

«Μας λες κύριους αλλά σίγουρα δεν μας συμπεριφέρεσαι σαν να είμαστε κύριοι» τον ειρωνεύτηκε ο συνταγματάρχης καθώς σηκονώτας.

«Ορίστε;»

«Νομίζετε πως αυτό που κάνετε είναι φιλοξενία; Πως τώρα βρισκόμαστε “υπό τη φροντίδα σας”;»

Ο μπάτλερ παραξενεύτηκε με το ξαφνικό ξέσπασμα του συνταγματάρχη.

«Αν δεν μας έλεγες τον όνομα του γεύματος κάθε φορά που έμπαινες μπορεί να νομίζαμε ότι βρισκόμασταν δύο ή και τρεις μέρες σε αυτό το υπόγειο» του είπε και πήρε την πρώτη δαγκωνιά στα φρούτα που μας είχε δώσει ο μπάτλερ.

«Ακολουθώ εντολές που έχουν σκοπό την ασφάλεια της οικογένειας» του είπε ο μπάτλερ και έκλεισε την πόρτα πίσω του.

«Την τελευταία φορά που θυμάμαι να υπόθηκαν όμοιες λέξεις, αθώες ζωές χάθηκαν μπροστά στα μάτια μου. Οπότε, αν έχετε την καλοσύνη και θέλετε όντως να μας βοηθήσετε, σας παρακαλώ δώστε μας ένα αξιοπρεπές δωμάτιο. Εξάλλου, τι θα γίνει αν μαθευτεί ότι μας έχετε οριακά αιχμάλωτους» τον προειδοποίησε ο συνταγματάρχης και επιτάχυνε τον ρυθμό με τον οποίο έτρωγε.

«Θα… δω τι μπορώ να κάνω» ακούστηκε η φωνή του μπάτλερ πίσω από την πόρτα μερικά δευτερόλεπτα αφότου είχε τελειώσει τον λόγο του ο συνταγματάρχης και ακούστηκε πάλι ο ήχος της κλειδαριάς.

«Αμάν» αναφώνησε ο συνταγματάρχης μπουκωμένος.

«Δεν νομίζεις ότι ήσουν πολύ απότομος;» τον ρώτησα παραξενεμένος και γω από το ξέσπασμα του.

«Θα ήταν ύποπτο αν δεν παραπονιόμασταν» μου απάντησε με τελείως διαφορετικό ύφος από αυτό που είχε πριν δέκα δευτερόλεπτα. «Μπορεί να ήταν κάποιο τεστ που έπρεπε να περάσουμε και να μας έχωναν όντως φυλακή αν δεν τους ζητούσαμε κάποια αλλαγή»

«Ναι…» του είπα διστακτικά, καθώς ο τόνος της φωνής του μπάτλερ δεν συνιστούσε πως η τρέχουσα κατάσταση μας ήταν κάποια δοκιμασία. Ή, τουλάχιστον, έτσι μου φαινόταν. «Όμως και το αντίθετο άκρο είναι επίσης ύποπτο. Αν είσαι πολύ επιθετικός θα καταλάβουν πως προσπαθείς να το παίξεις εκλεπτυσμένος ή κάτι»

«Και πως ακριβώς θα έπρεπε να φέρομαι;»

«Λίγο πιο χαλαρά. Εξάλλου, εσύ υποτίθεται ότι είσαι ο τραυματισμένος ψυχολογικά και σωματικά από τον πόλεμο ηλικιωμένος στο “Σχέδιο”, οπότε…»

Έκανε μία παύση ο συνταγματάρχης για να με κοιτάξει περιμένοντας τις επόμενες λέξεις που θα έβγαιναν από το στόμα μου.

«…προσπάθησε να είσαι λίγο πιο ήπιος και θα δω τι μπορώ να κάνω εγώ για να μας εξασφαλίσω τη δεύτερη φάση του “Σχεδίου” μέχρι τα σινιάλα»

«Θα… δω τι μπορώ να κάνω» μου απάντησε αυτός και κοίταξε τα λίγα φρούτα που είχε αφήσει στο πιάτο του. «Γιατί δεν τρως;» με ρώτησε. «Είναι πολύ ωραία»

Η αλήθεια είναι πως από μικρός δεν μου άρεσαν και τόσο τα φρούτα. Όσο η μητέρα μου με πίεζε να τα φάω, εγώ επέμενα και καθόμουν νηστικός και κοιτούσα τους υπόλοιπους να τρώνε τις δικές τους μερίδες για αρκετή ώρα. Στην τελική, βέβαια, υπέκυπτα στην πίεση της μητέρας και την έμμεση πίεση των υπολοίπων. Μία μέρα ειδικά, όταν ήμουν πολύ μικρός, αφότου πέθανε ο πατέρας μου, θυμάμαι που είχαμε μαζευτεί πολλές οικογένειες σε ένα σπίτι και ο οικοδεσπότης μας είχε ετοιμάσει δύο γαβάθες φρούτα και όλα τα παιδιά τσιμπολογούσαν με τα πιρούνια τους ενώ οι ενήλικες συζητούσαν μεταξύ τους.

Το περίεργο δεν ήταν βέβαια τα φρούτα, αλλά το γεγονός ότι κάπως ειπώθηκε μεταξύ των παιδιών που ήμασταν εκεί η κλασική ερώτηση που κάνουν οι μεγαλύτεροι μιας οικογένειας στους μικρούς όταν δεν έχουν κάτι να συζητήσουν: “Τι θέλεις να κάνεις όταν μεγαλώσεις;”, “Τι θες να είσαι όταν μεγαλώσεις;”, ή κάτι όμοιο τέλος πάντων. Δεν ξέρω ποιο παιδί άρχισε τη συζήτηση, αλλά όλοι έδιναν απαντήσεις όπως “θέλω να ζήσω με την οικογένεια μου”, “να σταματήσω τον πόλεμο”, “να ζήσω”, που ναι μεν είναι καλές απαντήσεις, αλλά παραξενεύτηκα που κανείς δεν είχε το ίδιο όνειρο με έμενα.

«Δεν πεινάς, ε;» μου είπε ο συνταγματάρχης καθώς άφηνε το πιάτο του στο πάτωμα. «Είσαι αγχωμένος;»

Δεν του απάντησα.

«Μη σε αγχώνει το δεύτερο μέρος του “Σχεδίου”. Πιστεύω όλα θα κυλήσουν φυσικά»

Πήρα μία δαγκωνιά από το κομμάτι μήλου που βρισκόταν στην άκρη του πιάτου μου. Όπως περίμενα δεν με εξέπληξε η γεύση του φρούτου και ακούμπησα το πιάτο μου κάτω μπροστά από τα πόδια μου. Οι αριστοκράτες τρώνε τα ίδια φρούτα με εμάς δηλαδή; Τώρα που το σκέφτομαι, και τα προηγούμενα δύο γεύματα που μας είχε φέρει ο μπάτλερ δεν είχαν κάτι το ξεχωριστό. Ίσα ίσα, πιστεύω πως η μητέρα μου μαγειρεύει όσο καλά όσο αυτοί.

«Λοιπόν» ακούσαμε τον μπάτλερ να ξεκλειδώνει την πόρτα. «Ωχ. Κύριε, εσείς δεν φάγατε;» με ρώτησε αυτός αντικρίζοντας το γεμάτο πιάτο μου.

«Δεν πεινούσα» του απάντησα.

«Καλώς, δεν πειράζει» είπε και σήκωσε τα δυο πιάτα μας από το πάτωμα. «Ήρθα να σας ανακοινώσω πως το αίτημα αλλαγής δωματίου σας έγινε δεκτό από την οικογένεια και σας έχουμε ετοιμάσει ένα δωμάτιο στον πρώτο όροφο του αρχοντικού για τη διαμονή σας»

«Στον πρώτο όροφο; Επίτηδες με κάνετε να ανεβαίνω σκάλες;» του απάντησε ο συνταγματάρχης. Του έκανα νεύμα με το χέρι μου να ηρεμήσει.

«Συγγνώμη, σας διαβεβαιώνω δεν ήταν…»

«Πλάκα κάνω» ξεφύσιξε και τέντωσε το χέρι του προς το μέρος μου για να τον σηκώσω.

«Χαίρομαι» ανακουφίστηκε ο μπάτλερ. «Μέχρι να ανεβείτε τις σκάλες εγώ θα περάσω από την κουζίνα για να ακουμπήσω τα πιάτα σας. Ας συναντηθούμε στις σκάλες για να σας συνοδεύσω στο δωμάτιο σας» είπε και προχώρησε μόνος του μπροστά μας.

«Το να τους ειρωνεύομαι μερικές φορές είναι καλό, γιατί δείχνει πως δεν έχω δεχτεί ακόμη την αλλαγή στον τρόπο ζωής μου» μου είπε σιγανά ο συνταγματάρχης πριν τον σηκώσω.

«Καλά, αλλά μη το παρακάνεις» του είπα και τον τράβηξα.

Φτάνοντας στην κορυφή της σκάλας του υπογείου, ποιος άλλος θα μπορούσε να μας περιμένει παρά ο συνομήλικος φίλος μου που εξακολουθούσε να έχει το ίδιο επίμονο βλέμμα στα μάτια του. Στεκόταν ακίνητος ακριβώς μπροστά μας εκνευρισμένος και δεν έβγαζε άχνα. Εγώ, ξεκινόντας στο μυαλό μου και επίσημα τη δεύτερη φάση του “Σχεδίου” του χαμογέλασα και στράφηκα προς τις σκάλες όπου στεκόταν ο μπάτλερ και μας περίμενε. Πολύ γρήγορος αυτός ο μπάτλερ.

«Ελάτε από εδώ» μας φώναξε ΄για να τον πλησιάσουμε.

«Μισό λεπτό» άκουσα τη φωνή του συνομήλικου μου από πίσω μου καθώς κουβαλούσα τον συνταγματάρχη.

Γύρισα και είδα τον νεαρό σκεπτόμενο.

«Πως… πως σε λένε;» με κοίταξε στα μάτια και με ρώτησε

«Πόλο. Και τον παππού μου Τζιοβάνι» του είπα χαμογελώντας. «Χάρηκα για τη γνωριμία…»

«Μάρκο»

«Χάρηκα Μάρκο»

Αυτός φάνηκε προβληματισμένος, λες και η γνωριμία μας του είχε προκαλέσει σκέψεις.

«Είσαι καλά;» τον ρώτησα.

«Πηγαίνετε στο δωμάτιο σας» μου απάντησε και αποχώρησε κοιτώντας το πάτωμα του αρχοντικού πάλι σκεπτόμενος. Τι τον είχε πειράξει;


\chapter{}

«Θέλετε τίποτα παραπάνω;» είπε ο Άλμπερτ αφήνοντας μας δύο από τα τέσσερα πιάτα που κρατούσε.

«Όχι Άλμπερτ μου, ευχαριστώ» του απάντησε η μητέρα και αυτός μας άφησε μόνους μας στην τραπεζαρία να φάμε το πρωινό μας. «Φάε Μάρκο»

Η νύχτα με τους νέους μας συγκάτοικους δεν ήταν και τόσο διαφορετική από τις υπόλοιπες. Αν εξαιρέσεις το αρχικό άγχος που είχα πως θα πεθάνω πριν αποκοιμηθώ, ο ύπνος μου κύλησε ομαλά, χωρίς να δω κάποιον εφιάλτη. Θα μου πεις, πώς να πέθαινα; Οι δυο τους είναι κλειδωμένοι στο υπόγειο. Μόνο και μόνο αυτό θα έπρεπε να μου έχει απαλύνει την ανησυχία.

«Φρούτα πάλι» σχολίασα καθώς έπαιρνα την πρώτη μπουκιά μου.

«Ναι, είναι υγιεινά, για αυτό τρώγε» μου απάντησε αυτή.

«Ο Άλμπερτ; Δεν θα κάτσει μαζί μας ποτέ ξανά να φάει;»

«Πήγε να ταΐσει τους συγκατοίκους μας. Είναι το ελάχιστο που μπορούμε να κάνουμε για να τους φροντίζουμε»

Εκνευρισμένος έβαλα τον αριστερό μου αγνώνα πάνω στο τραπέζι και άρχισα να κατεβάζω το γεύμα που είχα μπροστά. Ποτέ δε μου άρεσαν τα φρούτα. Ο μόνος λόγος που τα τρώω είναι για να μην ακούω τη γκρίνια της μητέρας μου… και για να μη μείνω νηστικός. Μία μέρα θυμάμαι δεν είχα προλάβει να φάω πρωινό επειδή προετοιμαζόμασταν για μια επίσκεψη σε άλλο αρχοντικό μερικές ώρες μακριά από το δικό μας, οπότε περίμενα μέχρι το δεκατιανό για να μπορέσω να τσιμπήσω κάτι, αλλά απογοητεύτηκα μόλις ο τότε οικοδεσπότης ακούμπησε ένα πιάτο φρούτα στο τραπέζι των παιδιών που καθόμουν. Το περίεργο εκείνη τη μέρα δεν ήταν μόνο αυτό, αλλά και κάτι άλλο που συνέβει όσο όλα τα παιδιά μασουλούσαμε. Κάπως, δεν θυμάμαι πως ακριβώς, ειπώθηκε αυτή η κλασική ερώτηση που σε ρωτάνε οι μεγαλύτεροι σου όταν είσαι παιδί και δεν ξέρουν τι να συζητήσουν μαζί σου: “Τι θες να γίνεις όταν μεγαλώσεις”, ‘Τι θες να κάνεις όταν μεγαλώσεις;“, όπως και να το πει κανείς οι υπόλοιποι καταλαβαίνουν το νόημα. Οι περισσότεροι απαντούσαν “θέλω να γίνω πλουσιότερος από τους γονείς μου”, “να έχω ένα μεγάλο σπίτι”, και γενικά όλοι τους σχολίαζαν υλικά αγαθά, αγνοώντας πλήρως τις υπόλοιπες πτυχές της ζωής. Είχα παραξενευτεί τελείως που κανένας δεν είχε όνειρο που να μοιάζει έστω και ελάχιστα στο δικό μου όνειρο.

«Πεινούσες φαίνεται» μου είπε η μητέρα μου βλέποντας το άδειο μου πιάτο που είχα κατεβάσει ασυνείδητα.

Δεν της απάντησα.

«Μην αγχώνεσαι, αν θες να φας παραπάνω…»

«Δεν πεινάω άλλο μητέρα» της απάντησα και σηκώθηκα από τη θέση μου.

«Καλά, άσε το πιάτο σου τότε…»

«Κυρία μπορώ να σας μιλήσω για ένα λεπτό;» εμφανίστηκε ο Άλμπερτ από την πόρτα που οδηγούσε στο χολ.

«Μπορείτε να σταματήσετε να με διακόπτετε» παραπονέθηκε η μητέρα και σηκώθηκε. «Πες μου Άλμπερτ»

«Μήπως θα μπορέσαμε να αλλάξουμε τους “συγκάτοικους” μας δωμάτιο;» είπε διστακτικά αυτός.

«Προς τι η ξαφνική αλλαγή Άλμπερτ;» τον ρώτησε η μητέρα.

«Υπήρχαν παράπονα από πλευράς τους πως δεν τους φερόμαστε όπως αρμόζει σε θύματα πολέμου» της απάντησε ο Άλμπερτ και με κοίταξε.

«Τι σου είπα χθες Άλμπερτ;»

«Πιστεύω πως αν οι δύο αυτοί άντρες ήταν στόχος των αντιστασιακών θα το γνωρίζαμε, καθώς έχουν περάσει αρκετές μέρες από την άλωση της πόλης τους. Τα νέα θα είχαν φτάσει τα αυτιά μας»

Η μητέρα εκνευρίστηκε και χτύπησε το τραπέζι πριν κάτσει μπροστά από το πιάτο της σκεπτόμενη.

«Καταλαβαίνω τον θυμό σας, όμως πρέπει να καταλάβετε πως…»

«Κάνε ό,τι νομίζεις Άλμπερτ, δεν θέλω να ασχοληθώ άλλο» είπε και συνέχισε να τρώει τα λίγα φρούτα που της είχαν απομείνει. «Αν συμβεί κάτι όμως, θα πάρεις εσύ την ευθύνη;» τον ρώτησε.

«Ναι» της απάντησε αποφασιστικά ο Άλμπερτ.

«Καλώς» του είπε η μητέρα. «Βάλ’ τους στο δωμάτιο δίπλα από του Μάρκο που είναι κενό τόσα χρόνια. Επιτέλους βρήκαμε χρήση για αυτό το δωμάτιο»

«Μάλιστα κυρία μου» είπε ο Άλμπερτ και μάζεψε το άδειο πιάτο μου από το τραπέζι.

«Μισό, μισό. Εγώ δεν συμφώνησα» τους είπα ετεροχρονισμένα.

«Μάρκο» με πλησίασε ο Άλμπερτ. «Αν ανησυχείς ακόμη για τους δύο άντρες, θέλω να σου επιβεβαιώσω πως η συμπεριφορά τους σήμερα μου έδειξε όσα χρειαζόμασταν να δούμε για να καταλάβουμε αν οι δυο τους λένε ψέματα ή όχι»

«Μα…»

«Δεν θέλω να τους ταλαιπορήσουμε περισσότερο Μάρκο» μου είπε ο Άλμπερτ και πήγε στην κουζίνα για να αφήσει το λερωμένο μου πιάτο. «Θα πάω να μαζέψω τα δικά τους πιάτα σε λίγο» μου είπε «να τους προσκαλέσω να ανέβουν στο δωμάτιο τώρα ή αργότερα; Εξάλλου το δωμάτιο δίπλα από του Μάρκο δεν χρειάζεται τακτοποίηση» φώναξε στην μητέρα από την κουζίνα.

«Κάνε ό,τι θες» του φώναξε αυτή πίσω, χωρίς να σταματήσει να τρώει.

Ακολούθησα τον Άλμπερτ όταν επέστρεψε από την κουζίνα οδεύοντας προς το υπόγειο για να ανακοινώσει τα νέα στους συγκατοίκους μας. Τι μπορεί να του είπαν καν για να τον πείσουν τόσο άσχημα; Δεν μπορώ να καταλάβω πως ο Άλμπερτ μπορεί και αγνοεί όλα τα υπόλοιπα σημάδια και εμπιστεύεται τους δύο ξένους. Μήπως κάνω λάθος; Μήπως όντως είναι με το μέρος μας και έχω παραισθήσεις; Δεν γίνεται. Δεν γίνεται να φαντάζομαι όλες αυτές τις λεπτομέρειες. Κάτι θα πρέπει να υπάρχει για να διακρίνω αν οι δυο άντρες είναι αντιστασιακοί ή όχι, απλά δεν μου έχει έρθει ακόμη… Τα ονόματα τους. Δεν μας είπαν ποτέ τα ονόματα τους. Δεν τους ρωτήσαμε βέβαια για την “προστασία” μας, αλλά το θεωρώ βλακεία. Ο Άλμπερτ εξάλλου κατέρριψε αυτό το σκεπτικό πριν λίγα λεπτά. Τα ονόματα τους είναι το τέλειο διαχωριστικό αριστοκρατών και αντιστασιακών, καθώς οι αντιστασιακοί έχουν ονόματα από τα βόρεια ενώ εμείς από τα νότια. Όχι όλοι βέβαια, υπάρχουν και εξαιρέσεις, αλλά ποιες οι πιθανότητες να είναι αντιστασιακοί και να έχουν και οι δύο όνομα από τα νότια;

Δεν ακολούθησα τον Άλμπερτ στο υπόγειο, περίμενα στην κορυφή της σκάλας. Δεν πέρασε πολλή ώρα πριν τον δω να επιστρέφει, ανεβαίνοντας βιαστικά τις σκάλες με τα πιάτα των δύο συγκατοίκων μας στο χέρι του, το ένα γεμάτο και το άλλο άδειο. Άλλο ένα στοιχείο, όχι τόσο πειστικό που θα έκρινε την υπόθεση αν βρισκόμασταν σε δικαστήριο, αλλά ύποπτο αρκετά που μπορεί να πείσει κάποιον με τα σωστά συμφραζόμενα. Πίσω του, οι δύο άντρες, με τον νεότερο να βοηθάει τον κουτσό γέρο του να ανέβει τις σκάλες αργά και σταθερά.

Φτάνοντας μετά από αρκετή ώρα στην κορυφή της σκάλας, οριακά όση ώρα συζήτησαν με τον Άλμπερτ για την αλλαγή του δωματίου τους, οι δύο άντρες με κοίταξαν και ο νεότερος μου χαμογέλασε πριν στραφεί προς τον Άλμπερτ που στεκόταν στο επόμενο σετ από σκάλες που θα έπρεπε να ανέβουν. Έπρεπε να δράσω, να τους ρωτήσω γρήγορα. Αν ήταν αντιστασιακοί και τους τοποθετούσαμε ακριβώς δίπλα από το δωμάτιο μου θα ήμουν ο πρώτος τους στόχος. 

«Πως…» αγχώθηκα λίγο «πως σε λένε;» ρώτησα τον συνομήλικο μου.

«Πόλο» μου απάντησε αυτός «Και τον παππού μου Τζιοβάνι» και με τις λέξεις αυτές όλη μου η πραγματικότητα κατέρρευσε. «Χάρηκα για την γνωριμία…»

«Μάρκο» του απάντησα αντανακλαστικά βλέποντας πως αυτός περίμενε το όνομα μου για να ολοκληρώσει την πρόταση του.

«Χάρηκα Μάρκο»

Τα ονόματά τους ήταν νότια. Και τα δύο. Οι πιθανότητες να ήταν αντιστασιακοί έτειναν στο μηδέν. Μα πως γίνεται αυτό; Ο τρόπος που έφτασαν στο αρχοντικό, το φρέσκο αίμα στο πόδι του ηλικιωμένου, η αποδοχή της στέγασης στο υπόγειο, έστω και για μια μέρα, τα μισοφαγωμένα φρούτα. Μήπως το έχω παρατραβήξει;

«Είσαι καλά;» με ρώτησε ο Πόλο.

«Πηγαίνετε στο δωμάτιο σας» του απάντησα και προσπάθησα να φύγω από εκεί όσο πιο γρήγορα μπορούσα.

Μήπως το έχω παρατραβήξει;


\chapter{}

Έχουν περάσει δύο μέρες από τη γνωριμία μου με τον Μάρκο και ακόμη φαίνεται προβληματισμένος. Έχει σταματήσει βέβαια να μας κοιτάζει κατάματα τόσο επίμονα όσο προηγουμένως, το οποίο είναι κάτι θετικό. Δυσκολεύομαι να βρω κάποιο περιθώριο όμως για να συνεχίσω τη δεύτερη φάση του “Σχεδίου”, καθώς ο συνομήλικος μου δε φαίνεται και πολύ κοινωνικός. Αν δεν αρχίσω σύντομα μπορεί να μας διώξουν από το αρχοντικό και μετά όλο το πλάνο της αντίστασης πάει στα σκουπίδια.

Ο Τζιοβάνι από την άλλη καθόταν χαλαρός δίπλα μου στην τραπεζαρία γνωρίζοντας πως δεν χρειαζόταν να βοηθήσει απαραίτητα στη δεύτερη φάση. Περίμενε τα επόμενα σινιάλα ξέγνοιαστος. Άραγε του περνούσε ποτέ από το μυαλό του η τελευταία φάση του “Σχεδίου”;

Ο Μάρκο καθόταν απέναντι μου και έπαιζε βαριεστημένος με μία οδοντογλυφίδα. Λύγιζε το ξύλο μέχρι αυτό να έπαιρνε μορφή που να μοιάζει με μέρος μιας έλλειψης και μετά το άφηνε να εκτιναχθεί μερικά εκατοστά μπροστά του. Επανάλαβε αυτή την κίνηση αρκετές φορές μέχρις ότου η μητέρα του σηκώθηκε από το τραπέζι. Προσπαθώντας βιαστικά να πετάξει την οδοντογλυφίδα άλλη μία φορά πριν την ακολουθήσει έξω από την τραπεζαρία, δεν υπολόγισε σωστά τη γωνία που την έστησε και αυτή πετάχτηκε προς το μέρος μου, πέφτοντας πάνω στα ρούχα που μας είχαν δώσει οι αριστοκράτες.

«Ωπ, δεν πειράζει» του είπα και έπιασα την οδοντογλυφίδα από το πάτωμα για να του την επιστρέψω. «Ορίστε»

Αυτός μισο έκλεισε τα μάτια του και πλησίασε τον Άλμπερτ, που είχε μόλις μπει στο δωμάτιο για να μαζέψει τα πιάτα μας. Ο νεαρός του ψιθύρισε κάτι και τον ανάγκασε να βγει μαζί του από το δωμάτιο πριν προλάβει να πάρει τα δικά μας πιάτα. Προσπάθησα να πλησιάσω για να κρυφακούσω, αλλά δεν μπορούσα να φτάσω αρκετά κοντά ώστε να μπορέσω να καταλάβω τον διάλογο τους και να μη γίνω αντιληπτός.

«Έλα δω» μου είπε ο Τζιοβάνι και μου έκανε νεύμα για να τον πλησιάσω.

«Είναι πιο δύσκολο από ότι περίμενα» του είπα και κάθησα πάλι στη θέση από την οποία είχα σηκωθεί.

«Μην σε ανησυχεί, όλα θα έρθουν φυσικά. Είστε συνομήλικοι εξάλλου. Αυτός είναι και ο λόγος που είσαι στο “Σχέδιο”» προσπάθησε να με καθησυχάσει αυτός.

«Όταν λες φυσικά…» δεν πρόλαβα να τον ρωτήσω γιατί μπήκε ο Μάρκο πάλι στην τραπεζαρία και κάθισε στην ίδια θέση με πριν.

«Δεν είσαι πολύ ομιλητικός, ε; Ή, τουλάχιστον, με εμάς» μου είπε και άρχισε να συνέχισε το παιχνίδι του με την οδοντογλυφίδα του.

Η αλήθεια είναι πως τρεις μέρες στο αρχοντικό οι μόνες κουβέντες που είχαμε ανταλλάξει με τους οικοδεσπότες μας ήταν η πρώτη μας επαφή στον κήπο, οι δύο ή τρεις φορές που ο Τζιοβάνι μίλησε στον Άλμπερτ για την αλλαγή του δωματίου μας, και, το πιο σημαντικό, η γνωριμία μας. Από τότε, οι μόνες λέξεις μας που έχουν απευθυνθεί προς αυτούς είναι ευχαριστίες για τα γεύματα που μας προσφέρουν τρεις φορές τη μέρα. Και τα ευχαριστώ μας κυρίως απευθύνονται στον μπάτλερ, οπότε λογικό ο Μάρκο να νιώθει πως δεν του μιλάω. Δηλαδή, προσπάθησα να του μιλήσω τώρα που του έδωσα την οδοντογλυφίδα, όχι;

«Απέδειξες μόλις το επιχείρημα μου» μου είπε και έσπασε την οδοντογλυφίδα του. 

Ο μπάτλερ μπήκε στο δωμάτιο και κατευθύνθηκε προς το μέρος μας.

«Πάντα έτσι περνάτε τον χρόνο σας;» τον ρώτησα και έκανα λίγο πίσω την καρέκλα μου ώστε να μπορέσει ο μπάτλερ να μαζέψει τα πιάτα μας.

«Τι εννοείς;» με ρώτησε και πέταξε την οδοντογλυφίδα προς το μέρος μου, όμως αυτή τη φορά είχε σημαδέψει το πιάτο που σήκωνε εκείνη τη στιγμή ο μπάτλερ του, το οποίο και πέτυχε με εκπληκτική ακρίβεια και πανηγύρισε με κλειστή τη γροθιά του. «Γιατί δεν έχει γίνει άθλημα αυτό;»

«Επειδή υπάρχει πόλεμος Μάρκο» του απάντησε ο μπάτλερ ξερά.

«Καλά» είπε αυτός και έκατσε στην καρέκλα του ξενερωμένος. «Τι εννοούσες εσύ;» μου απευθύνθηκε ο Μάρκο πάλι.

«Εννοώ… δεν κάνετε κάτι άλλο εκτός από το…» σταμάτησα για λίγο να σκεφτώ κάτι να πω το οποίο δεν θα τους πρόσβαλε.

«Να κοιτάμε τους τοίχους και το πάτωμα;» μου απάντησε αυτός.

«Μπορείς να το πεις και έτσι»

«Γιατί εσείς τι κάνατε όταν δεν υπήρχε κάποια σοβαρή δουλειά πίσω στην πόλη σας;»

«Εγώ…» αγχώθηκα, έπρεπε να βρω γρήγορα μία απάντηση «Εγώ συνήθως γράφω» είπα το πρώτο πράγμα που ήρθε στο μυαλό μου.

«Γράφεις;» παραξενεύτηκε.

«Ναι γράφω»

«Τι δηλαδή; Μυθιστορήματα και βιβλία και τέτοια;»

«Ναι ναι» συμφώνησα μαζί του παρόλο που δεν είχα γράψει ποτέ πάνω από πέντε προτάσεις συνεχόμενα.

Το βλέμμα του έδειχνε πως είχε εντυπωσιαστεί λίγο από το ψεύτικο χόμπι μου. Ο νεαρός σηκώθηκε και κατευθύνθηκε προς το χολ.

«Έλα» μου είπε και στάθηκε στην πόρτα.

«Που;»

«Στο δωμάτιο μου» μου απάντησε.

Γύρισα και κοίταξα τον Τζιοβάνι πριν σηκωθώ και αυτός μου έγνεψε θετικά με το κεφάλι του.

«Μην ανησυχείς, δεν μπορώ να πάω και πουθενά» μου είπε αυτός παιχνιδιάρικα. «Πήγαινε με τον νεαρό να περάσεις τον χρόνο σου» μου χαμογέλασε και με βοήθησε να σηκωθώ από την καρέκλα μου. Του χαμογέλασα πίσω. 

Ακολούθησα το Μάρκο στις σκάλες και στη συνέχεια στο δωμάτιο του που βρισκόταν δίπλα από το δικό μας. Ο χωρος του ήταν πολύ διαφορετικός. Αρχικά, το δωμάτιο του ήταν αρκετά πιο πλατύ από το δικό μας, είχε δύο παραπάνω παράθυρα, το ένα εκ των οποίων ήταν κάθετο στα υπόλοιπα επειδή το δωμάτιο βρισκόταν στη γωνία του αρχοντικού, ενώ τα άλλα δύο παράθυρα ήταν σας το δικό μας, πολύ μεγαλύτερα από το τρίτο και κοιτούσαν προς την πόλη. Ανάμεσα από τα δύο αυτά παράθυρα βρισκόταν ένα γραφείο με αρκετά φύλλα χαρτί, μερικές πένες και δύο τετράδια ακουμπισμένα πάνω. Το κρεβάτι του ήταν όμοιο με το δικό μας σε μέγεθος, παρόλο που αυτός ήταν ένας και εμείς ήμασταν δύο. Η ντουλάπα του ήταν τεράστια, πόσα ρούχα να είχε μέσα; Το πάτωμα του δωματίου του δεν ήταν και τόσο τακτοποιημένο, με στρατιωτάκια, ξύλινα αυτοκίνητα και σπιτάκια σκορπισμένα από ζωγραφισμένες σελίδες χαρτί. Χαρτί που μάλλον αντιπροσώπευε χάρτη. Υπέθεσα δηλαδή ότι ήταν χάρτης, καθώς από ένα σημείο του μου φαινόταν αρκετά γνώριμο, καθώς το είχα ξαναδεί πρόσφατα.

«Θα έρθεις;» με ρώτησε ο Μάρκο που είχε προχωρήσει πέρα από τα στρατιωτάκια καθώς εγώ παρατηρούσα το δωμάτιο του από την πόρτα.

«Ναι, αλλά… τι κάνουμε εδώ;» τον ρώτησα.

«Έλα στο γραφείο μου» μου είπε και κάθισε στο κρεβάτι του.

Πέρασα, λοιπόν, πάνω από τα στρατιωτάκια του και βρέθηκα μπροστά του να στέκομαι περιμένοντας την επόμενη του εντολή καθοδήγησης.

«Κάθισε» μου είπε και τράβηξε την καρέκλα του. «Πάρε μια πένα και γράψε να περάσει η ώρα» μου είπε δείχνοντας το γραφείο του.

«Ωραία…» είπα και ξεφύσηξα αγχωμένος με την κατάσταση που είχα φέρει τον εαυτό μου καθώς τέντωνα το χέρι μου για να πιάσω ένα τετράδιο και μία από τις πένες του συνομήλικου μου.

«Όχι αυτό το τετράδιο…» μου είπε και το πήρε από τα χέρια μου. «Αυτό… θα το πάρω εγώ» μου είπε και άνοιξε τις πρώτες σελίδες του.

«Τι είναι αυτό;»

«Στρατηγικές» μου απάντησε αυτός.

«Στρατηγικές;»

«Για τον πόλεμο. Αυτό κάνω στον ελεύθερο μου χρόνο» μου απάντησε και έκλεισε το τετράδιο. «Αυτό το τετράδιο είναι γεμάτο, για αυτό το πήρα. Το άλλο τετράδιο πιο δεξιά είναι μισοάδειο. Πάρε αυτό και γράψε αν θες» με παρότρεινε.

«Εντάξει» του απάντησα και πήρα το δεύτερο τετράδιο.

Το τετράδιο αυτό επίσης είχε στις πρώτες σελίδες σχέδια όπως το προηγούμενο. Τα σχέδια ήταν κυρίως κύκλοι με βέλη να τους συνδέουν και ονόματα πόλεων και στρατοπέδων στο παρασκήνιο. Πέρασα αρκετές σελίδες και εντυπωσιάστηκα στον όγκο “στρατηγικών” που είχε ζωγραφίσει ο Μάρκο στα τετράδια του. Κάθε σελίδα ήταν περιείχε και κάτι καινούργιο. Μερικές φορές, αντί για ονόματα πόλεων, είχε υποθετικά σενάρια με διάφορα πλεονεκτήματα και μειονεκτήματα για τους δύο στρατούς γραμμένα στην κάτω άκρη της σελίδας με λεπτομέρεια.

Ο Μάρκο είχε καθίσει στο πάτωμα πάνω στον αυτοσχέδιο χάρτη του και μετακινούσε στρατιωτάκια σκεπτόμενος. Φτάνοντας σε μία κενή σελίδα επιτέλους, σκέφτηκα πως έπρεπε να κάνω και γω κάτι, γιατί αλλιώς θα του κινήσω υποψίες και θα καταλάβει πως είπα ψέματα. Σκεπτόμενος τι να γράψω, είπα στον εαυτό μου ότι θα προσπαθούσα να γράψω “κάτι” που θα περιέγραφε κάποιο συμβάν ή νέο που έγινε πρόσφατα. Έτσι, δεν θα μπορέσω να ξεμείνω από ιδέες, αφού θα έγραφα μία ιστορία με σημείο αναφοράς κάτι υπαρκτό. Δεν μπορούσα όμως και να γράψω λεπτομέρειες για το γεγονός που θα περιέγραφα, καθώς πάλι θα μου χαλούσε την κάλυψη μου, οπότε έπρεπε να είμαι όσο πιο γενικός γίνεται. Από την άλλη, δεν μπορούσα να γράψω για κάτι το οποίο ήξερε ο Μάρκο ότι μου συνέβει, καθώς πάλι θα έγειρε υποψίες.

«Το βρήκα» είπα φωναχτά και ξεκίνησα να γράφω.

Είδα με την άκρη του ματιού μου τον Μάρκο να με κοιτάζει όταν με άκουσε ξαφνικά να μιλάω μετά από μερικά λεπτά σιωπής, όμως δεν με ενόχλησε. Είχα βρει τι μπορώ να γράψω ώστε να περάσει η ώρα και να μη γίνω ξαφνικά ύποπτος, οπότε ξεκίνησα. Είχα πάρει φόρα είναι η αλήθεια, δεν μπορούσα να σηκώσω την πένα μου από το χαρτί. Όσο παραπάνω έγραφα, τόσο περισσότερα πράγματα σκεφτόμουν που ήθελα να αλλάξω ή να προσθέσω και αυτό μου δημιουργούσε μία επιθυμία να συνεχίσω.

Μετά από αρκετή ώρα, ο Μάρκο σηκώθηκε και με πλησίασε για να πάρει μάτι. 

«Όχι, συνέχισε» μου είπε όταν σταμάτησα καθώς στάθηκε δίπλα μου.

Έβλεπα τα μάτια του να πηγαίνουν από τη μία άκρη της σελίδας στην άλλη πολύ πιο γρήγορα από το ρυθμό που μπορούσα να διαβάσω εγώ, κάνοντας μικρές παύσεις ενδιάμεσα.

«Γιατί γράφεις λες και απαριθμείς σε λίστα;» με ρώτησε αυτός.

«Επειδή… είναι σαν ένα προσχέδιο» του απάντησα ότι μου ήρθε εκείνη τη στιγμή για να φανώ λες και ήξερα τι έκανα.

«Ενδιαφέρον» μου απάντησε αυτός και συνέχισε να διαβάζει.

«Πρέπει να βελτιώσεις τον γραφικό σου χαρακτήρα» μου είπε και συνέχισε να διαβάζει. «Αυτό ήταν;» μου είπε μετά από λίγη ώρα, όταν είχε περάσει όλα όσα είχα γράψει.

«Ναι;» παραξενεύτηκα με τη βιασύνη του. «…για την ώρα»

«Ενδιαφέρον» μου απάντησε αυτός και αποχώρησε. «Συνέχισε με την ιστορία σου. Θέλω να τη διαβάσω όταν τελειώσεις» με ενθάρρυνε πριν καθίσει και πάλι με τα στρατιωτάκια του.


\chapter{}

Άλλη μία βαρετή μέρα στο σπίτι, με τους δυο μας συγκάτοικους να τρώνε στην τραπεζαρία μαζί μας άλλο ένα γεύμα. Η “φιλοξενία” μας έχει φτάσει κυριολεκτικά σε άλλο επίπεδο. Τους έχουμε αφήσει ελεύθερους να τριγυρνάνε στο σπίτι μας και μάλιστα δεν έχουμε θέσει ούτε ένα όριο για το που μπορούν να βρίσκονται κάθε χρονική στιγμή, το οποίο το θεωρώ ακραίο. Μέχρι στιγμής δεν μας έχουν κάνει τίποτα όσον αφορά θέματα κλοπής ή καταπάτησης στα δωμάτια μας, αλλά αυτό δεν αποδεικνύει τίποτα. Αν μου έχει μείνει κάτι από τα μαθηματικά που έχω διδαχθεί είναι αυτό.

Όσο περίμενα τη μητέρα να σηκωθεί ώστε ο Άλμπερτ να εισέλθει στην τραπεζαρία και να μαζέψει τα πιάτα μας, έπαιζα με μία οδοντογλυφίδα που είχα αρπάξει από την κουζίνα πριν έρθω και καθίσω για το γεύμα. Μου άρεσε ο τρόπος που λύγιζε το ξύλο όταν το πίεζα από πάνω κάθετα στο επίπεδο του τραπεζιού και στη συνέχεια ο τρόπος που πεταγόταν φαινομενικά απρόβλεπτα σαν ελατήριο. Εγώ προσπαθούσα να προβλέψω από μέσα μου σε ποια κατεύθυνση θα εκτινασσόταν η οδοντογλυφίδα. Μπορώ να πω πως αυτό το παιχνίδι με κράτησε απασχολημένο για αρκετή ώρα.

Όταν σηκώθηκε η μητέρα, προσπάθησα βιαστικά να κάνω άλλη μία επανάληψη του παιχνιδιού αυτού που είχα εφεύρει, αλλά από τη βιασύνη μου λύγισα την οδοντογλυφίδα προς τη λάθος μεριά και εκτινάχθηκε στην άλλη άκρη του τραπεζιού όπου καθόταν ο συνομήλικος μου με τον παππού του.

«Ωπ, δεν πειράζει» μου είπε όταν η οδοντογλυφίδα μου έσκασε στα ρούχα του και μου την επέστρεψε. «Ορίστε»

Δεν πειράζει; Εγώ θα είχα θυμώσει αν κάτι τόσο βρώμικο είχε πέσει πάνω στα δικά μου ρούχα και θα ζητούσα από τον Άλμπερτ να μου φέρει αμέσως αλλαξιά. Τι εννοεί δεν πειράζει; Είδα τον Άλμπερτ να εισέρχεται εκείνη τη στιγμή στο δωμάτιο αφού η μητέρα μου είχε σηκωθεί και ήξερα πως έπρεπε να του μιλήσω. 

«Μπορώ να σου πω κάτι;» του ψιθύρισα καθώς αυτός σήκωνε τα πιάτα μας από το τραπέζι.

«Τι συμβαίνει;» μου απάντησε αυτός και στράφηκε προς τους δύο συγκατοίκους μας για να μπορέσει να μαζέψει και τα δικά τους πιάτα.

«Είναι σημαντικό» του είπα και τον τράβηξα προς το χολ.

«Τι συμβαίνει Μάρκο;» νευρίασε ο Άλμπερτ με τη ξαφνική συμπεριφορά μου.

«Είναι για τους συγκατοίκους μας»

«Όχι, δεν θα ακούσω άλλη λέξη» μου είπε και τίναξε το χέρι με το οποίο τον κρατούσα.

«Δεν νομίζεις ότι είναι λίγο περίεργος ο τρόπος που συμπεριφέρονται ακόμη και αφότου μάθαμε τα ονόματα τους; Μόλις τώρα έπεσε η οδοντογλυφίδα μου πάνω στον νέο άντρα και αυτός αντί να νευριάσει που του λέρωσα τα ρούχα, δεν είχε καμία αντίδραση»

«Αρκετά» μου είπε αυτός και με κοίταξε στα μάτια. «Πιστεύεις πως αν ήταν όντως εχθροί, δεν θα μας είχαν ήδη επιτεθεί στον ύπνο μας ή κάτι;»

«Ναι αλλά…»

«Τους έλεγξα, δεν κουβαλούσαν τίποτα όταν μπήκαν στο αρχοντικό μας, εκτός από ένα σακούλι με νομίσματα και πεταμένα φάρμακα, μάλλον για τον μεγαλύτερο άντρα» προσπάθησε να μου πει με ήρεμο ύφος ο Άλμπερτ. «Οπότε, σε παρακαλώ, σταμάτα να κάνεις αρνητικές σκέψεις και κάθισε αν θες στην τραπεζαρία ή κάνε ό,τι θες τέλως πάντων. Θα επιστρέψω στην τραπεζαρία σε λίγο» 

«Καλά…» του είπα και επέστρεψα στην τραπεζαρία, παρόλο που είχε σηκωθεί η μητέρα μου, για να συνεχίσω το παιχνίδι μου.

Εκεί, είδα προφανώς τους δυο φιλοξενούμενους μας να κάθονται περιμένοντας μάλλον τον Άλμπερτ να τους μαζέψει και τα δικά τους πιάτα. Οι δυο τους κάθονταν κοντά ο ένας στον άλλον σαν να μιλούσαν μεταξύ τους πριν μπω εγώ και τους διακόψω.

«Δεν είσαι πολύ ομιλητικός, ε;» είπα στο συνομήλικο μου όταν κάθησα στην καρέκλα μου. «Ή, τουλάχιστον, με εμάς»

Αυτός δεν μου απάντησε. Κάθισε στη θέση του παρατηρώντας εμένα που λύγιζα για πολλοστή φορά την οδοντογλυφίδα μου και εκτινάσσοντας την στην αντίθετη κατεύθυνση από την οποία είχα προβλέψει. Που να πάρει η ευχή.

«Απέδειξες μόλις το επιχείρημα μου» του είπα και έσπασα την οδοντογλυφίδα από τα νεύρα μου.

Ο Άλμπερτ εισήλθε στο δωμάτιο και κατευθύνθηκε προς το μέρος τους.

«Πάντα έτσι περνάτε τον χρόνο σας;» με ρώτησε αυτός ξερά, χωρίς να θέσει άλλα συμφραζόμενα.

«Τι εννοείς;» τον ρώτησα και πέταξα την οδοντογλυφίδα προς το μέρος του Άλμπερτ ευελπιστώντας ότι θα έπεφτε μέσα σε ένα από τα δύο πιάτα, πανηγυρίζοντας όταν έγινε πραγματικότητα. «Γιατί δεν έχει γίνει άθλημα αυτό;» αναρωτήθηκα φωναχτά νιώθοντας την αδρεναλίνη από την επιτυχής βολή μου στο μετακινούμενο πιάτο.

«Επειδή υπάρχει πόλεμος Μάρκο» μου απάντησε ο Άλμπερτ νευριασμένος, λες και δεν μπορούσε να συνυπάρχουν.

«Καλά» είπα για να μην εντείνω τον θυμό του και κάθισα πάλι στην καρέκλα μου. «Τι εννοούσες εσύ;» απευθύνθηκα στον συνομήλικο μου πάλι.

«Εννοώ…» σκέφτηκε λίγο πριν απαντήσει «δεν κάνετε κάτι άλλο εκτός από το…»

«Να κοιτάμε τους τοίχους και το πάτωμα;» συμπλήρωσα την πρόταση του.

«Μπορείς να το πεις και έτσι» ενέκρινε αυτός.

«Γιατί εσείς τι κάνατε όταν δεν υπήρχε κάποια σοβαρή δουλειά πίσω στην πόλη σας;»

«Εγώ…» έκανε πάλι μία μικρή παύση. «Εγώ συνήθως γράφω» μου απάντησε.

«Γράφεις;»

«Ναι γράφω»

«Τι δηλαδή; Μυθιστορήματα και βιβλία και τέτοια;» τον ρώτησα για να διευκρινίσει.

«Ναι ναι» μου απάντησε αυτός κατευθείαν, αρκετά ύποπτα έχω να πω.

Τον κοίταξα για λίγο δείχνοντας εντυπωσιασμένος που κάποιος άλλος στην ηλικία μου έκανε κάτι έστω και κάπως “δημιουργικό” όσο ο χρόνος μας περνούσε, χωρίς καμία ελπίδα να τελειώσει ο πόλεμος. Ταυτόχρονα, σκέφτηκα πως θα μπορούσα να ελέγξω κρυφά αν έλεγε ψέματα. Σηκώθηκα, λοιπόν, και κατευθύνθηκα προς τις σκάλες για τον πρώτο όροφο.

«Έλα» του είπα περιμένοντας τον στην πόρτα.

«Που;» αναρωτήθηκε αυτός.

«Στο δωμάτιο μου» είπα γνωρίζοντας πως είναι λίγο ριψοκίνδυνο αν βρισκόμουν στο χειρότερο σενάριο, αλλά άξιζε τον κόπο.

Ο συνομήλικος μου δίστασε για λίγο και κοίταξε τον παππού του, λες και αυτός θα του έδινε κάποια συμβουλή ή κάτι. Αυτός έγνεψε θετικά το κεφάλι του.

«Μην ανησυχείς» του είπε «δεν μπορώ να πάω και πουθενά. Πήγαινε με τον νεαρό να περάσεις τον χρόνο σου» του χαμογέλασε και τον βοήθησε να σηκωθεί από την καρέκλα του.

Στρίβοντας δεξιά στην κορυφή της σκάλας και φτάνοντας στο υπνοδωμάτιο μου, αυτός φάνηκε να είχε εκπλαγεί από το μέγεθος του δωματίου. Έχω να πω πως μου έχει δοθεί ίσως το δεύτερο μεγαλύτερο υπνοδωμάτιο στο σπίτι, μετά των γονιών μου προφανώς. Ποτέ δεν κατάλαβα γιατί το δωμάτιο της αδελφής μου, το οποίο ήταν τώρα υπνοδωμάτιο των φιλοξενούμενων μας, ήταν το υποδιπλάσιο σε μέγεθος, ίσως οι κατασκευαστές δεν υπολόγισαν σωστά τις μετρήσεις τους ή κάτι. Προχώρησα περνώντας πάνω από από τα στρατιωτάκια και τους χάρτες μου και στάθηκα ανάμεσα στο γραφείο και το κρεβάτι μου.

«Κάθισε» τον παρακάλεσα να πλησιάσει και τράβηξα την καρέκλα του γραφείο μου ελάχιστα πιο πίσω ώστε να καταλάβει που ακριβώς τον προσκαλούσα.

«Ωραία» απάντησε αυτός λίγο αγχωμένος και πλησίασε.

Στο γραφείο μου είχα ακουμπισμένα τα δύο τετράδια με τις στρατηγικές μου. Αυτός, μάλλον από ευγένεια, πήρε το τετράδιο που βρισκόταν πιο μακριά από την καρέκλα ελπίζοντας μάλλον πως θα είναι αυτό που θα είχα χρησιμοποιήσει λιγότερο.

«Όχι αυτό το τετράδιο» το πήρα από τα χέρια του. «Αυτό… θα το πάρω εγώ» του είπα και άρχισα να ξεφυλλίζω τις σελίδες για να δει πως ήταν γεμάτες.

«Τι είναι αυτό;» παραξενεύτηκε με τα σχέδια μου.

«Στρατηγικές» του απάντησα.

«Στρατηγικές;»

«Για τον πόλεμο. Αυτό κάνω στον ελεύθερο μου χρόνο» του απάντησα πάλι και έκλεισα το τετράδιο. «Αυτό το τετράδιο είναι γεμάτο, για αυτό το πήρα. Το άλλο τετράδιο πιο δεξιά είναι μισογεμάτο. Πάρε αυτό αν θες» του πρότεινα.

«Εντάξει» μου απάντησε και πήρε το δεύτερο τετράδιο από το γραφείο μου.

Τον άφησα λίγο να επεξεργαστεί το γραφείο μου μέχρι να αρχίσει να γράφει. Δεν είχα κάτι να κρύψω, σε αντίθεση πιθανότατα με αυτόν, οπότε του έδωσα την ελευθερία να κοιτάξει λίγο το τετράδιο μου που είχε στα χέρια του μέχρι να σκεφτεί κάποια ιδέα να γράψει, αν δεν είχε ήδη κάποια ιστορία. Αυτός γυρνούσε τις σελίδες σκεπτόμενος, προσπαθώντας να αποκωδικοποιήσει τα σχέδια μου, μάλλον αποτυχημένα. Τα παράτησε αρκετά γρήγορα και άφησε το τετράδιο στο γραφείο και έσκηψε από πάνω του, ακουμπώντας τον δεξί του αγκώνα δίπλα στο τετράδιο και κουνώντας μία πένα δίπλα από το κεφάλι του, με το βλέμμα του στραμμένο στην κενή σελίδα που είχε μπροστά του.

«Το βρήκα» είπε και ξεκίνησε να γράφει.

Δεν γύρισα να τον κοιτάξω, ήμουν αφοσιωμένος στη δική μου δραστηριότητα. Η περιέργεια όμως με έτρωγε, ήθελα να δω τι ακριβώς είχε γράψει. Συγκράτησα τον εαυτό μου για αρκετή ώρα, ρίχνοντας κλεφτές ματιές πίσω μου για να δω πόσες σελίδες είχε γεμίσει ο συνομήλικος μου και έχω να πω πως έγραφε πολύ σιγά. Τα χέρια του έτρεμαν, η ανάσα του δεν ήταν σταθερή, το σώμα του είχε μία αφύσικη υπερβολική καμπούρα πάνω το χαρτί. Πολλά παράξενα πράγματα στη συμπεριφορά του και πάλι, αλλά τον άφησα να συνεχίσει.

Όταν τον είδα να αλλάζει για πρώτη φορά σελίδα αποφάσισα να τον πλησιάσω. Αυτός σήκωσε την πένα του και μετακίνησε την καρέκλα του γραφείο λίγο πιο δεξιά βλέποντας με να στέκομαι δίπλα του.

«Όχι, συνέχισε» του είπα γυρνώντας πίσω τη σελίδα για να διαβάσω τι είχε γράψει.

Τα γράμματα του ήταν άθλια, λες και είχε να γράψει σε χαρτί πάρα πολύ καιρό, ή λες και δεν είχε μάθει να γράφει σωστά. Αγνοώντας αυτή την ενόχληση, προσπάθησα να διαβάσω την ιστορία του συνομήλικου μου, αλλά αυτή είχε δομή λίστας, λες και απαριθμούσε γεγονότα.

«Γιατί γράφεις λες και απαριθμείς σε λίστα;» τον ρώτησα.

«Επειδή… είναι σαν ένα προσχέδιο» μου απάντησε και κοίταξε το πάτωμα ντροπιασμένος.

«Ενδιαφέρον» του απάντησα για να τον “ενθαρρύνω” και συνέχισα να διαβάζω.

“Στρατόπεδο μετώπου”

“Ο στρατηγός βρίσκεται στη σκηνή του και οργανώνει τη γραφειοκρατία εκείνου του πρωινού”

Μπλα, μπλα, μπλα.

“Ο στρατηγός νιώθει ένα… εκείνο… εμένα;”

«Πρέπει να βελτιώσεις τον γραφικό σου χαρακτήρα» έκανα μια παύση και του είπα.

“Ήταν η κόρη του” τι;

“Είχε εισβάλει στο στρατόπεδο κρυφά γιατί ήθελε να βρίσκεται κοντά στον πατέρα της”

Σταμάτησα και τον κοίταξα για λίγο, αλλά αυτός ακόμη είχε το βλέμμα του στραμμένο προς το πάτωμα και περίμενε να τελειώσω την ανάγνωση.

“Συζητάει με την κόρη του” μπλα, μπλα, μπλα.

“Ξαφνικά… οι εχθροί τους τους αιφνιδιάζουν εκτός ώρας μάχης…”

Γύρισα την σελίδα για να δω μήπως είχε προλάβει να γράψει τίποτα από πίσω, αλλά δεν υπήρχε σημάδι μελανιού ακόμη.

«Αυτό ήταν;» τον ρώτησα.

«Ναι… για την ώρα»

«Ενδιαφέρον» του απάντησα και αποχώρησα από κοντά του. «Συνέχισε με την ιστορία σου. Θέλω να τη διαβάσω όταν τελειώσεις» του είπα πριν καθίσω πάλι μαζί με τα στρατιωτάκια μου.

Μία ιστορία για έναν στρατηγό με την κόρη του στο στρατόπεδο του μετώπου. Λες να με ειρωνεύεται; Ή μήπως…


\chapter{}

Οι επόμενες μέρες πέρασαν κάπως όμοια. Εγώ έγραφα, ο Μάρκο μου ζήταγε που και που σελίδες από το τετράδιο μου για να μπορέσει να σχεδιάσει τα δικά του. Γενικά επικρατούσε ησυχία στο δωμάτιο διακοπτόμενη από συχνές επισκέψεις στο γραφείο από τον Μάρκο για να διαβάσει τι νέο είχα προλάβει να γράψω και αντιδρώντας συνήθως μόνο με την ίδια λέξη: «ενδιαφέρον». Η μόνη φορά που δεν είχε αυτή την αντίδραση ήταν όταν στο “προσχέδιο” της ιστορίας που αφορούσε την αδελφή μου έφτασα στον θάνατο της. Τότε ήταν που με κοίταξε πλάγια και άφησε την σελίδα που κρατούσε ελεύθερη. 

«Όλα καλά;» τον ρώτησα παραξενεμένος όταν τον είδα να αποχωρεί χωρίς την συνήθη αντίδραση του, παρατηρώντας πως τα μάτια του είχαν ανοίξει διάπλατα και το στόμα του είχε μείνει μισάνοιχτο.

«Συνέχισε» μου είπε αποχωρόντας χωρίς να με κοιτάξει.

Είμαι τόσο καλός συγγραφέας; Δεν νομίζω. Πρόσθεσα βέβαια μία εκτεταμένη μάχη μεταξύ των δύο πλευρών της ιστορίας πριν πεθάνει η κόρη του στρατηγού στην ιστορία, για να φανεί πιο ηρωικός ο θάνατος της. Είχα ξεκινήσει να γράφω και τη συνέχεια της ιστορίας. Ήθελα να συνέχιζε με έναν νέο στρατιώτη που είχε τον στρατηγό σαν πατέρα του και την κόρη του σαν αδελφή του να τους βλέπει νεκρούς και να ορκίζεται πως θα πάρει εκδίκηση. Τουλάχιστον αυτό είχα γράψει πριν έρθει ο Μάρκο από πάνω μου να διαβάσει και με διακόψει. Μπορεί στη συνέχεια να το άλλαζα.

«Πόλο» μου απευθύνθηκε ο Μάρκο καθιστός με την πλάτη του γυρισμένη σε μένα. «Σκέφτεσαι ποτέ πως η ζωή κάποιου άλλου προσώπου που ξέρεις μπορεί να είναι όσο σύνθετη και περίπλοκη όσο η δική σου;»

«Τι;» παραξενεύτηκα.

«Συνήθως μένουμε κλεισμένοι στις δικές μας ζωές και ξεχνάμε πως όλοι γύρω μας επίσης ζούνε, έχουν τους δικούς τους προβληματισμούς, τις δικές τους προτιμήσεις, τις δικές τους ιστορίες που λένε όταν κάποιος που δεν έχουν δει μετά από πολύ καιρό τους ρωτάει “Τι έγινε;”»

«Πως σου ήρθε τώρα αυτό;» τον ρώτησα.

«Πόλο» γύρισε να με κοιτάξει «Νιώθεις ποτέ μόνος; Αλλά όχι με την τυπική έννοια» έκανε μια παύση. «Νιώθεις ποτέ πως κανένας δεν έχει την ίδια ζωή με σένα; Και κατά συνέπεια κανείς δεν μπορεί να καταλάβει τι περνάς;»

«Ναι…» του απάντησα μετά από μερικά δευτερόλεπτα σιωπής σκεπτόμενος αν οι εμπειρίες μου ήταν όντως μοναδικές.

«Και… εγώ» μου απάντησε και έστρεψε το βλέμμα του στο πάτωμα. 

«Πως σου ήρθε όμως τώρα αυτό, ο καθένας μας μπορεί να έχει διαφορετική ζωή, αλλά μπορούμε να είμαστε συμπονετικοί ο ένας για τον άλλον»

«Πόλο» σηκώθηκε και με πλησίασε.

«Ναι» του είπα καθώς αυτός κάθισε στο κρεβάτι του και γύρισα την καρέκλα για να μπορώ να τον βλέπω καλύτερα.

«Ποια είναι η άποψη σου για τον πόλεμο» με ρώτησε ο Μάρκο και ξάπλωσε στο κρεβάτι του ανάσκελα, με τα γόνατα του να εξέχουν από το στρώμα.

«Πιστεύω…» έκανα μία μικρή παύση για να μην αφήσω τα συναισθήματα μου να με κυριεύσουν. «Δεν πιστεύω πως έχω δει ποντίκι να κατασκευάζει ποντικοπαγίδα ποτέ στη ζωή μου. Όμως ακούμε συνέχεια ανθρώπους να σκοτώνονται μεταξύ τους επειδή πιστεύουν σε διαφορετικούς θεούς. Βλέπουμε κατοίκους της ίδιας χώρας να βρίσκονται αντιμέτωποι λόγω διαφορετικών πολιτικών απόψεων, φίλους να χωρίζονται επειδή η οικογένεια τους έχει διαφορετική ετικέτα από τις ετικέτες των άλλων.»

Ο Μάρκο γύρισε το κεφάλι του για να με κοιτάξει πιο καλά λες και ήθελε να συγκεντρωθεί περισσότερο στα λόγια μου.

«Έχουμε μάθει ως άνθρωποι να μισούμε κάποιον πριν σκεφτούμε λογικά για μία κατάσταση. Είναι πιο εύκολο εξάλλου, πιο εύκολο και από απλά να ακούς τον άλλον, πιο εύκολο από το να αποδεχτείς ότι φοβάσαι πως είσαι λάθος. Δεν έχω δει ποντίκι να φτιάχνει ποντικοπαγίδα, αλλά έχω δει ανθρώπους να φτιάχνουν συστήματα, ιδεολογίες που εν τέλη τους οδηγούν στη βία. Είναι πάντα “εγώ εναντίον εσένα”, “αριστερός ή δεξιός”, “καλός ή κακός” ανάλογα με το ποιος μιλάει. Οι ηγέτες του πολέμου χτίζουν επιχειρήματα τόσο ισχυρά ώστε να μπορούν να αφήνουν παιδιά… Παιδιά να πεθαίνουν, και να μπορούν να κοιμούνται ήσυχοι τα βράδια. Έχω δει με τα δικά μου μάτια την κακή μεριά του πολέμου»

«Άρα τι πιστεύεις;» σηκώθηκε ο Μάρκο καθιστός μπροστά μου και με κοίταξε κατάματα.

«Πιστεύω πως δεν χρειάζεται να συμφωνήσουμε να μη σταματήσουμε ο ένας τον άλλον. Χρειάζεται απλά να θυμηθούμε πως ο άνθρωπος απέναντι μας δεν είναι ο εχθρός»

«Ενδιαφέρον» είπε ο Μάρκο και έπιασε το πηγούνι του εντυπωσιασμένος. «Ειρωνικό, όμως, δεν νομίζεις;»

«Τι ακριβώς είναι ειρωνικό;» παραξενεύτηκα.

«Είναι ειρωνικό που χρειάστηκε μόνο ένα μήλο να πέσει από ένα δέντρο και ο άνθρωπος ανακάλυψε τη βαρύτητα. Όμως, έχουν πέσει χιλιάδες, αν όχι εκατομμύρια πτώματα, αλλά κανείς δε γνωρίζει ποιο είναι το νόημα της ανθρωπότητας»

«Είναι αρκετά ειρωνικό είναι η αλήθεια»

«Τι πιστεύεις για αυτό;»

«Αλήθεια… αν το καλοσκεφτείς, είμαστε απλά προσωρινές οντότητες σε έναν δανεισμένο κόσμο, οι οποίες προσπαθούν να μάθουν να ζουν για πρώτη φορά»

«Σωστά…»

«Για το σύμπαν σαν σύνολο δεν σημαίνουμε τίποτα, αλλά για εμάς και τους συνανθρώπους μας η ζωές μας είναι το παν»

«Άρα πιστεύεις ότι το νόημα της ζωής είναι να σε θυμούνται; Να γνωρίζεις ανθρώπους οι οποίοι σε εκτιμούν;»

«Δεν ξέρω. Απλά είπα να αναφέρω κάποιες παρατηρήσεις που έχω κάνει όσον αφορά…» κοίταξα έξω από το παράθυρο στα δεξιά μου γιατί νόμιζα πως άκουσα φασαρία και είδα έναν όχλο να μαζεύεται έξω από ένα σπίτι κοντά στο αρχοντικό.

«Τι έγινε;» πλησίασε ο Μάρκο βλέποντας την ξαφνική παύση μου. «Ωχ. Προς τι το πλήθος;»

«Δεν ξέρω» απάντησα στη ρητορική ερώτηση του.

«Θες να πάμε να δούμε;»

«Τι;»

«Πάμε έλα» σηκώθηκε ο Μάρκο με όρεξη να κουτσομπολέψει. «Να αλλάξουμε ρούχα κιόλας πριν φύγουμε»

«Να αλλάξουμε;»

«Ναι…» έκανα μία μικρή παύση και με κοίταξε πλάγια «…να αλλάξουμε»


\chapter{}

Κατεβαίνοντας από τις σκάλες του αρχοντικού και πλησιάζοντας το πλήθος, όλοι οι περαστικοί με χαιρετούσαν. “Μάρκο μου τι κάνεις;”, “Πώς είσαι Μάρκο;”, “Μα καλά μεγάλωσες πολύ” είναι από τις κλασικές ατάκες που ακούω κάθε φορά που κατεβαίνω στην πόλη. Ο πατέρας μου ήταν πολύ καλός με τους μεγαλοαστούς όσον αφορά θέματα οργάνωσης της πόλης μας, με αποτέλεσμα όλοι τους να έχουν σε υψηλή εκτίμηση την οικογένεια μας. Κάτι περίεργο όμως ήταν πως ο Πόλο δεν φάνηκε καθόλου παραξενεμένος με το ποσοστό των ανθρώπων που μου απευθύνονταν καθώς περπατούσαμε. Λες να είχα κάνει λάθος για αυτόν;

Φτάνοντας στον όχλο, χώθηκα ανάμεσα από το πλήθος προσπαθώντας να φτάσω στο θέαμα που όλοι κοιτούσαν στο κέντρο και ο Πόλο με ακολούθησε. 

«Φτιάξτε ένα διάδρομο» άκουσα μία φωνή από το κέντρο και ένιωθα τον κόσμο από τα αριστερά μου να μας σπρώχνει.

«Τι γίνεται» με ρώτησε ο Πόλο.

«Ένας άντρας αυτοκτόνησε» του είπε ένας άλλος άντρας με τον οποίο βρισκόταν ώμο με ώμο.

«Τι;» φώναξα αγχωμένος.

«Έτσι νομίζω τουλάχιστον. Τώρα τον πήραν οι γιατροί και τον μεταφέρουν στο νοσοκομείο» μας είπε ο άντρας και εμείς επιταχύναμε προς το κέντρο του πλήθους.

Δεν ήθελα σε καμία περίπτωση να δω τους γιατρούς να κουβαλάνε τον νεκρό άντρα, αυτό με αηδιάζει. Ήθελα απλά να μάθω γιατί. Γιατί να αυτοκτονήσει κάποιος οριακά από το πουθενά; Η πόλη μας δεν έχει άμεση εμπλοκή με τον πόλεμο. Μήπως είχε χάσει την οικογένεια του; Μήπως είχε ψυχολογικά; Δύσκολο, δεν νομίζω. Η ποιότητα ζωής τα τελευταία χρόνια στην πόλη μας είναι στα ύψη. Γιατί;

Πλησιάζοντας το κέντρο, είδαμε δύο νεαρές κοπέλες που έμοιαζαν φατσικά, αλλά μία ξανθιά και μία μελαχρινή με μπούκλες, να διαβάζουν ένα γράμμα προβληματισμένες.

«Καλησπέρα» τους απευθύνθηκα καθώς εγώ και ο Πόλο της πλησιάζαμε. «Μήπως γνωρίζατε τον άντρα;»

«Μισό λεπτό» τον απέρριψε η μία κοπέλα.

«Πάρτε τον χρόνο σας, δεν θα σας ενοχλήσουμε»

«Όχι δεν μας ενοχλείται κύριε Μπερνούλι, απλά διαβάζουμε ξανά το γράμμα του πατέρα μας για να δούμε αν όντως γράφει αυτό που νομίζουμε»

«Άρα γνωρίζατε τον άντρα»

Η ξανθιά με κοίταξε πλάγια λες και δεν ήταν αυτονόητο από τα λόγια της πως γνώριζαν τον άντρα που.

«Τα συλλυπητήρια μου» της είπα και έσκυψα ελάχιστα το κεφάλι ως ένδειξη σεβασμού.

«Δεν βγάζει νόημα» είπε η αδελφή της.

«Ναι ούτε εγώ καταλαβαίνω γιατί να το κάνει αυτό» συμφώνησε η ξανθιά.

«Θα μπορούσα να ρωτήσω τι σας προβληματίζει ή μήπως είναι προσωπικό θέμα;» τους είπα.

«Ο πατέρας μας ήταν αρκετά μεγάλος ήξερε πως θα πεθάνει σύντομα και έγραψε αυτό το γράμμα πριν από κάποιες μέρες ώστε να μοιράσει την λίγη κληρονομιά του, το σπίτι και τα δύο, τρία χρυσαφικά του δηλαδή, στην οικογένεια του. Όμως, σε εμάς δεν άφησε τίποτα»

«Τι;» είπαμε εγώ και ο Πόλο ταυτόχρονα.

«Ναι καλά ακούσατε. Τα άφησε όλα σε έναν ανηψιό του ονόματι Δούρειους Τίρο, γιο της αδελφής του. Λέει ότι αυτός έχασε το σπίτι του στην άλωση της πόλης βορειοανατολικά από τη δική μας. Θα σταλεί λέει επιστολή που τον καλεί για να κληρονομήσει το σπίτι σύντομα μετά τον θάνατο του, αλλά δεν λέει πως. Τι, δηλαδή, αυτόματα; Δεν γίνεται»

«Και δεν γνωρίζετε αυτόν τον άντρα;» την ρώτησα προσπαθώντας να βρω κάποια λογική βάση στα λεγόμενα της. 

«Όχι βέβαια, πρώτη φορά ακούω το όνομα του» μου απάντησε η ξανθιά και έδωσε το γράμμα στην αδελφή της εκνευρισμένη.

«Άρα ο πατέρας σας δεν αυτοκτόνησε;» 

«Όχι καλέ, γιατί να αυτοκτονήσει;» άκουσα. «Ήταν χαρούμενος άνθρωπος. Χωρισμένος μεν, αλλά δεν είχε κανένα σοβαρό θέμα στη ζωή του, τουλάχιστον απ’ όσο ξέρω»

«Και εσείς που μένετε τώρα;»

«Καλά ανάκριση μας κάνεις;» μου νευρίασε η κοπέλα, πιθανώς επειδή της έκανα πολλές ερωτήσεις. «Με την μητέρα μας και τον αδελφό μας φυσικά, τι σχέση όμως έχει αυτό;»

«Πίστευα πως θα μένατε με τις οικογένειας σας ή κάτι. Δεν έχετε σύζηγο κυρία μου;»

«Μόλις σε είπε μεγάλη» σχολίασε η μελαχρινή ψιθυρίζοντας στην αδελφή της.

«Μόνο εμένα;» είπε η ξανθιά και την κοίταξε; «Κύριε Μπερνούλι δεν μας τα λέτε καλά. Είμαστε πολύ νέες για να κάνουμε οικογένεια. Είμαστε περίπου στην ίδια ηλικία κύριε»

«Αλήθεια;»

«Βεβαίως»

«Μάλιστα…» είπα και κοίταξα τον Πόλο. 

Αυτός κοιτούσε το παράθυρο του σπιτιού του άντρα που είχε πεθάνει. Έστρεψα το βλέμμα μου στο ίδιο παράθυρο αλλά δεν μπόρεσα να καταλάβω τι κοιτούσε, μάλλον είχε αφαιρεθεί.

«Τι σκέφτεσαι;»

«Τίποτα» μου απάντησε.

«Ξέρεις αυτόν τον Δούρειο που λένε οι κοπέλες;»

«Κάτι μου θυμίζει, αλλά δεν είμαι και απολύτως σίγουρος, οπότε μη με φάτε αν κάνω λάθος» μου απάντησε αυτός χαμογελώντας.

«Πέντε μέρες σε ξέρω και ακόμη να καταλάβω πως σκέφτεσαι» παραξενεύτηκα με τη χαρά του σε αυτή την τραγική κατάσταση που είχαμε μπλέξει τους εαυτούς μας.

«Και συ πώς τον ξέρεις; Ξέρεις κάτι παραπάνω που δεν γνωρίζουμε εμείς;» τον ρώτησε με απειλητικό ύφος η μελαχρινή κοπέλα.

«Καταγόμαστε μάλλον από την ίδια πόλη, το όνομα του κάτι μου θυμίζει» της απάντησε ήρεμα αυτός.

«Α… συγγνώμη» του απάντησε αυτή και έκανε ένα βήμα πίσω από την αδελφή της από την ντροπή της.

«Πόλο» του τράβηξα την προσοχή. «Εσύ που είσαι παρατηρητικός, απ’ ότι φαίνεται, ποια είναι η γνώμη σου;»

«Δεν είμαι και ντετέκτιβ για να μπορώ να σκεφτώ γιατί ο πατέρας τους αποφάσισε να δώσει το σπίτι του στον ανιψιό του. Πιστεύω πως πρέπει να αφήσουμε τις κοπέλες στην ησυχία τους και να επιστρέψουμε σπίτι» πρότεινε ο Πόλο καθώς οι άνθρωποι γύρω από το σπίτι είχαν αρχίσει να αποχωρούν από τη στιγμή που αποχώρησαν οι γιατροί.

«Καλά, ό,τι πεις. Κυρίες μου χάρηκα για την γνωριμία» 

«Ποια γνωριμία; Αφού δεν μας ρώτησες καν πως μας λένε κύριε Μπερνούλι» μου σήκωσε το φρύδι η ξανθιά κοπέλα και σταύρωσε τα χέρια της.

«Α… πώς σας…»

«Πλάκα κάνω» με διέκοψε η ξανθιά κοπέλα. «Με λένε Κλάρα και αυτή είναι η αδελφή μου η Γκρέτα» μας συστήθηκαν με μία υπόκλιση.

«Κλάρα…»

«Τι συνέβει κύριε Μπε…»

«Σταμάτα να με φωνάζεις κύριο, είμαστε συνομήλικοι» της είπα εκνευρισμένος. «Συγγνώμη που σας φώναξα» απολογήθηκα αμέσως «εγώ με τον Πόλο θα επιστρέψουμε στο σπίτι μας. Αν υπάρχουν εξελίξεις στην κατάσταση σας ενημερώστε μας την επόμενη φορά που θα συναντηθούμε»

«Βεβαίως» μου απάντησε η ξανθιά. «Ελπίζω να ξανασυναντηθούμε» μας αποχαιρέτησαν οι κοπέλες και αποχώρησαν.

Άρα ο άντρας δεν αυτοκτόνησε. Μάλλον ο περαστικός που μας ανέφερε την αυτοκτονία είχε μάθει ψευδείς πληροφορίες και μας παραπληροφόρησε κατά λάθος. Ο ψίθυρος και το κουτσομπολιό μεταξύ των ανθρώπων είναι επικίνδυνα πράγματα. Ο καθένας μας έχει δικές του προκαταλήψεις, προδιαθέσεις, προτιμήσεις και τα λόγια κάποιου αλλάζουν ελάχιστα από στόμα σε στόμα, με αποτέλεσμα αυτοί που βρίσκονται στις άκρες του δέντρου του κουτσομπολιού να έχουν πολύ διαφορετικές πληροφορίες για τα γεγονότα σε σχέση με άλλους ψηλότερα στην ιεραρχία προς την πηγή της “γνώσης”.

«Τι έγινε Μάρκο; Γιατί βρισκόσασταν έξω;» με ρώτησε ο Άλμπερτ με το που μπήκαμε στο σπίτι, ο οποίος κουβαλούσε στον ώμο του τον παππού του Πόλο.

«Είδαμε ένα πλήθος να μαζεύεται και είπαμε να πάμε να δούμε προς τι όλη η φασαρία» είπα κλείνοντας την κύρια είσοδο του σπιτιού.

«Μάλιστα» μου απάντησε ο Άλμπερτ. «Περάσατε καλά; Τι είδατε»

Εκείνη τη στιγμή κοίταξα τον Πόλο από περιέργεια και τον είδα να κοιτάζει τον γέρο του με το ίδιο χαμόγελο που είχε προηγουμένως, όταν συζητούσα με τις κοπέλες.


\chapter{}

Δεν πήρε τόσο πολύ καιρό τελικά το πρώτο σινιάλο. Δεν ήρθε και γρήγορα, μη παρεξηγηθούμε, αλλά υπολογίζαμε πως θα ερχόταν λίγο αργότερα, ίσως μία μέρα πριν το δεύτερο σινιάλο. Καλά, έτσι κι αλλιώς δεν έχει τόση σημασία, από το δεύτερο σινιάλο και μετά θα κριθούν όλα. Μέχρι τότε, πρέπει να συνεχίσω να συμπεριφέρομαι όπως έχω συμπεριφερθεί ως τώρα, παρόλο που αρχικά ήμουν εναντίον της ιδέας του “Σχεδίου”. Ο αντισυνταγματάρχης είπε θα με καλύψει αν δεν μπορέσω να αποτελειώσω την δουλειά μου, οπότε είμαστε καλά.

Οι επόμενες δύο μέρες πέρασαν πάλι στο δωμάτιο του Μάρκο, με εμένα να γράφω την “ιστορία” μου, την οποία, για να πω την αλήθεια, είχα βαρεθεί λίγο. Αυτός σχεδίαζε τα γνωστά κυκλάκια και βελάκια στα φύλλα που μου ζητούσε από το δικό μου τετράδιο επειδή είχαμε μόνο ένα τετράδιο και βαριόταν να το αναφέρει στους υπόλοιπους για να του πάρουν καινούργιο. Που και που τον ρωτούσα και ερωτήσεις όπως “Γιατί σε ξέρουν όλοι στην πόλη;” ή “Γιατί μόνο έναν μπάτλερ;”, έστω για να κάνουμε μικρές συζητήσεις και να μη καθόμαστε μουγκοί όλη τη μέρα. Οι απαντήσεις του ήταν κάπως απλές: “Ε, λόγω του πατέρα μου” ή “Δεν χρειαζόμαστε πάνω από έναν, γιατί να προσλάβουμε κι άλλον;”. Αυτός δεν με ρώταγε πολλές ερωτήσεις από μόνος του, μόνο απέκρουε τις δικές μου πίσω σε μένα όποτε μπορούσαν οι ερωτήσεις να εφαρμοστούν και στη δική μου περίπτωση. 

Το απόγευμα της δεύτερης μέρας μετά το περιστατικό με τον ηλικιωμένο άντρα, είχα κλειστό το τετράδιο μου και είχα ακουμπήσει το κεφάλι μου στα χέρια μου πάνω στο γραφείο του Μάρκο και κοίταζα την πλάτη του Μάρκο κάτω από τον αριστερό μου αγκώνα.

«Δεν έχεις βαρεθεί;» τον ρώτησα και τον άκουσα να ξεφυσάει.

«Ε, λίγο» μου απάντησε αυτός και ξάπλωσε πάνω στα ακουμπισμένα στο πάτωμα χαρτιά.

«Δεν κάνουμε τίποτα άλλο; Κάθε μέρα τα ίδια» σήκωσα το σώμα μου και τον κοίταξα που καθόταν ανάσκελα.

«Θες να βγούμε πάλι στην πόλη;» 

«Στην πόλη;»

«Ναι γιατί όχι. Μπορεί να πετύχουμε και τις κοπέλες από την προηγούμενη μέρα»

«Γιατί; Σου άρεσαν;»

«Όχι ιδιαίτερα…» τον πρόδωσαν τα κοκκινισμένα μάγουλα του «απλά… μπορεί να έχουν νέα όσον αφορά τον άντρα από την πόλη σου που κληρονόμησε το σπίτι»

«Α, μπορεί» του απάντησα και σηκώθηκα. «Που θα τις βρούμε όμως;»

«Μην αγχώνεσαι για αυτό, λογικά θα τις βρούμε αν είμαστε έξω αρκετή ώρα» σηκώθηκε και αυτός ξεσκονίζοντας την πλάτη του. «Συμφωνείς οπότε;»

«Ναι αμέ»

«Και αν δεν τις βρούμε δεν πειράζει. Θα έχουμε κάνει τη βόλτα μας»

Η πόλη του Μάρκο ήταν αρκετά μεγάλη σε μέγεθος, αλλά τα περισσότερα σπίτια ήταν τεράστια και πολύ πιο εντυπωσιακά από το συνηθισμένο, όχι και σαν το αρχοντικό, αλλά δεν συγκρίνονταν με αυτά που είχαμε στην δική μας μεριά. Σε κάθε γωνία της είχε κόσμο. Δύο γιαγιάδες να πίνουν καφέ στο μπαλκόνι τους και να συζητάνε, ένα γκρουπ μεσήλικων να στέκονται σε μία γωνία και να σπάνε πλάκα μεταξύ τους, πολλά παιδιά με τους γονείς τους να περπατάνε στους δρόμους. Όλοι τους χαρούμενοι προχωρούσαν με τις ζωές τους, αγνοώντας τελείως το γεγονός ότι ένας πόλεμος εξελίσσεται στο παρασκήνιο, μακριά μεν από το σπίτι τους.

Περνώντας από την πλατεία της πόλης, είδαμε ένα μεγάλο αριθμό από αγοράκια πολύ μικρότερα από εμάς να παίζουν ποδόσφαιρο, που μου φάνηκε περίεργο. Νόμιζα πως οι πλουσιότεροι δεν είχαν τέτοια σκληρά χόμπι όπως το ποδόσφαιρο. Η εικόνα μικρών καλοντυμένων παιδιών να κλωτσούν μία μπάλα όμοια με αυτή που έχουμε πίσω στις δικές μας πόλεις ήταν αρκετά περίεργη στο μάτι μου. Τα μισά παιδιά είχαν μάλιστα λερώσει τα καλά τους ρούχα με χώμα και άλλα είχαν σκίσει μέρη των παντελονιών και των παπουτσιών τους. Τι θα έλεγαν οι γονείς τους για αυτό;

«Σου αρέσει το ποδόσφαιρο;» με ρώτησε ο Μάρκο βλέποντας με να παρακολουθώ τον μικρό αγώνα που εξελισσόταν μπροστά μας.

«Έχω παίξει μερικές φορές» το έπαιξα διπλωματικός.

«Θες να παίξουμε μαζί τους για λίγο;»

«Τι;»

«Έλα πάμε» είπε και έτρεξε προς το μέρος των μικρών. «Ε, μπορείτε να μου δώσετε πάσα;»

Το αγοράκι που είχε την μπάλα εκείνη τη στιγμή είδε τον Μάρκο να τους πλησιάζει και χαρούμενος του πλάσαρε την μπάλα στα πόδια του. Εγώ έτρεχα από πίσω του και αυτός όταν πήρε την μπάλα σταμάτησε μπροστά μου και έκανε προσποίηση με το σώμα του προς τα δεξιά και με απέφυγε τσιμπώντας την μπάλα προς τα αριστερά και κάνοντας μισή στροφή με το σώμα του προς την ίδια κατεύθυνση.

«Έλα» τον άκουσα να λέει από πίσω μου και έστριψα να τον κοιτάξω. «Παίξε άμυνα» είπε και τσίμπησε την μπάλα προς τα αριστερά μου.

Τον ακολούθησα κάνοντας πλάγια βήματα έχοντας το τέρμα πάντα στην πλάτη μου. Αυτός προσπαθούσε να επιβραδύνει και να επιταχύνει απότομα ώστε να με κάνει να χάσω την ισορροπία μου, αλλά εγώ δεν ξεγελιόμουν από τις κινήσεις του σώματος του. Βλέποντας ότι δεν μπορούσε να με περάσει, έκανε ένα βήμα πίσω με τη μπάλα και σήκωσε το κεφάλι του κοιτώντας τα παιδιά που βρίσκονταν πίσω μου. Έκανα και εγώ ένα βήμα προς το μέρος του αλλά αυτός με ξεγέλασε και έδωσε πάσα σε ένα αγοράκι που βρισκόταν στα αριστερά του και, πριν προλάβω να αντιδράσω, έτρεξε από την αντίθετη μεριά ζητώντας πάλι την μπάλα. Το αγόρι του έδωσε μια εξαιρετική προωθημένη πάσα που κατέληξε ακριβώς στα πόδια του και αυτός την έστειλε στην αριστερή μεριά του τέρματος με το δεξί του πόδι, εκεί που ο τερματοφύλακας δεν μπορούσε ούτε αν αντιδρούσε εγκαίρως να αποκρούσει την μπάλα, καθώς είχε κάτσει στο δεξί μέρος του τέρματος και παρακολουθούσε την μικρή διαμάχη μας.

«Γκολ» πανηγύρισε ο Μάρκο.

«Δεν ήταν δίκαιο» του είπα «Πάσαρες στο αγόρι εκεί»

«Η ζωή δεν είναι δίκαιη Πόλο» μου είπε παιχνιδιάρικα μέσα στην χαρά του.

«Πάλι καλά που η ζωή είναι άδικη, γιατί αλλιώς θα ήμασταν όλοι κρεμασμένοι σε έναν σταυρό»

«Ναι ναι, παίξτο μας και θρήσκος τώρα» μου είπε πλησιάζοντας με. «Ευχαριστούμε παιδιά, καλή συνέχεια»

«Καλή συνέχεια» του απάντησαν όλα τα παιδιά μαζί και ο τερματοφύλακας έδωσε πάσα σε έναν διπλανό του συμπαίκτη για να συνεχίσουν το παιχνίδι τους.

«Ωχ κοίτα ποιες είναι εδώ» είπε ο Μάρκο καθώς απομακρυνόμασταν από τα αγοράκια δείχνοντας μου τις δύο κοπέλες από την προηγούμενη μέρα να κάθονται και να τρώνε παγωτό. «Σου είπα, δεν θα χρειαστεί να ψάξουμε τόσο»

«Πως είσαι τόσο σίγουρος, ήξερες που μένουν;»

«Όχι»

«Τότε;»

«Αν αναζητάς κάτι, το σύμπαν θα στο φέρει. Αυτό έχω καταλάβει τις τελευταίες μέρες»

«Τι;»

«Καλησπέρα σας Κλάρα, Γκρέτα» πλησίασε τις κοπέλες ευγενικά ο Μάρκο. «Πώς είναι η μέρα σας;»

«Πάλι εσείς;» του είπε η Κλάρα.

«Προς τι ο θυμός κυρία μου;»

«Ειρωνικά το είπα κύριε Μπερνούλι» του απάντησε αυτή αφήνοντας το μπολ από το οποίο μοιράζονταν το παγωτό τους στο παγκάκι.

«Σας είχα πει να μη με φωνάζετε κύριο, κυρία μου» είπε παιχνιδιάρικα αυτός.

«Ω, και πως να σας φωνάζω;»

«Μάρκο» της έπιασε το χέρι και φίλησε την παλάμη της.

Εκείνη τη στιγμή ένιωσα ένα δάχτυλο να μου ακουμπάει απαλά τον ώμο. Γυρίζοντας, είδα την δεύτερη κοπέλα, η οποία είχε το βλέμμα της στραμένο στο έδαφος, ενώ το σώμα της κουνιόταν ελάχιστα πέρα δώθε. 

«Συγγνώμη για την προηγούμενη μέρα» μου απολογήθηκε για κάποιο λόγο.

«Για ποιο πράγμα;»

«Α… δεν το θυμάσαι καν» χάρηκε αυτή με την άγνοια μου.

«Όχι σοβαρά, για ποιο πράγμα;»

«Πες μου ότι δεν θυμάσαι καν το όνομα μου;» 

«Όχι βέβαια, πως να ξεχάσω το όνομα σου, αφού είναι το ίδιο με το όνομα της αδελφής μου»

«Ενδιαφέρον» μου απάντησε αυτή και στάθηκε έκπληκτη μπροστά μου.

Κοίταξα προς τον Μάρκο και τον είδα να συζητάει με την άλλη κοπέλα. Αυτή γελούσε με τα λόγια του, μάλλον της άρεσε ή κάτι. Με είχε αφήσει μόνο μου με την αδελφή της η οποία ήταν μεν χαμογελαστή, αλλά δεν φαινόταν να ήθελε να συζητήσει με εμένα.

«Έχετε κανένα νέο για εκείνον τον ξάδερφο;» την ρώτησα για να ξεκινήσω μία μικρή συζήτηση.

«Ναι βασικά» μου απάντησε. «Ήρθε χθες στην πόλη. Ήταν ψηλός, ξανθός και κατσούφης»

«Τι ηλικία τον έκανες;»

«Πολύ μεγαλύτερος από εμάς σίγουρα, ίσως σαρανταπέντε με πενήντα»

«Όντως;»

«Είναι αυτός που ήξερες;»

«Όχι» της απάντησα και αυτή γέλασε.

«Έ τότε τι μου λες “όντως”;»

«Το εννοούσα… με την αντίθετη έννοια. Περίμενα να μου περιέγραφες κάποιο νεότερο και παραξενεύτηκα»

«Α… εντάξει τότε»

«Ε, εσείς οι δύο» άκουσα την φωνή της ξανθιάς κοπέλας από πίσω μου να μας τραβάει την προσοχή. «Αργότερα, εγώ με την Γκρέτα θα πηγαίναμε μία βόλτα προς την παραλία για να δούμε το ηλιοβασίλεμα. Μήπως θα θέλατε ε»

Παραλία; δεν έχει παραλία εδώ τριγύρω. Κοίταξα τον Μάρκο και φαινόταν και αυτός προβληματισμένος.

«Αφού δεν έχει καμία παραλία κοντά στην πόλη, μόνο την λίμνη δέκα λεπτά νότια» αναρωτήθηκε φωναχτά ο Μάρκο.

«Ναι, αυτό εννοώ, τη λίμνη» του απάντησε η ξανθιά.

«Τότε ως προς τι ο ονοματισμός “παραλία”;» την ρώτησε ο Μάρκο.

«Έχει ένα ωραίο μέρος που κοιτάει δυτικά και είναι σαν παραλία θάλασσας με βότσαλα και ωραία θέα»

«Πόλο εσύ τι λες;» μου απευθύνθηκε ο Μάρκο.

«Θα χρειαστεί να ανάψουμε φωτιά;»

«Όχι καλέ» μου είπε η Γκρέτα. «Δέκα λεπτά περπάτημα προς τα κει είναι ο προορισμός μας και τα φώτα της πόλης είναι αρκετά για να μας κατευθύνουν πίσω»

«Ωραία λοιπόν» πήρε πρωτοβουλία ο Μάρκο «θα σας συνοδεύσουμε μέχρι την λίμνη. Είσαι μέσα;»

Αν μου ξεφύγει κάτι ενώ είμαστε μακριά από την πόλη, δεν θα μπορέσω να σώσω τον εαυτό μου, ούτε το “Σχέδιο”. Ήδη παραλίγο να ξεφύγω ενώ συζητούσα με την μελαχρινή κοπέλα. Από την άλλη, είναι ευκαιρία να αναπτύξω το “Σχέδιο”. Είχα δει την μητέρα του Μάρκο πίσω στο αρχοντικό ήδη να παραπονιέται για την συνύπαρξη μας στο αρχοντικό της, αλλά ο Μάρκο και ο Άλμπερτ την καθησύχασαν λέγοντας της πως είναι πολύ νωρίς για να μπορέσουμε να συνεχίσουμε εγώ και ο Τζιοβάνι με τις ζωές μας. Δεν θα περάσουν πολλές μέρες ακόμη μέχρι το τέλος, οπότε ελπίζω το ρίσκο που θα πάρω τώρα να μην καταστρέψει τα πάντα.

«Τι σκέφτεσαι;» με ρώτησε ο Μάρκο.

«Τίποτα, πάμε»


\chapter{}

«Και πως γνωριστήκατε εσείς οι δύο;» με ρώτησε η Κλάρα δείχνοντας τη μία στιγμή τον Πόλο και την άλλη στιγμή εμένα.

«Είναι μεγάλη ιστορία»

«Δεν έχεις όρεξη να μου πεις, ε;»

«Όχι, όχι»

«Τότε γιατί δεν μου λες;»

«Μία μέρα τον βρήκαμε με την οικογένεια μου στην πίσω μεριά του κήπου μας αγκαλιά με τον παππού του. Στην αρχή μου φαινόταν πολύ περίεργη η κατάσταση, επειδή σκεφτόμουν πως στο καλό βρέθηκαν στο σπίτι μετά από αρκετές μέρες της άλωσης της πόλης του από τους αντιστασιακούς»

«Τι σημασία έχει αυτό;»

«Δεν έχει… Ίσως εγώ ήμουν υπερβολικός και προσπαθούσα να βρω στοιχεία εκεί που δεν υπήρχαν για να αποδείξω ότι οι δυο τους ήταν εισβολείς»

«Ναι αλλά γιατί;»

«Έχω κάτι στο νου μου, αλλά αυτό όντως δεν έχω όρεξη να στο πω» της είπα έχοντας την εικόνα του Άλμπερτ να μου αφηγείται τον θάνατο του πατέρα και της αδελφής μου.

«Καλά, ό,τι νομίζεις» μου θύμωσε αυτή.

«Ξέρεις, έχεις ίδιο όνομα με την αδελφή μου»

«Το γνωρίζω. Πιο περίεργο θα ήταν να μη το ήξερα» μου απάντησε και η κοιλιά μου γουργούρισε. «Κάποιος πείνασε» μου είπε.

«Όχι…»

«Δεν πεινάς;»

Στάθηκα για λίγο.

«Ε, Πόλο» του φώναξα καθώς είχε περατήσει αρκετά πιο μπροστά από εμάς συζητώντας με την Κλάρα. «Θέλετε να φάμε τίποτα;»

Αυτός έκανε μια γρήγορη κίνηση με το δεξί του χέρι προς την τσέπη του. Αγχωμένος, άρχισε να ψάχνει βιαστικά και τις δύο τσέπες του για κάτι, σηκώνοντας ακόμη και το κοστούμι του για να μπορέσει να ελέγξει και με τα μάτια του.

«Τι ψάχνεις;» τον ρώτησα.

«Το βρήκα. Το είχα βάλει στην πίσω τσέπη. Το σακούλι με τα νομίσματα μου, πάλι καλά»

«Δε θα το χρειαστείς. Θα σας κεράσω όλους εγώ»

«Αλήθεια;» ενθουσιάστηκε η Κλάρα.

«Βεβαίως»

Τα περισσότερα μαγαζιά βέβαια που περνούσαμε στην άκρη της πόλης ήταν κλειστά εκείνη την ώρα. Το μόνο που είδα ήταν ένας αρτοπωλητή με το κάρο του. Δεν ήταν και η καλύτερη επιλογή γεύματος, αλλά όλοι πιστεύω θα συμβιβάζονταν με ένα κομμάτι ψωμί ο καθένας μας για την ώρα. Εξάλλου, οι κοπέλες είχαν τσιμπήσει και το παγωτό τους πριν τις συναντήσουμε, που τώρα που το σκέφτομαι, που άφησαν το μεταλλικό μπολ που κρατούσαν; Τέλος πάντων, είμαστε λίγο μακριά τώρα για να γυρίσουμε πίσω. Δεν πειράζει, σίγουρα έχουν κι άλλα στο σπίτι τους.

«Συγγνώμη» πλησίασα τον αρτοπωλητή «δεν είναι πολύ αργά για να πουλάτε ψωμί;»

«Περίμενα έναν γνωστό μου από την άλλη μεριά της πόλης σήμερα, αλλά δεν θα με επισκεφτεί απ’ ότι φαίνεται» μου απάντησε αυτός. «Τι σας φέρνει εδώ κύριε Μπερνούλι;»

«Άντε πάλι “κύριε”. Μήπως θα μπορούσατε να μας κόψετε μισή φρατζόλα για να μοιραστούμε όσο περπατάμε;»

«Ευχαρίστως» μου απάντησε αυτός.

«Μάλλον…» τον σταμάτησε η Κλάρα «…μπορείτε να μας δώσετε ολόκληρη φρατζόλα; Μπορεί να πετύχουμε πάπιες στην λίμνη. Δεν γίνεται να μην έχουμε τίποτα να της ταΐσουμε»

«Έχεις ένα δίκιο» της απάντησα.

«Να κάνω αυτό που είπε η νεαρή;» στάθηκε με το μαχαίρι του πάνω από τη φρατζόλα ο αρτοπωλητής περιμένοντας την απάντηση μου.

«Βεβαίως» του είπα και αυτός έβαλε την φρατζόλα σε μία χάρτινη σακούλα και μου την έδωσε κατευθείαν στο χέρι. «Ορίστε» είπα και έκοψα ένα κομμάτι στην Κλάρα.

«Ω, τι ευγενικός»

«Ορίστε και σε σας» έδωσα άλλο ένα κομμάτι στην αδελφή της και στον Πόλο, του οποίου τα μάτια άνοιξαν διάπλατα όταν είδε το ψωμί στην παλάμη του.

«Ευχαριστούμε Μάρκο» μου είπε ο Πόλο και πήρε μία δαγκωνιά από το κομμάτι του.

«Πάμε τώρα» είπε η κλάρα μπουκωμένη. «Θα χάσουμε το ηλιοβασίλεμα»\newline\newline

Φτάνοντας στη λίμνη, οι κοπέλες έτρεξαν μπροστά από εμάς. Σχολίαζαν συνέχεια στη διαδρομή πως ήθελαν να βουτήξουν τις πατούσες τους στο νερό πριν κάτσουμε όλοι μαζί στα βότσαλα. Ο Πόλο και γω δεν τις σταματήσαμε. Καθίσαμε κοντά στο νερό, αλλά όχι αρκετά κοντά για να μη μας πλατσουρίζουν οι κοπέλες που έτρεχαν με το ψωμί που είχαμε αγοράσει ψάχνοντας για πάπιες. Ο ήλιος ήταν αρκετά ψηλά ακόμη, μπορεί και να είχαμε άλλη μισή ώρα για το ηλιοβασίλεμα, οπότε χαλαρώσαμε στις θέσεις μας απολαμβάνοντας την θέα.

«Ξέρεις» μου είπε ο Πόλο. «Γιατί πρέπει να παλεύουμε;»

«Πώς σου ήρθε τώρα αυτό;»

«Είναι παράλογο. Κανένας σε ολόκληρο τον κόσμο δεν αξίζει να πληγώνεται. Θέλεις να μου πεις πως ο άνθρωπος είδε αυτό…» είπε δείχνοντας τις φίλες μας που έτρεχαν στην άκρη της λίμνης με τη λίμνη στο παρασκήνιο και ακόμη πιο πίσω τα ψηλά δέντρα του δάσους της αντίπερα όχθης «…και εφηύρε όπλα, εκρηκτικά, ξίφη, τρόπους για να μπορέσει να σκοτώσει τον συνάνθρωπο του»

«Έχεις δίκιο» του είπα και ξανακοίταξα το τοπίο για να νιώσω στην καρδιά μου τις λέξεις του. «Ούτε εγώ μπόρεσα να το καταλάβω ποτέ αυτό. Λέμε “οι αντιστασιακοί είναι κακοί” ή “οι αριστοκράτες είναι οι καλοί”. Αυτές οι πεποιθήσεις υπερισχύουν καθόλη τη διάρκεια του πολέμου, τουλάχιστον από την πλευρά μας. Αλλά τα παιδιά που δεν έχουν ζήσει πόλεμο και τα παιδιά που δεν έχουν ζήσει περίοδο ειρήνης έχουν διαφορετικές αξίες, χωρίς να εξαιρώ και τον εαυτό μου» είπα και ο Πόλο έστρεψε το βλέμμα του προς το μέρος μου με το που ολοκλήρωσα την τελευταία μου πρόταση. «Αυτοί που στέκονται στην κορυφή, βέβαια, καθορίζουν τι είναι σωστό και τι όχι. Στην τελική λένε “η δικαιοσύνη θα επικρατήσει”, το οποίο φαίνεται προφανές. Όποιος κερδίσει τον πόλεμο θα γίνει η δικαιοσύνη…»

«Δεν μπορούσες να έχεις περισσότερο δίκιο…» μου είπε και πέταξε ένα βότσαλο στην λίμνη.

«Είπες, όμως, ότι κανείς δεν αξίζει να πληγώνεται. Είναι, όμως, όλοι οι άνθρωποι ίσοι;»

«Τι εννοείς;»

«Δηλαδή, θα έβαζες στην ίδια κατηγορία έναν δολοφόνο με ένα παιδί; Ή με έναν μαγαζάτορα;»

«Όχι, είναι τελείως διαφορετικές περιπτώσεις. Σε μερικούς αρέσει να λένε πως ο Θεός δεν δημιουργεί κανέναν άνθρωπο πάνω από κανέναν άλλο. Όμως αυτό είναι χαζό. Πιστεύω πως όλοι μας είμαστε ίσοι στην γέννα μας, όμως οι πράξεις μας κατά τη διάρκεια της ζωής μας κρίνουν αν είμαστε πραγματικά “καλοί” ή όχι»

«Άρα τι πιστεύεις είναι καλύτερο. Κάποιος που γεννήθηκε “καλός” ή κάποιος που ξεπερνά την “κακή” του φύση μέσω των πράξεων του;»

«Όπως είπα, δεν υπάρχει κανείς που γεννιέται πραγματικά “καλός”»

«Άρα πιστεύεις πως όλοι μας γεννιόμαστε σε ένα “μεσαίο επίπεδο” και πρέπει όλοι μας να προσπαθούμε να γίνουμε η καλύτερη εκδοχή του εαυτού μας;»

«Ναι…» πήρε άλλο ένα βότσαλο «…μπορείς να το πεις και έτσι» είπε και το πέταξε στη λίμνη.

«Ενδιαφέρον»

«Τι συζητάμε;» ήρθαν προς το μέρος μας η Κλάρα με την αδελφή της και κάθισαν αριστερά και δεξιά μας.

«Τι πιστεύεις, υπάρχουν καλοί και κακοί άνθρωποι; Ή καλά και κακά πράγματα που συμβαίνουν στον κόσμο γενικότερα;» τη ρώτησα.

«Πιστεύω πως η προοπτική παίζει μεγάλο ρόλο στην απόφαση κάποιου» μου απάντησε η Κλάρα. «Μπορείς να παραπονιέσαι ότι τα τριαντάφυλλα έχουν αγκάθια, ή μπορείς να χαίρεσαι που οι θάμνοι με αγκάθια έχουν τριαντάφυλλα. “Κακά” πράγματα συμβαίνουν σε “καλούς” ανθρώπους επειδή συνέχεια συμβαίνουν πράγματα, έτσι δουλεύει ο κόσμος. Είναι στα χέρια των ανθρώπων να κρίνουν αν αυτά που συμβαίνουν είναι “καλά” ή “κακά”»

Έστρεψα την προσοχή μου στον Πόλο, ο οποίος είχε μαζέψει κάπως το σώμα του και δεν ήταν πλέον χαλαρός όπως πριν. Από τα μάτια του κύλησαν λίγα δάκρυα τα οποία χτύπησαν στα βότσαλα και άφηναν βρεγμένα σημάδια.

«Τι σκέφτεσαι Πόλο;»

Αυτός δεν μου απάντησε.

«Είσαι καλά;» τον ρώτησε η αδελφή της Κλάρας πλησιάζοντας τον και ακουμπώντας τον στο γόνατο συμπονετικά. 

Αυτός σηκώθηκε εκνευρισμένος.

«Νομίζω… πρέπει να επιστρέψουμε σπίτι Μάρκο»


\chapter{}

«Μπορείς να χαίρεσαι που οι θάμνοι με αγκάθια έχουν τριαντάφυλλα» οι χθεσινές λέξεις της κοπέλας επαναλαμβάνονταν στο μυαλό μου συνέχεια.

Γιατί ο δικός μου θάμνος με αγκάθια στον οποίο έχω μπλεχτεί δεν έχει λουλούδια; Το “Σχέδιο”, είχε πει ο στρατηγός, είναι για να τελειώσει τον πόλεμο και να επικρατήσει επιτέλους δικαιοσύνη. Όμως, όντως, όποιος κερδίσει τον πόλεμο θα είναι η δικαιοσύνη. Όλα είναι υποκειμενικά, ακόμη και το σωστό και το λάθος, το καλό και το κακό. Υπάρχει μία αντικειμενική δόση στα πάντα, αλλά οι προτιμήσεις και οι προκαταλήψεις μας κάνουν να κλίνουμε διαφορετικά ο ένας από τον άλλο.

Και πες πως όντως είναι υποκειμενικά όλα, και δεν υπάρχει τίποτα αντικειμενικό. Αυτό που δεν μου αρέσει στον πόλεμο είναι πως δε μπορούμε να χτίσουμε μία γκρι ζώνη, μία ενδιάμεση λύση με την οποία μπορούν να συμβιβαστούν και οι δύο πλευρές. Για να έχει διαρκέσει τόσα χρόνια ο πόλεμος, προφανώς και οι ηγέτες των δύο πλευρών θα έχουν προσπαθήσει να συζητήσουν πάρα πολλές φορές χωρίς αποτέλεσμα. Συνεχίζουμε όμως να παλεύουμε, επειδή αυτό έχουμε μάθει. Αν κάποιος δεν συμφωνεί να θυσιάσει κάτι για να συζήσει με εμάς, τον εξουδετερώνουμε, είτε μεταφορικά, είτε κυριολεκτικά, και το αντίθετο. Μερικοί άνθρωποι δεν είναι φτιαγμένοι να συζούν.

«Πόλο» άκουσα τη φωνή του Μάρκο να ηχεί πίσω από την πόρτα του χολ.

Είχα καθίσει στο τραπέζι του δωματίου που βρισκόμασταν την πρώτη μέρα και περιμέναμε τον Άλμπερτ και την κυρία του αρχοντικού να μας ετοιμάσουν το “δωμάτιο” μας. Ήθελα να περάσω λίγο χρόνο μόνος μου σήμερα και να σκεφτώ. Να σκεφτώ αν αυτό που γίνεται είναι σωστό, αν το “Σχέδιο” είναι όντως ο τέλειος τρόπος να κερδίσουμε τον πόλεμο, αν μπορούν να σωθούν περισσότερες ζωές, έστω και αυτές οι ελάχιστες που είμαστε σταλμένοι να εξουδετερώσουμε. Ίσως είμαι πολύ νέος και οπτιμιστής. Ίσως αυτοί που ζουν περισσότερα χρόνια στις πρώτες γραμμές του πολέμου γνωρίζουν περισσότερα από εμένα, αλλά οι σκέψεις μου δεν είναι σε κανένα σενάριο λάθος. Γιατί να χρειαστεί να πεθάνουν όλοι οι στόχοι μας και όχι, έστω, αυτοί που δεν μπορούν να συμβιβαστούν με τη συνεργασία; Τη συνεργασία των δύο πλευρών του πολέμου.

«Εδώ είσαι;» άνοιξε την πόρτα ο Μάρκο και με πλησίασε. «Τι κάνεις εδώ μόνος σου»

Δεν του απάντησα. Τον κοίταξα για λίγο σκεπτόμενος τα γεγονότα των προηγούμενων ημερών.

«Πες μου ότι σε πείραξε αυτό που είπε η Κλάρα χθες»

«Όχι…» του απάντησα «…δεν είναι αυτό απαραίτητα»

«Κοίτα» μου είπε και κάθισε στη διπλανή καρέκλα. «Ξέρω δεν έπρεπε να κάνουμε εκείνη τη συζήτηση χθες, ίσως δεν ταίριαζε στην περίσταση, γνωρίζοντας και το παρελθόν σου»

«Ναι…» του είπα γελώντας.

«Φταίμε και οι δύο εν μέρη, και εσύ που την άρχισες, αλλά και εγώ παραπάνω που τη συνέχισα. Ελπίζω να μην είπα κάτι χθες που να σε πείραξε»

«Όχι…» τον κοίταξα για μια στιγμή στα μάτια «…όχι, όχι. Εσύ δεν φταις σε τίποτα. Έχω απλά… κάτι στο μυαλό μου τις τελευταίες μέρες»

«Θες να το συζητήσουμε;» μου έπιασε την παλάμη του χεριού που είχα ακουμπισμένο στο τραπέζι.

«Ό… όχι, δεν είναι τόσο σημαντικό»

«Μα για να λες ότι σε ταλαιπωρεί τις τελευταίες μέρες λογικά θα είναι»

«Δεν πειράζει. Θα το ξεχάσω εξάλλου σε λίγες μέρες» του είπα και τράβηξα αργά την παλάμη μου προς τα πίσω.

«Όπως νομίζεις» μου απάντησε και σηκώθηκε. «Αν αλλάξεις γνώμη και θελήσεις να περάσεις από το δωμάτιο μου αργότερα, η πόρτα είναι ανοιχτή»

«Έρχομαι μαζί σου» του είπα και τον ακολούθησα στο χολ.

Καθώς ανεβαίναμε τις σκάλες, ακούσαμε έναν χτύπο στην πόρτα της κύριας εισόδου του αρχοντικού.

«Το αναλαμβάνω εγώ» είπε ο Άλμπερτ που άρχισε να κατεβαίνει τις σκάλες εκείνη τη στιγμή τρέχοντας για να φτάσει γρήγορα στην είσοδο. 

Εμείς κοιταχτήκαμε για ένα δευτερόλεπτο και μετά τον ακολουθήσαμε. Ανοίγοντας την πόρτα, αντικρίσαμε έναν αριστοκράτη στρατιώτη, ο οποίος κρατούσε στο χέρι του μερικές σελίδες χαρτί με γραμμένα μερικά ονόματα και στην τσέπη του στήθους της στολής του είχε κρεμασμένη μία πένα. 

«Καλημέρα σας» τον καλωσόρισε ο Άλμπερτ.

«Καλημέρα, δεν θα κάτσω για πολύ. Ήρθα να σας ενημερώσω πως οι πρώτες γραμμές του πολέμου καταστράφηκαν σε μία δύσκολη για εμάς μάχη χθες το βράδυ. Είμαι από τους ελάχιστους που καταφέραμε να επιβιώσουμε  αυτή τη μάχη χωρίς σοβαρό τραυματισμό και στάλθηκα στο αρχοντικό σας κύριε Μπερνούλι…» έστρεψε το βλέμμα του στον Μάρκο «…για να σας ενημερώσω πως χρειαζόμαστε επειγόντως νέους ανθρώπους οι οποίοι είναι διατεθειμένοι να ηγηθούν του στρατού μας, ή να γίνουν υψηλόβαθμοι συνεργάτες μας»

«Τι;» αναφώνησε ο Μάρκο.

«Ακριβώς αυτό που ακούσατε. Το ιστορικό της οικογένειας σας με τον πρώην στρατηγό, καθώς και η θητεία του πατέρα σας, σας κάνουν πολύ καλό υποψήφιο για την ηγεσία του στρατού. Πλέον, με ανορθόδοξο τρόπο βέβαια, οι αιτήσεις της οικογένειας σας μπορούν να γίνουν πραγματικότητα» είπε ο στρατιώτης και σήκωσε το χαρτί που κρατούσε για να το ακουμπήσει στον τοίχο. «Τι λέτε λοιπόν, κύριε Μπερνούλι;» τον ρώτησε ο στρατιώτης περιμένοντας με την πένα μερικά χιλιοστά πάνω από το χαρτί.

«Πόσο… πόσο καιρό μπορείτε να μου δώσετε για να αποφασίσω;» τον ρώτησε πίσω ο Μάρκο, με την φωνή του να τρέμει φοβισμένη.

«Κατά το μέγιστο μία εβδομάδα. Οι στρατιώτες μας έχουν υποχωρήσει τρεις πόλεις νότια από τη δική σας. Εγώ και ο άλλος αγγελιοφόρος που έχουμε σταλθεί να ενημερώσουμε όλα τα αρχοντικά ψάχνοντας νέους αρχηγούς πιστεύουμε πως θα έχουμε ενημερώσει όλα τα αρχοντικά μέχρι αύριο το απόγευμα. Αν αποφασίσετε πως θέλετε να συμμετάσχετε στον πόλεμο, ξέρετε που να μας βρείτε»

«Θα… θα το σκεφτώ» είπε ο Μάρκο και το βλέμμα του στράφηκε κάτω.

«Χαίρομαι πολύ. Χαίρομαι που επίσης τα αρχοντικά φιλοξενούν ανθρώπους»

«Τι εννοείτε;» παραξενεύτηκε ο Άλμπερτ.

«Εσύ…» είπε ο στρατιώτης δείχνοντας με «Εσύ δεν είσαι μέλος της οικογένειας ούτε μπάτλερ;»

«Όχι» του απάντησα.

«Να βλέπεις. Και το άλλο αρχονιτκό από το οποίο πέρασα φιλοξενούσε δύο άντρες, έναν ηλικιωμένο και έναν λίγο πιο νέο, ίσως τριαντάχρονο. Ο δικός τους ηλικιωμένος βέβαια φαινόταν ετοιμοθάνατος. Ελπίζω ο δικός σας να είναι κάπως καλύτερα»

«Βρίσκεται σε καλή κατάσταση» του απάντησε ο Άλμπερτ κοφτά. «Οι δύο ξένοι που φιλοξενούμε αυτές τις μέρες είναι θύματα της άλωσης της πόλης βορειοανατολικά από τη δική μας»

«Ενδιαφέρον. Και το κάνετε από συμπόνια ή…»

«Αυτός είναι ένας από τους λόγους…» του απάντησε ο Άλμπερτ. «Αν δεν έχετε κάτι παραπάνω, να μη σας κρατάμε χωρίς λόγο»

«Ωραία. Πάντως αν θέλετε να πολεμήσετε και εσείς κύριε δεν υπάρχει κανένα απολύτως θέμα» μου απευθύνθηκε ο στρατιώτης και εγώ απέφυγα να τον κοιτάξω στα μάτια. «Καλά… πάω να συνεχίσω να ενημερώνω και τους υπόλοιπους στη λίστα μου. Μπορείτε να μου κάνετε όμως μία χάρη;»

«Βεβαίος» του είπε ο Άλμπερτ.

«Προσπαθούμε να κρατήσουμε την κατάσταση μυστική για λίγες μέρες ακόμη. Πιστεύω πως σε τρεις ή τέσσερις μέρες θα βγει ανακοίνωση αναζήτησης νέου στρατεύματος για την απόπειρα αντεπίθεσης μας. Αν μπορείτε, μέχρι τότε, κρατήστε τη συζήτηση μας σε χαμηλό προφίλ»

«Μάλιστα. Ευχαριστούμε για την ενημέρωση σας. Καλή σας συνέχεια»

«Καλή συνέχεια φίλτατοι» απάντησε ο στρατιώτης στον Άλμπερτ και αυτός έκλεισε την πόρτα.

«Μάρκο πως νιώθεις;» τον ρώτησε ο Άλμπερτ.

Αυτός δεν του απάντησε.

«Πως νιώθεις που το όνειρο σου θα γίνει πραγματικότητα;»

Αυτός δεν του απάντησε.

«Μάρκο;»

Ακούγοντας το όνομα του, αυτός έτρεξε προς τον πρώτο όροφο ανεβαίνοντας τις σκάλες βαριανασαίνοντας και κρύβοντας το πρόσωπο του με το αριστερό του χέρι. Κοίταξα τον Άλμπερτ, αυτός είχε ένα περίεργο χαμόγελο στο πρόσωπο του. Τι το αστείο είχε η αποστολή του αφέντη του στις πρώτες γραμμές; Έβρισκε ευχάριστη την τρέχουσα κατάσταση;

«Συνεχίστε ό,τι κάνατε κύριε, αφήστε τον Μάρκο να σκεφτεί τι συνέβει μόλις και θα το συζητήσω μαζί του όταν δεν υπερισχύει το συναίσθημα μέσα του» μου είπε ο Άλμπερτ και αποχώρησε από την είσοδο με προορισμό την κουζίνα.

Δεν ήθελα όμως τον Μάρκο μόνο του να σκεφτεί Ήθελα να μάθω την αιτία του ξαφνικού ξεσπάσματος του. Ανέβηκα, λοποίν, τις σκάλες του αρχοντικού και στρίβοντας το κεφάλι μου δεξιά προς τα δωμάτια μας, είδα τον μάρκο καθισμένο απέναντι από την πόρτα του δωματίου του. Είχε πάρει κουλουριασμένη πόζα, κρατώντας με τα χέρια του τα γόνατα του. Το κεφάλι του ήταν ακουμπισμένο ανάμεσα από τα δύο πόδια του και το στήθος του ανεβοκατέβαινε άρρυθμα. Πλησιάζοντας τον άρχισα να ακούω καλύτερα τις αναπνοές του. Ο Μάρκο έκλεγε, δεν γελούσε. Η εισπνοές του ήταν πολύ μεγαλύτερες σε διάρκεια από τις εκπνοές, καθώς οι δεύτερες ήταν πιο κοφτές. Ο περίεργος τρόπος που ανέσαινε τον είχε κάνει το κεφάλι του να κοκκινίσει, που προφανώς σήμαινε και πως οι παλμοί του επίσης είχαν ανέβει.

«Εδώ είσαι;» τον πλησίασα και γονάτισα δίπλα του.

Δεν μου απάντησε.

«Είσαι… εντάξει;»

«Πως να είμαι εντάξει. Ο χειρότερος μου φόβος θα γίνει σε λίγες μέρες πραγματικότητα» μου απάντησε εκνευρισμένος αυτός.

«Ο χειρότερος σου φόβος;» παραξενεύτηκα. «Εσύ δεν σχεδιάζεις στρατηγικές για πλάκα; Και ο στρατιώτης είπε πως κάνατε κάποιες αιτήσεις ή κάτι παρόμοιο, οπότε ήθελες να μπεις στον στρατό ως υψηλόβαθμος έτσι κι αλλιώς, όχι;»

«Αυτά είναι οι χαζομάρες τις μητέρας. Νόμιζε πως, επειδή ο πατέρας ήταν ο προηγούμενος στρατηγός , μας είχε ταχθεί η επόμενη αρχηγεία όταν θα γινόμουν δεκαοκτώ. Εγώ δεν εξέφρασα ποτέ επιθυμία για κάτι τέτοιο»

«Ίσως οι δραστηριότητες σου στον ελεύθερό σου χρόνο την πείσανε και αυτή το πήρε λίγο πιο σοβαρά απ’ ότι έπρεπε»

«Το να φτιάχνεις στρατηγικές από την άνεση του σπιτιού σου είναι πολύ πιο εύκολο από το να βρίσκεσαι στο μέτωπο του πολέμου, όλοι μπορούν να κάνουν το πρώτο…»

«Μη το λες» του είπα και κάθισα δίπλα του, κρατώντας μία μικρή απόσταση για να μην τον ακουμπήσω κατά λάθος και νευριάσει παραπάνω.

«Πόλο»

«Ναι;»

«Φοβάμαι» είπε και κουλουριάστηκε πάλι ανάμεσα στα πόδια του.

«Τι εννοείς;»

«Φοβάμαι τον θάνατο. Δεν δίνω δεκάρα για το γεγονός ότι θα έχω υπερβολικές ευθύνες ως υποψήφιος στρατηγός ή ότι θα βρίσκομαι μακριά από το σπίτι μου. Το μόνο που με φοβίζει είναι το άγνωστο, το ανεξερεύνητο, αυτό από το οποίο κανένας άνθρωπος δεν μπόρεσε να επιστρέψει…»

Επικράτησε μία στιγμή σιωπής.

«Γιατί ο θάνατος είναι κάτι που πρέπει να φοβόμαστε;» του είπα κοιτώντας το ταβάνι και σκεπτόμενος όλες τις ζωές που έχουν χαθεί σε αυτόν τον πόλεμο. «Ζούμε μόνο και μόνο επειδή δε θέλουμε να πεθάνουμε; Μας προσφέρει ποτέ κάτι καλό η ύπαρξη μας σε αυτόν τον κόσμο;» είδα με την άκρη του ματιού μου τον Μάρκο να γυρίζει το κεφάλι του προς το μέρος μου. «Εμένα δεν μου έχει προσφέρει κάτι…»

«Δεν…» ξεκίνησε να πει κάτι ο Μάρκο και σταμάτησε. Άπλωσε τα πόδια του στο πάτωμα και τοποθέτησε τα χέρια του πάνω στα μπούτια του. «Δεν νομίζω ότι φοβάμαι τον θάνατο αυτό καθαυτό, αλλά φοβάμαι σίγουρα την σκέψη ότι δεν θα προλάβω να γίνω ο άντρας που ονειρεύομαι»

«Δηλαδή;»

«Νομίζω είναι προφανές αυτό που εννοώ…» μου είπε και με κοίταξε πλάγια.

Επικράτησε άλλη μία στιγμή σιωπής.

«Θυμάσαι που σου είχα πει πως νιώθω μόνος;» με ρώτησε.

«Ναι, πως κανείς δεν έχει τις ίδιες εμπειρίες με εσένα, οπότε δεν μπορούν να περάσουν από το μυαλό του οι σκέψεις σου ή κάτι τέτοιο»

«Ναι…» χαλάρωσε το κεφάλι του στον τοίχο. «Πόλο, αν είχες μία ευχή αυτή τη στιγμή… τι θα επιθυμούσες;»

«Αυτή τη στιγμή;... Μάλλον να μπορούσα να δω την αδελφή και τον πατέρα μου μία τελευταία φορά…»

Ο Μάρκο με κοίταξε πάλι πλάγια, αυτή τη φορά με τελείως διαφορετικό ύφος. 

«Και γω…»


\chapter{}

Όλη την υπόλοιπη μέρα, καθώς και την επόμενη, ο Μάρκο βρισκόταν στο δωμάτιο του. Δεν είχε κατέβει στην τραπεζαρία ούτε για πρωινό, ούτε για μεσημεριανό, ούτε για βραδινό και από πλευράς μας δεν τον είχε ενοχλήσει κανένας, μέχρι να τον προσεγγίσω εγώ το δεύτερο βράδυ, ακριβώς μετά το γεύμα μας με την κυρία του αρχοντικού. Δεν μπορούσα να αντέξω άλλο τη μουρμούρα της για τον γιο της, οπότε πήρα την απόφαση να περάσω από το δωμάτιο του μετά το βραδινό. Εξάλλου, όλα θα είχαν κριθεί μέχρι αύριο.

Έτσι, όταν ανέβασα τον Τζιοβάνι στο δωμάτιο μας, έκανα μία μικρή λοξοδρόμηση όταν τον ακούμπησα στο κρεβάτι μας για τον βραδινό του ύπνο και στράφηκα προς το δωμάτιο του Μάρκο

Αυτός καθόταν στην καρέκλα του γραφείου του σκεπτόμενος. Αγνάντευε την πόλη από το δεξί παράθυρο του δωματίου έχοντας πάρει μία χαλαρή πόζα, ξαπλωμένος οριακά στην καρέκλα, στερεώνοντας το κεφάλι του στο δεξί του χέρι. Τα δάχτυλα του αριστερού του χεριού ανεβοκατέβαιναν πανω στα σταυρωμένα του πόδια σαν να ήταν παλμοί από μικρά κύματα. Κάθισα στο κρεβάτι του και τον κοίταξα. Αυτός δεν είχε γυρίσει καθόλου για να με κοιτάξει όσο βρισκόμουν στο δωμάτιο του.

«Τι σκέφτεσαι;» τον ρώτησα.

«Έχω αυτό το περίεργο αίσθημα… ότι θα πεθάνω σύντομα» μου απάντησε αυτός.

«Τι;» παραξενεύτηκα.

«Τι θα γίνει αν… δεν τα καταφέρω; Τι θα γίνει αν… αποτύχω;»

«Τι θα γίνει αν δεν αποτύχεις;» του απάντησα. «Μη σε τρώει η σκέψη της αποτυχίας. Στο λέω, είμαι και εγώ πεσιμιστής. Ένας σοφός μου είχε πει πως τα πλοία είναι πάντα ασφαλή όταν κάθονται στο λιμάνι, αλλά αυτός δεν είναι ο σκοπός που φτιάχτηκαν» είπα και ο Μάρκο γύρισε ελάχιστα το κεφάλι του προς το μέρος μου και με κοίταξε.

«Ναι, ε;...» μου απάντησε αυτός και έστρεψε πάλι του βλέμμα του στην πόλη.

«Ξέρεις, δεν θα υπάρχουν πάντα άνθρωποι γύρω σου για να σε παρακινούν, όπως τώρα. Μερικές φορές πρέπει να πάρεις αποφάσεις μόνος σου»

«Όταν ήμουν παιδί…» ξεκίνησε να λέει κάτι φαινομενικά άσχετο ο Μάρκο «…πάντα αναρωτιόμουν τι είδους άντρας θα γινόμουν στο μέλλον»

«Τι;»

«Ποτέ δεν πίστευα πως θα καταλήξω έτσι. Από μικρός, το μόνο που ήθελα να επιτύχω είναι να είμαι ελεύθερος» είπε και έστρεψα το βλέμμα μου προς αυτόν. «Ήθελα να είμαι ελεύθερος, να μη μου ασκείται κριτική από κανέναν, να μη χρειάζεται να νοιάζομαι για κανέναν, να μπορώ να ζω χωρίς να επηρεάζω τη ζωή κανενός με τις πράξεις μου, και το αντίστροφο. Να μπορώ να ξεφύγω από την κοινωνική συνήθεια που επικρατεί. Μικρός νόμιζα πως αν βάλω τέλος στον πόλεμο, θα μπορούσα να χτίσω το όνειρο μου, όμως, τώρα που μου δίνεται η ευκαιρία… δεν πιστεύω πως έχω το θάρρος να θυσιάσω τις απαραίτητες ζωές για να το κάνω, τις ζωές των αντιστασιακών που εδώ και τόσα χρόνια πολεμάνε και αυτοί με στόχο την ελευθερία τους. Ένα διαφορετικό είδος ελευθερίας που θα τους εξασφάλιζε την ανεξαρτησία τους από την αριστοκρατική κυβέρνηση που επικρατούσε πριν τον πόλεμο. Γιατί; Γιατί είναι και αυτοί άνθρωποι Πόλο, σας και μας, είτε το πιστεύεις, είτε όχι»

«Και γιατί πιστεύεις πως χρειάζεται να χυθεί παραπάνω αίμα απ’ ό,τι έχει ήδη χυθεί;»

«Δεν ξέρω… λογικά αν υπήρχε ενδιάμεση λύση μεταξύ των δύο πλευρών, πιστεύω οι προηγούμενοι ηγέτες θα την είχαν βρει»

«Είναι κάπως ειρωνικό, δεν νομίζεις;» άλλαξα θέση στο κρεβάτι του ώστε να βλέπω καλύτερα την πόλη από το αριστερό παράθυρο του δωματίου του.

«Τι ακριβώς;»

«Σκεφτόμαστε όλη μας τη ζωή για το μέλλον μας, πως θα εξελιχθεί ο κόσμος γύρω μας, αλλά σχεδόν ποτέ δεν πάνε τα “σχέδια” μας σωστά. Κοιτώντας πίσω όμως, μπορούμε να βρούμε το λόγο που υπήρξε η αναταραχή»

«Που θες να καταλήξεις;» με ρώτησε ανυπόμονα ο Μάρκο.

«Είναι ειρωνικό που είμαστε αναγκασμένοι να ζούμε τη ζωή προς τα εμπρός, αλλά αυτή βγάζει νόημα μόνο ανάποδα. Δεν έχουμε καμία δεύτερη ευκαιρία, δεν έχουμε κανέναν τρόπο να αναιρέσουμε τις επιλογές μας, αλλά παρόλα αυτά καλούμαστε να ζήσουμε μέσα στην αγωνία της επόμενης μας κίνησης. Από την άλλη, έχουμε μόνο μία ζωή Μάρκο. Γιατί δεν τρέχουμε προς τα μεγαλύτερα μας όνειρα λες και έχουμε πάρει φωτιά;»

«Είναι εύκολο να το λες απλά» μου απάντησε αυτός. «Αλλά έχεις δίκιο. Ίσως αυτό είναι η διαφορά της θεωρίας και της πράξης» είπε και σηκώθηκε από την καρέκλα του. «Παρόλα αυτά, σε ευχαριστώ για τη βοήθεια σου Πόλο» άρχισε να περπατάει προς την πόρτα του δωματίου του.

«Χαι… χαίρομαι που ήμουν χρήσιμος»

«Θες να έρθεις μήπως να φάμε μαζί βραδινό;»

«Ο… όχι… έχω φάει ήδη. Και μάλλον ο παππούς μου θα με περιμένει για να κοιμηθούμε»

«Δεν πειράζει, και αύριο μέρα είναι»

«Ναι…»

«Α, πριν φύγω, δεν θέλω να το ξεχάσω. Είχες ξεχάσει στο κοστούμι που σου είχα δανείσει το σακούλι με τα νομίσματα σου. Το έχω στη δεξιά μεριά του γραφείου. Πάρτο πριν φύγεις, αν κάτσεις στο δωμάτιο μου και γράψεις ή κάτι»

«Α… ευχαριστώ» του είπα και τον κοίταξα καθώς άνοιγε την πόρτα.

«Α, και αν δεν στο έχω πει μέχρι σήμερα Πόλο…» έκανε μία παύση πριν αποχωρήσει από το δωμάτιο «…είσαι καλός φίλος» είπε και έκλεισε την πόρτα πίσω του σιγανά, πριν προλάβω να του απαντήσω κάτι.

Έμεινα άναυδος. Όσες σκέψεις με προβλημάτιζαν εξαφανίστηκαν τη στιγμή που ο Μάρκο ξεστόμησε της λέξεις αυτές. Φαινόταν από το ύφος του πως το εννοούσε με όλη του την καρδιά, αλλά εγώ δεν μπορούσα να κάνω τίποτα πλέον. Κατάπληκτος όπως ήμουν, σηκώθηκα από το κρεβάτι του Μάρκο, στρώνοντας λίγο την μεριά στην οποία είχα καθίσει για να μην έχει ζάρες το πάπλωμα του, και κατευθύνθηκα προς το δωμάτιο μου, όπου είδα τον γέρο να είναι ξαπλωμένος στην αριστερή μεριάτου κρεβατιού, δίνοντας μου τη δυνατότητα να καθίσω δεξιά, όπου βρισκόταν το παράθυρο και να απολαύσω τη θέα για μία τελευταία φορά.

Κάθισα για αρκετή ώρα ήσυχος κοιτώντας την πόλη και τα διάφορα σπίτια της που κατα πάσα πιθανότητα δεν θα ξαναέβλεπα ποτέ. Τα φώτα της πόλης ήταν μερικά αναμένα, μερικά κλειστά, και ένα συγκεκριμένο μπροστά στο σπίτι του άντρα που είχε αποβιώσει τις προηγούμενες μέρες. Σημάδι; Μπορεί. Προκανονισμένο; Όχι, σίγουρα όχι. Η αποστολή μας θα τελείωνε σύντομα, όμως δεν μπορούσα να δω το φως στο τέλος του τούνελ, τα τριαντάφυλλα στην άκρη του θάμνου με τα αγκάθια στον οποίο είχα μπλεχτεί. Το τέλος ήταν τόσο κοντά, αλλά παράλληλα πιο μακριά από ποτέ. Ο χρόνος κυλούσε και το μόνο που περνούσε από το μυαλό μου ήταν η συνέχεια του παραμυθιού, του ψέματος που είχαμε χτίσει τις τελευταίες μέρες. Και γιατί; Για να κερδίσουμε τον πόλεμο; Ο στόχος αυτός μου φαίνεται ασήμαντος πλέον. Οι άνθρωποι αλλάζουν, το πιστεύω με όλη μου την καρδιά. Γιατί πρέπει ο θάνατος να είναι η πρώτη λύση; Μακάρι να μην είχα διαλεχτεί ποτέ για αυτή την αποστολή. Μακάρι…

«Μάρκο» άκουσα τη φωνή του Γιόχαν από πίσω μου να με καλεί.

«Ναι;» 

«Νομίζω ήρθε η ώρα…» είπε και άνοιξε το συρτάρι του.


\chapter{}

Πρώτη φορά είχε μαζευτεί τόσος πολύς κόσμος στο στρατόπεδο το οποίο βρισκόμουν. Στοιχισμένοι όλοι σε μία σειρά, περιμέναμε τον στρατηγό να προσέλθει για να μας ανακοινώσει το νέο του πλάνο. Είναι από τις λίγες φορές που είχα δει τον στρατηγό με τα μάτια μου, καθώς αυτός βρισκόταν συνήθως στις πρώτες γραμμές του πολέμου, ενώ εγώ ήμουν τοποθετημένος δύο στρατόπεδα πίσω. 

Κοιτώντας γύρω μου, έβλεπα μόνο στρατιώτες, δεκανέες και λοχίες, όλοι τους από είκοσι μέχρι τριάντα ετών, μαζί τους και εγώ. Κανένας υψηλόβαθμος στο οπτικό μας πεδίο, όμως όλοι μας περιμέναμε ακίνητοι τον στρατηγό. Ο μόνος θόρυβος που ακουγόταν στην αυλή που βρισκόμασταν ήταν από το τρίψιμο της στολής στα σώματα των στρατιωτών που ήμασταν περίεργοι και κοιτούσαμε δεξιά και αριστερά τους υπόλοιπους συναδέλφους μας.

Ο στρατηγός δεν άργησε πολύ να φτάσει στο σημείο που βρισκόμασταν. Εμφανίστηκε πίσω από μία σκηνή από τα δεξιά μας ακολουθούμενος από είκοσι περίπου συνταγματάρχες, αντισυνταγματάρχες, ταγματάρχες και λοχαγούς, στοιχισμένοι σε δύο σειρές ανάλογα με την ηλικία τους. Οι μισοί από αυτούς φαίνονταν μεγαλύτεροι ηλικιακά, σίγουρα άνω των εξήντα ετών, ενώ οι υπόλοιποι ήταν κατά μέσο όρο μεσήλικες, γνωστοί στο στρατό για τις ικανότητες τους με οποιοδήποτε όπλο στο χέρι.

«Λοιπόν, δεν θα μακρυγορήσω» είπε ο στρατηγός και τίναξε το χαρτί που κρατούσε. «Όσοι ακούτε το όνομα σας κάνετε ένα βήμα μπροστά. Βόλφγκανγκ Γκρίμμερ»

«Μάλιστα» είπε ένας στρατιώτης που βρισκόταν στα αριστερά μου και έκανε ένα βήμα μπροστά.

«Χμ…» αναφώνησε σκεπτόμενος ο στρατηγός. «Όχι. Παρακαλώ επιστρέψτε στη θέση σας. Ντίτερ Μέσνερ»

«Μάλιστα» φώναξε ο στρατιώτης που καθόταν ακριβώς δεξιά μου.

«Λοιπόν…» σημείωσε κάτι στο χαρτί του ο στρατηγός «…ακολούθα τους δύο αυτούς λοχαγούς» έδειξε τους πρώτους δύο άντρες που βρίσκονταν πίσω του.

«Μάλιστα» του απάντησε ο στρατιώτης και ακολούθησε τους δύο άλλους άντρες που τον συνόδευσαν πέρα από τις σκηνές που είχε πίσω του ο στρατηγός.

«Επειδή δεν θέλω να επαναλαμβάνομαι, θα σας λέω απλά ναι ή όχι και εσείς πιστεύω είστε αρκετά έξυπνοι ώστε να δράσετε κατάλληλα. Κατανοητό;»

«Μάλιστα» ακούστηκε ομόφωνα η απάντηση μας.

«Ωραία λοιπόν. Μάρκο… Βάινμπαχ»

«Μάλιστα» ακούστηκε ένας στρατιώτης κοντά στην δεξιά άκρη της σειράς.

«Όχι» του απάντησε ο στρατηγός.

Η κατάσταση αυτή συνέχισε για αρκετά λεπτά, με τον στρατηγό να απορρίπτει πάρα πολλούς στρατιώτες στην αρχή. Δεν φώναζε και τα ονόματα αλφαβητικά, οπότε δεν ήμουν σίγουρος πότε θα ακουγόταν το δικό μου. Μετά από τις τρεις ή τέσσερις απόδοχές του, ο στρατηγός άρχισε να είναι πιο απλόχερος με τα “ναι” του και οι υψηλόβαθμοι από πίσω του άρχισαν να μειώνονται πιο γρήγορα.

«Μάρκο Μπελνούρι» φώναξε το όνομα μου ο στρατηγός μετά από δύο συνεχόμενες απορρίψεις και με τέσσερις άντρες να έχουν μείνει πίσω του.

«Μάλιστα» φώναξα και προχώρησα μπροστά από τους υπόλοιπους.

«Ωραία» μου απάντησε αυτός και μου έκανε νεύμα με το κεφάλι του να ακολουθήσω τους δύο επόμενους άντρες.

«Από δω νεαρέ» μου απευθύνθηκε ο γηραιότερος άντρας για να τους πλησιάσω και να τους ακολουθήσω.

Οι δύο άντρες με οδήγησαν σε μία άλλη αυλή όπου όλες οι τριάδες που είχαν αποχωρήσει πριν από εμάς βρίσκοντας σε άλλη μία ουρά και περίμεναν, με την πρώτη τριάδα να στέκεται ακριβώς έξω από την σκηνή του στρατηγού. Επικρατούσε ένας μικρός ψίθυρος στο σημείο τούτο, καθώς όλοι οι νεότεροι ρωτούσαν τους ανώτερό τους ως προς τι όλη αυτή η “δοκιμασία”, αν μπορεί να την πει κάποιος έτσι. Το μόνο που γνωρίζαμε εμείς ήταν πως ο στρατηγός είχε καλέσει όλους τους στρατιώτες κάτω τον τριάντα δύο ετών στο στρατόπεδο μου για την ανακοίνωση του νέου πλάνου. Κανείς μας δεν περίμενε αυτό που προηγήθηκε.

«Μήπως έχει γίνει κάποιο λάθος;» ρώτησα τους δύο άντρες που με συνόδευσαν μέχρις εδώ.

«Γιατί το λες;» με ρώτησε πίσω ο γηραιότερος, που είδα από το διακριτό που φορούσε πως ήταν συνταγματάρχης.

«Γιατί ο στρατηγός επιλέγει άτομα; Δεν θα γινόταν ανακοίνωση του νέου πλάνου;»

«Αυτό γίνεται…» με πλησίασε ο συνταγματάρχης «…απλά είναι άκρως μυστικό» μου ψηθύρισε στο αυτί.

«Τι του λες του παιδιού» διαμαρτηρήθηκε ο άλλος άντρας που βρισκόταν στην τριάδα μας.

«Γιατί, δεν είναι μυστικό»

«Όχι, απλώς θα το γνωρίζουμε μόνο εμείς»

«Αν αυτό δεν είναι ο ορισμός του μυστικού, δεν ξέρω τι είναι μυστικό» του είπε ο συνταγματάρχης. «Μην ανησυχείς, όμως, μικρέ. Ίσα ίσα, αν ολοκληρωθεί το “Σχέδιο” με επιτυχία μπορεί να πάρεις και προαγωγή»

«Το “Σχέδιο”;» παραξενεύτηκα

«Ναι, έτσι το λέει ο στρατηγός» μου απάντησε ο άλλος άντρας που από το διακριτό του βαθμού του φαινόταν να είναι αντισυνταγματάρχης. 

«Και γιατί το “Σχέδιο” χρειάζεται τριάδες από άντρες διαφορετικών βαθμών και ηλικίας; Γιατί ο στρατηγός απέρριπτε στρατιώτες δεξιά κι αριστερά; Περάσατε και εσείς κάποια διαδικασία επιλογής;»

«Ηρέμησε μικρέ. Όλα θα βγάλουν νόημα όταν μας εξηγήσει τις λεπτομέρειες ο στρατηγός» με καθησύχασε ο συνταγματάρχης.

«Άρα ούτε εσείς ξέρετε για το “Σχέδιο”;»

«Δεν γνωρίζουμε τίποτα παραπάνω από το γεγονός ότι θα χρειαστεί να δουλεύουμε σε τριάδες αντρών διαφορετικών ηλικιών» μου απάντησε ο αντισυνταγματάρχης.

«Ωραία λοιπόν» ακούστηκε η φωνή του στρατηγού από πίσω μου ακολουθούμενη από ένα παλαμάκι. Ακολουθόταν από την τελευταία τριάδα αντρών που είχε επιλέξει. «Θα εισέρχεστε μέσα τρεις-τρεις, με τη σειρά που είστε στοιχισμένοι και θα σας εξηγώ τους ρόλους σας για τις επόμενες… δεκαεπτά μέρες, ή τουλάχιστον έτσι υπολογίζω» μας είπε ο στρατηγός χαρούμενος και μπήκε στη σκηνή του ακολουθούμενος από την πρώτη τριάδα οι οποίοι αρχικά δίστασαν να τον ακολουθήσουν. 

«Μάλιστα» είπε ο αντισυνταγματάρχης και σταύρωσε τα χέρια του.

«Μάλλον θα χρειαστεί να περιμένουμε αρκετή ώρα, ε;» σχολίασε ο συνταγματάρχης.

«Αφού είμαστε προτελευταίοι. Λογικό είναι»

«Ναι…»

«Συγγνώμη» τους διέκοψα. «Μήπως θα μπορούσα να έχω τα ονόματα σας;»\newline\newline

«Γιόχαν Βάγκνερ, Ράινερ Γουλφ και Μάρκο Μπελνούρι» μας καλωσόρισε ο στρατηγός. «Ελάτε γύρω από το γραφείο μου» μας έκανε νεύμα να πλησιάσουμε.

Πάνω στο γραφείο του είχε τοποθετημένο έναν χάρτη με έντεκα πινέζες στις άκρες πόλεων που κατοικούσαν μεγαλοαστοί και αριστοκράτες, καθώς και έντεκα κόκκινα χ σε κοντινές πόλεις που είχαν πινέζες. Το στρατόπεδο που βρισκόμασταν ήταν πολύ μακριά από τις πινέζες αυτές προφανώς, αλλά ήταν και αυτό σημαδεμένο με έναν μαύρο κύκλο με τις λέξεις “είμαστε εδώ”, το οποίο μου φάνηκε λίγο ειρωνικό. Ο στρατηγός έβγαλε από το συρτάρι του τέσσερα χαρτιά, τρία εκ των οποίων μας τα έδωσε. Ήταν λίστες από ονόματα.

«Αυτοί είναι οι στόχοι σας» μας εξήγησε αφότου μας άφησε λίγο χρόνο να τις επεξεργαστούμε. «Όχι όλοι δικοί σας. Ο δικός σας στόχος είναι ο προτελευταίος, ονόματι… Μάρκο Μπερνούλι. Χα, ειρωνικό, ε;»

«Μάλιστα» είπε ο αντισυνταγματάρχης. «Έχουμε λεπτομέρειες στο που θα τον βρούμε ή…»

«Όλα είναι μπροστά σου Ράινερ» του απάντησε ο στρατηγός.

Ο αντισυνταγματάρχης παραξενεύτηκε.

«Πλάκα κάνω, προφανώς και θα σας εξηγήσω τώρα» είπε ο στρατηγός και ακούμπησε τα χέρια του πάνω στον χάρτη. «Τα ονόματα που έχετε μπροστά σας ανήκουν στους αριστοκράτες που έχουν επιβιώσει μέχρι στιγμής και δεν βρίσκονται στις πρώτες γραμμές του πολέμου, αλλά πίσω στα αρχοντικά τους, στις πόλεις τους. Κάποια αρχοντικά έχουν επιζώντες πατεράδες και γιους, ενώ σε άλλα, σαν το δικό σας, κατοικεί μόνο ο γιος της οικογένειας, είτε επειδή ο πατέρας είναι υψηλόβαθμος του στρατού στις πρώτες γραμμές, είτε επειδή είναι νεκρός. Μη σας νοιάζει όμως αυτό, εσείς θα πρέπει να επικεντρωθείτε στο αρχοντικό αυτό τις επόμενες δυόμιση εβδομάδες» είπε και μας έδειξε μία από τις έντεκα πινέζες που ήταν τοποθετημένες στον χάρτη, η οποία βρισκόταν αρκετά κοντά σε μία τεράστια λίμνη. «Σε αυτή την πόλη, στα βόρεια προάστια, το πρώτο σπίτι που θα συναντήσετε είναι το αρχοντικό της οικογένειας Μπερνούλι»

«Συγγνώμη που αποσπώμαι από τον στόχο…» διέκοψε τον στρατηγό ο αντισυνταγματάρχης «…αλλά οι υπόλοιπες δέκα πινέζες τι συμβολίζουν; Τους υπόλοιπους δέκα στόχους;»

«Ακριβώς Ράινερ» του απάντησε ο στρατηγός.

«Και γιατί κάθε πόλη με πινέζα έχει μία διπλανή πόλη που είναι κυκλομένη;»

«Μη βιάζεσαι, όλα θα έρθουν στην ώρα τους. Επειδή, όμως, σε τρώει η περιέργεια, θα σας εξηγήσω λίγο περιληπτικά όλο το “Σχέδιο” και μετά θα μπω σε λεπτομέρεια. Θα στείλουμε έντεκα τριάδες στα έντεκα αυτά αρχοντικά που βλέπετε σημαδεμένα. Η κάθε τριάδα θα έχει διαφορετικό πλήθος από αντιπάλους να αντιμετωπίσει, οπότε στο αρχοντικό τελικά θα φτάσουν δύο ή τρία άτομα, ανάλογα την περίσταση. Στη δική σας περίπτωση θα είστε δύο»

«Και θέλετε, λοιπόν, να μπουκάρουμε αιφνιδιαστικά και να εξουδετερώσουμε τους αριστοκράτες;»

«Όχι, Ράινερ, αυτό δεν θα μας βοηθούσε και πολύ από μόνο του. Θέλω να μείνετε μαζί τους στο αρχοντικό μέχρις ότου να ηχήσουν τα δύο σημαντικά σινιάλα, συν άλλο ένα αν χρειαστεί» είπε ο στρατηγός και του έδωσε το τέταρτο χαρτί που κρατούσε στο χέρι του.

«Τι να το κάνω αυτό; Έχει μόνο ένα όνομα» παραπονέθηκε ο Ράινερ κοιτώντας μπρος και πίσω το χαρτί που του δώθηκε.

«Το πρώτο σινιάλο…» συνέχισε ο στρατηγός αγνοώντας τον αντισυνταγματάρχη «…θα είναι λογικά κάποιος αγγελιοφόρος που θα έρθει στο αρχοντικό σας ενημερώνοντας την οικογένεια πως οι πρώτες γραμμές του πολέμου έπεσαν»

«Ψεύτικος αγγελιοφόρος;» ρώτησε ο συνταγματάρχης.

«Όχι, κανονικός. Γενικός μας σκοπός είναι να χρησιμοποιήσουμε ως πλεονέκτημα μας την υπεραριθμία μας έναντι των αριστοκρατών και το τραγικό τους λάθος να μαζεύουν όλο τους τον στρατό στο μέτωπο του πολέμου. Αφού κάνουμε την επίθεση μας στις πρώτες γραμμές του πολέμου, με απειροελάχιστη πιθανότητα να ηττηθούμε, θα αφήσουμε ελάχιστους αντίπαλους στρατιώτες, είτε επίτηδες, είτε κατά λάθος, δεν ξέρω, να ξεφύγουν ώστε να μπορέσουν να ενημερώσουν τον υπόλοιπο κόσμο της μεριάς τους. Αυτό θα είναι το πρώτο σινιάλο σας.

«Ενδιαφέρον» είπε ο αντισυνταγματάρχης με το δεξί του χέρι να τρίβει το πιγούνι του.

«Δύο μέρες αφότου ακουστεί το πρώτο σινιάλο, οι δύο που θα βρίσκονται μέσα στο αρχοντικό θα κάνετε το δεύτερο σινιάλο, το οποίο…»

«Συγγνώμη…» τον διέκοψα διστακτικά «…αλλά… τι εννοείται θα είμαστε μέσα στο αρχοντικό;»

«Α δεν το εξήγησα; Με συγχωρείτε, είστε από τους τελευταίους και θεωρώ κάποια πράγματα αυτονόητα. Όταν λέω θα είστε “μέσα” στο αρχοντικό, εννοώ πως θα μένετε μαζί με τους αριστοκράτες»

«Πως;» παραξενεύτηκα.

«Βλέπετε αυτές τις πόλεις που έχουμε κυκλώσει; Συγκεκριμένα αυτή εδώ» είπε ο στρατηγός και μου έδειξε την πιο κοντινή πόλη στο αρχοντικό-στόχο μας, η οποία βρισκόταν βορειοανατολικά του στόχου μας. «Αυτή η πόλη δεν θα υπάρχει σε λίγο καιρό»

«Τι εννοείται;» τον ρώτησα.

«Όλες αυτές οι πόλεις θα είναι θύμα επίθεσης από μεριάς μας, πριν εσείς προλάβετε να φτάσετε στα αρχοντικά σας. Εσύ και ο Γιόχαν θα το χρησιμοποιήσετε αυτό ως δικαιολογία για να μπείτε στο αρχοντικό ως…»

«Πώς;»

«Θα σας βάλουν, μην ανησυχείς νεαρέ. Το μόνο που χρειάζεται να κάνετε είναι να τραυματίσετε τον Γιόχαν σε ένα μη θανάσιμο σημείο μία μέρα πριν, ώστε να μη ρέει αίμα από την πληγή του και οι αριστοκράτες θα σας βοηθήσουν»

«Γιατί να μας βοηθήσουν;»

«Για να καλυτερεύσουν την εικόνα τους. Πιθανώς θα φοβούνται πως ο υπόλοιπος κόσμος των μεγαλοαστών ασκεί κρυφά κριτική προς αυτούς επειδή δεν συμμετέχουν ενεργά στον πόλεμο. Φιλοξενώντας δύο “θύματα πολέμου” θα τους βοηθήσει να καλυτερεύσουν την εικόνα που νομίζουν ότι έχουν στον υπόλοιπο λαό τους»

«Και πως είμαστε σίγουροι πως αυτό θα γίνει;» τον ρώτησε ο συνταγματάρχης.

«Δεν είμαστε, απλά ποντάρουμε με βάση τις πιθανότητες και την ψυχολογία τους. Αν δεν σας δεχτούν, που δε νομίζω, γιατί μετά χαλάει σχεδόν όλο το “Σχέδιο”, θα μπορέσετε να μείνετε με τον Ράινερ μέρες. Όταν, τέλος πάντων, λάβουν χώρα και τα δύο σινιάλα, θα χρειαστεί ο Μάρκο και ο Ράινερ να αναλάβουν την νύχτα μετά το τελευταίο σινιάλο την βρώμικη δουλειά. Αυτό θα είναι και το τέλος της αποστολής σας. Τώρα… γιατί το κάνουμε όλο αυτό. Πιστεύουμε πως μετά την υποχώρηση των αριστοκρατών, θα τους πάρει το πολύ δύο εβδομάδες να ξαναφτιάξουν στρατό έτοιμο για αντεπίθεση, τους έχουμε ξαναδεί στο παρελθόν να το ξανακάνουν. Όμως, η πηγή ηγετών τους είναι τα αρχοντικά και ειδικότερα οι νέοι…» με κοίταξε ο στρατηγός κάνοντας μία παύση. «Χτυπώντας την πηγή πριν αυτοί προλάβουν να αναδιαταχτούν, θα τους φέρουμε σε ευάλωτη θέση, την οποία προφανώς θα εκμεταλλευτούμε για να κερδίσουμε τον πόλεμο» έκανε άλλη μία παύση ο στρατηγός τελειώνοντας τον τελευταίο μονόλογο του για την ώρα. «Τώρα για τις λεπτομέρειες…»

«Επιτέλους» αναφώνησε ο Ράινερ.

«Οι τέσσερις ομάδες που θα πάνε σε μακρινά αρχοντικά θα μεταφερθούν μέχρις ένα σημείο με αυτοκίνητα, αλλά δεν σας ενδιαφέρει αυτό. Τα δύο τρίτα περίπου των ομάδων θα ταξιδέψετε μαζί…» ξεκίνησε ο στρατηγός να λέει κυλώντας σιγά σιγά το δάχτυλο του από το στρατόπεδο μας προς τις πινέζες «…κάνοντας πεζοπορία μέχρις ότου να φτάσετε σε αυτή τη διασταύρωση. Τότε είναι που εσείς θα χωριστείτε από αυτούς και θα κόψετε δρόμο προς το αρχοντικό» μας έδειξε χρησιμοποιώντας τους δείκτες και των δύο χεριών του για να καταλάβουμε την αλλαγή πορείας. «Συνολικά το ταξίδι αυτό θα πρέπει να σας πάρει επτά μέρες. Μην ανησυχείτε, θα σας δοθούν προμήθειες για δέκα, οπότε το φαγητό δεν θα είναι πρόβλημα»

«Και πως θα τα κουβαλήσουμε όλες αυτές τις προμήθειες;» τον ρώτησε ο αντισυνταγματάρχης.

«Παραδόξως είναι πιο ελαφριές απ’ ό,τι νομίζεις» του απάντησε ο στρατηγός. «Φτάνοντας στο αρχοντικό, λοιπόν, θα πρέπει, όπως είπαμε να έχετε τραυματίσει τον Γιόχαν, ώστε το επιχείρημα “Ω, είμαστε θύματα πολέμου” να είναι πιο πειστικό. Πίσω από το αρχοντικό υπάρχει δάσος, καθώς και σε όλη την διαδρομή σας θα ακολουθήσετε το δάσος. Δεν σας το είπα αυτό, και ήταν σημαντικό. Επειδή οι αριστοκράτες έχουν χτίσει δρόμους για να μπορούν να μετακινούνται και αυτοί με αυτοκίνητα μεταξύ περιοχών. Είναι επικίνδυνο, όμως, εσείς να περπατάτε κοντά σε αυτούς τους δρόμους, καθώς αν σας δει κάποιος αριστοκράτης πριν φτάσετε στο αρχοντικό, το χάσαμε το παιχνίδι. Κατανοητό;»

«Μάλιστα» είπαμε και οι τρεις λίγο ετεροχρονισμένα.

«Που είχα μείνει; Α, φτάνοντας στο αρχοντικό, που έχει πίσω δάσος, μπλα, μπλα, Γιόχαν και Μάρκο θα αφήσετε κανένα χιλιόμετρο πίσω τα όπλα σας και ό,τι προμήθεια σας έχει απομείνει και θα κατευθυνθείτε προς το αρχοντικό οι δυο σας κουβαλώντας μόνο ένα σακούλι που θα έχει μέσα μερικά νομίσματα, για πειθώ, και μερικά χάπια που θα σας χρησιμεύσουν αργότερα. Ο Ράινερ θα σας αποχεραιτήσει και θα κατευθυνθεί με το όπλο και τις προμήθειες του στο σπίτι του κυρίου που αναγράφεται στο χαρτί που του έδωσα λίγο πιο πριν» ο Ράινερ κοίταξε πάλι πεταχτά το τέταρτο χαρτί. «Αυτός ο άντρας μάθαμε πως ζει μόνος του και το σπίτι του είναι κοντά στο αρχοντικό, οπότε θα είναι καλή κρυψώνα για να τοποθετηθεί ο Ράινερ»

«Και πως δεν θα μάθει πως ένας ακόμη άντρας μένει σπίτι του;» ρώτησε ο Ράινερ.

«Μα θα το ξέρει. Στόχος είναι να τον απειλήσεις κάπως, ξέρεις εσύ από τέτοια, και να τον “πείσεις” να σε φιλοξενήσει για τις επόμενες μέρες»

«Και αν δεν μπορέσει να κρατήσει το στόμα του κλειστό;»

«Ε τότε θα χρειαστείς και συ ένα σακούλι με φάρμακα. Αν πιστεύεις ότι η συμπεριφορά του ξεφεύγει, ανάγκασε τον να πάρει υπερβολική δόση φαρμάκων και να αυτοκτονήσει. Πριν το κάνει όμως, βεβαιώσου, με όποιον τρόπο, ότι θα γράψει μία διαθήκη, στην οποία θα κληρωνομείς εσύ το σπίτι του»

«Μάλιστα… ενδιαφέρον.»

«Αυτό είναι το σινιάλο “μηδέν”. Αν δεν χρειαστεί, δεν θα λάβει χώρα, αλλά αν συμβεί, οι άλλοι δύο να είστε σίγουροι πως ο τρίτος συνεργάτης σας ζει και είναι στην πόλη»

«Ω, θα χρειαστεί» σχολίασε ο αντισυνταγματάρχης. «Δεν είμαι και πολύ χαλαρός στις απειλές μου»

«Δεν χρειάζεται να μπούμε σε παραπάνω λεπτομέρεια για αυτό το κομμάτι, νομίζω πιο σημαντικό είναι…»

«Συγγνώμη, αλλά μπερδεύτηκα» διέκοψα τον στρατηγό. «Άρα υπάρχουν τρία σινιάλα;»

«Ναι αλλά τα δύο άλλα είναι τα σημαντικά, και μπορεί το τρίτο που σας είπα τώρα να έρθει πριν ή ενδιάμεσα των δύο άλλων, για αυτό δεν το εντάσσω μαζί με τα άλλα. Αν σε βολεύει να το σκέφτεσαι καθολικά, πες ότι έχουμε τρία σινιάλα»

«Θα έχουμε τρία σινιάλα» επιβεβαίωσε ο αντισυνταγματάρχης.

«Καλά, ό,τι πιστεύετε είναι απαραίτητο. Πάμε τώρα σε εσάς τους δύο. Στο αρχοντικό θα χρειαστεί να περάσετε, αν έχουμε υπολογίσει καλά, τουλάχιστον άλλες επτά ημέρες μαζί με την οικογένεια των αριστοκρατών, οι οποίοι λογικά θα είναι από τρεις μέχρι πέντε, ανάλογα αν έχουν έναν οι τρεις “μπάτλερ”, ή όπως τους φωνάζουν τέλος πάντων. Αυτές τις μέρες, Μάρκο, σε θέλω να δεθείς με τον στόχο μας, τον γιο των Μπερνούλι»

«Τι; Γιατί;»

«Είναι πολύ πιθανό πως η οικογένεια κάποια στιγμή θα θελήσει να σας διώξει κάπως έμμεσα από το αρχοντικό. Επειδή δεν πρέπει να ρισκάρουμε την πρόωρη αποχώρηση σας από το αρχοντικό, προσπάθησε να συζητάς με τον συνομήλικο σου. Είπα ότι είναι συνομήλικος σου έτσι; Και ότι σε διάλεξα επειδή μοιάζεται λίγο έτσι;»

«Όχι…»

«Δεν πειράζει, το είπα τώρα. Γίνε φίλος με τον συνομήλικο σου. Έτσι, δύσκολα θα σας διώξουν από το σπίτι τους»

«Ναι αλλά…» έκανα μία μικρή παύση για να προσπαθήσω να σκεφτώ αν οι επόμενες μου λέξεις έβγαζαν νόημα μόνο στο κεφάλι μου ή αν είχα όντως βρει ένα λάθος στο “Σχέδιο” του στρατηγού. «Είπατε πριν πως εγώ και ο Ράινερ θα αναλάβουμε την βρώμικη δουλειά, που, θέλω να ξεκαθαρίσω, είμαι απολύτως εντάξει με αυτό. Αν, όμως, δεθώ πολύ με τον στόχο μας… τότε τι;»

«Δεν πρόκειται να γίνει αυτό, ούτε σε εκατό χρόνια. Παρά τις εμφανησιακές ομοιότηες σας, σε διαβεβαιώνω πως δεν θα έχετε κάτι με το οποίο μπορείτε να δεθείτε»

«Και στη χειρότερη…» πρόσθεσε ο αντισυνταγματάρχης «…θα αναλάβω εγώ» μου είπε κλείνοντας το μάτι του.

«Εντάξει…» τους απάντησα, χωρίς να έχω πειστεί εξ ολοκλήρου.

«Ωραία, λοιπόν. Επτά μέρες. Γίνεστε φίλοι. Την όγδοη μέρα, θα έρθει το πρώτο ή δεύτερο σινιάλο σας, ανάλογα τι θα κάνει ο Ράινερ. Υπολογίζουμε πως οι αγγελιοφόροι θα πάρουν μιάμιση μέρα να ενημερώσουν και τα έντεκα αρχοντικά, οπότε, επειδή εσείς θα βρίσκεστε και στο δεύτερο ή τρίτο, αν υπολογίζω καλά, αρχοντικό από τις πρώτες γραμμές του πολέμου, θα ενημερωθείτε για την πτώση της πρώτης γραμμής την πρώτη μέρα. Το βράδυ της ένατης ημέρας θα ξεκινήσετε το τρίτο σινιάλο»

«Θα ξεκινήσουμε;» παραξενεύτηκα για πολλοστή φορά.

«Ναι, θα ξεκινήσετε, γιατί το σινιάλο ουσιαστικά θα είναι ορατό το πρωί της δέκατης μέρας. Θυμάστε αυτό το σακούλι που σας είπα; Με τα φάρμακα;»

«Ναι»

«Ε, αυτό, λοιπόν θα είναι το σινιάλο σας»

«Τι;»

«Όπως και το θύμα του Ράινερ, μάλλον, τις προηγούμενες μέρες…» είπε ο στρατηγός και έριξε μια πλάγια ματιά στον αντισυνταγματάρχη «…έτσι και ο Γιόχαν θα καταναλώσει υπερβολική δόση φαρμάκων ώστε να αυτοκτονήσει, κάνοντας το να φαίνεται λες και είναι φυσικός θάνατος»

«Τι;» ξαφνιάστηκα. «Συμφωνήσατε εσείς σε κάτι τέτοιο» ρώτησα τον συνταγματάρχη.

Αυτός δεν μου απάντησε. Δεν με κοίταξε καν στα μάτια. Απλώς μου χαμογέλασε.

«Το σινιάλο αυτό είναι απαραίτητο, καθώς οι αριστοκράτες κατά πάσα πιθανότητα θα θελήσουν να κρατήσουν κρυφή την ήττα τους από το πλήθος για να μην προκαλέσουν άμεσο πανικό. Με τον θάνατο του Γιόχαν, όμως, η οικογένεια του αρχοντικού μπορεί να φωνάξει γιατρούς το επόμενο πρωί, ή μπορείς να φωνάξεις εσύ γιατρούς, ό,τι τύχει εκείνη τη στιγμή. Ο Ράινερ προφανώς θα δει τον πανικό που επικρατεί στην πόλη με τους γιατρούς να τρέχουν να προλάβουν να σώσουν έναν ήδη νεκρό άνθρωπο και θα καταλάβει πως ήρθε η ώρα να δράσετε. Τα μεσάνυχτα της δέκατης μέρας, θα ανοίξεις την είσοδο στον Ράινερ, που, γνωρίζοντας προφανώς που είχατε αφήσει τα όπλα σας, θα σου φέρει το δικό σου όπλο σαν ταχυδρόμος και, αφού είστε μέσα πλέον στο αρχοντικό, η τελευταία φάση της αποστολής σας είναι προφανής»

Επικράτησε μία παύση στη σκηνή μετά τις τελευταίες λέξεις του στρατηγού.

«Ανακεφαλαιώνοντας: ταξίδι επτά ημέρες. Πρώτη φάση είναι να μπείτε στο αρχοντικό. Δεύτερη φάση είναι να περιμένετε για τα σινιάλα, χτίζοντας δεσμό με τον στόχο σας. Τελευταία φάση, σινιάλο του Γιόχαν και εξουδετέρωση της οικογένειας. Α, και το βράδυ αφότου ολοκληρώσετε την τελευταία φάση θα υπάρχει όχημα διαφυγής κοντά στο αρχοντικό με το οποίο θα μπορέσετε να φύγετε. Κατανοητό;»

«Μάλιστα» είπαμε όλοι μαζί ετεροχρονισμένα πάλι.

«Μία τελευταία λεπτομέρεια που θέλω να συζητήσετε στο ταξίδι σας προς το αρχοντικό είναι οι λεπτομέρειες όπως το που θα τραυματιστεί ο Γιόχαν, για παράδειγμα, που θα κρύψετε τα φάρμακα στο αρχοντικό ή η χρήση ψεύτικων ονομάτων. Το τελευταίο είναι σημαντικό βασικά, μη τολμήσετε και χρησιμοποιήσετε τα κανονικά σας ονόματα. Οι δύο από τους τρεις έχετε βόρεια ονόματα, οπότε πρέπει οι δυο σας, αν όχι όλοι σας, να βρείτε ψεύτικα ονόματα, ώστε να μην κινήσετε υποψίες. Αν σκεφτείτε κι άλλες μικρές προσθήκες σαν αυτές που σας είπα προηγουμένως, προφανώς συζητήστε και αυτοσχεδιάστε. Ο σκελετός του “Σχεδίου”, όμως, παρακαλώ να ακολουθηθεί κατα γράμμα. Κατανοητό;»

«Μάλιστα» είπαμε όλοι μαζί ταυτόχρονα αυτή τη φορά.

«Ωραία, λοιπόν. Υπάρχει καμία τελευταία ερώτηση;»

«Ίσως λίγο λεπτομέρεια…» πήρα πρωτοβουλία να ρωτήσω τον στρατηγό. «...αλλά γιατί θα χρειαστεί να βρισκόμαστε δέκα ημέρες στο αρχοντικό; Γιατί δεν κάνετε επίθεση στο μέτωπο νωρίτερα;»

«Γιατί, φοβάσαι μη δεθείς υπερβολικά με τον συνομήλικο σου;» σχολίασε ο αντισυνταγματάρχης.

«Όχι… απλά θέλω να μάθω ως προς τι η μεγάλη αυτή αναμονή…» του απάντησα.

«Για να σιγουρευτούμε πως όλες οι τριάδες έχουν φτάσει στα αντίστοιχα αρχοντικά» μου είπε ο στρατηγός.

«Α… τόσο απλό;»

«Ναι. Εσείς έχετε αναλάβει από τα πιο βόρεια αρχοντικά, οπότε θα χρειαστεί να περιμένετε περισσότερο από άλλους»

«Μάλιστα» απάντησα στον στρατηγό και μαζεύτηκα στην μεριά μου.

«Κάποια άλλη απορία;»

«Όχι» του είπε ο συνταγματάρχης. «Νομίζω είμαστε καλά»

«Ωραία, λοιπόν. Αυτά από εμένα. Όλες οι ομάδες θα αποχωρήσετε αύριο το πρωί από τα νότια. Καλή ξεκούραση μέχρι τότε»


\chapter{}

Ξόδεψα όλο το πρωινό της επόμενης ημέρας πλάι στους  γιατρούς, υπογράφοντας χαρτιά ύπο το όνομα μιας ψεύτικης περσόνας, χαρτιά σχετικά με το θάνατο ενός ψεύτικου προσώπου που από πίσω του έκρυβε μία πραγματική ψυχή. Όσο περνούσαν οι ώρες μου στο νοσοκομείο, τόσο περισσότερο έβρισκα τον θάνατο του Γιόχαν ανούσιο. Γιατί δεν μπορούσαμε να ενημερώσουμε κάπως αλλιώς τον Ράινερ. Γιατί δεν μπορούσαμε να επιστρέψουμε και οι τρεις σπίτι ήρωες… Ήρωες…

Στους δρόμους της πόλης στον δρόμο μου πίσω στο αρχοντικό παρατηρούσα το πλήθος. Όλοι τους ζούσαν τις ζωές τους όπως κανονικά, κανένα σημάδι υποψίας. Πως να με βρει κανείς όμως ύποπτο; Το “Σχέδιο” είχε φροντίσει για την προσωρινή ένταξη μας στην κοινωνία των αριστοκρατών, ακόμη και στην περίπτωση μας ίσως είχα υπερβάλει λιγάκι. Μερικοί περαστικοί με χαιρετούσαν κιόλας, μάλλον επειδή ήξεραν ότι βρίσκομαι κοντά στον Μάρκο, οπότε λογικά ήθελαν να έχουν καλές σχέσεις μαζί μου, μαζί με την ψεύτικη αυτή περσόνα των τελευταίων ημερών. Δεν θα προλάβαιναν όμως να χτίσουν καμία σχέση με την περσόνα αυτή. Ο Πόλο θα εξαφανιζόταν σύντομα χωρίς λόγο και αιτία, μαζί με την οικογένεια του αρχοντικού.

Κοντά στο σπίτι του άντρα που είχε πεθάνει πέντε μέρες πριν, είδα τις κοπέλες από εκείνη τη μέρα να κάθονται σε ένα παγκάκι και να συζητάνε. Δεν είχα, όμως, το κουράγιο να τους μιλήσω. Δεν θα πάθαιναν τίποτα τις επόμενες εικοσιτέσσερις ώρες, δεν ήταν θύματα του “Σχεδίου”, όμως οι τύψεις μου δεν με άφηναν να προχωρήσω προς το μέρος τους, έστω και για μία καλημέρα, ένα γεια, μία μικρή συζήτηση. Εξάλλου, ο Ράινερ βρισκόταν πάντα μερικά βήματα πίσω μου, σαν τίγρης που παρακολουθούσε το θήραμα του από απόσταση. Αν με έβλεπε να κάνω δημόσιες σχέσης την πιο σημαντική μέρα του “Σχεδίου” δεν θα ήταν και πολύ χαρούμενος το βράδυ μετά την ολοκλήρωση της αποστολής μας. Στάθηκα, παρόλα αυτά, κοντά στην είσοδο του κήπου του αρχοντικού, παρατηρώντας τις κοπέλες, ή τουλάχιστον ό,τι μπορούσα να δω από τόσο μακριά, ρίχνοντας και κλεφτές ματιές στον Ράινερ ο οποίος είχε στηθεί πίσω από μία γωνία κρυμμένος και αυτός από τις δύο κοπέλες, καθώς γνώριζε πως αν τον έβλεπαν θα του έπιαναν πάλι συζήτηση. Η εικόνα περιέγραφε τις τελευταίες δέκα μέρες στην εντέλεια. Ο κυνηγός περιμένει ώσπου το θήραμα του να βρεθεί σε ευάλωτη θέση για να το εξουδετερώσει αιφνιδιαστικά.

Ξεφυσώντας μία τελευταία φορά, έκανα νεύμα με το κεφάλι στον Ράινερ για να επιβεβαιώσω πως είχε λάβει το μήνυμα και αυτός μου απάντησε με ένα νεύμα πίσω, δείχνοντας μου την κατανόηση του. Τα σινιάλα εντωμεταξύ γιατί ήταν απαραίτητα; Αφού γνωρίζαμε από πριν πως την έβδομη μέρα θα πέσει η πρώτη γραμμή των αριστοκρατών. Μπορεί να πήγαινε κάτι λάθος και να αργούσε ξέρω γω η επίθεση ή κάτι. Μπορεί. Πιθανότατα όχι, αλλά ο στρατηγός μάλλον αποφάσισε να μην πάρει αυτό το ρίσκο. Το ρίσκο όμως για να μπούμε στο αρχοντικό; Γιατί έπρεπε να το πάρουμε; Ήταν λιγότερο σημαντικό; Δεν θα μπορέσω να καταλάβω ποτέ πως κάποιος άνθρωπος… ΆΝΘΡΩΠΟΣ σκέφτηκε αυτό το σατανικό “Σχέδιο”.

Μέσα στο αρχοντικό τα πράγματα κύλησαν όπως τις υπόλοιπες μέρες. Πέραν της συμπόνιας του μπάτλερ της οικογένειας, δεν μπόρεσα να μιλήσω σε κανέναν άλλο όλη την υπόλοιπη μέρα. Ο Μάρκο βρισκόταν στο δωμάτιο του, όπως τις άλλες δύο μέρες, κοιτώντας έξω από το παράθυρο του σκεπτόμενος. Όταν άνοιξα την πόρτα του και τον είδα να έχει την πόζα αυτή που είχε και τις προηγούμενες μέρες, αποφάσισα να τον αφήσω στην ησυχία του. Είχαμε και οι δύο πράγματα να σκεφτούμε την υπόλοιπη μέρα έτσι κι αλλιώς, οπότε λογικά δεν θα με ενοχλούσε ούτε αυτός αν καθόμουν στο δωμάτιο μου και τους συνόδευα στην τραπεζαρία μόνο για τα σημαντικά γεύματα. Άλλο ένα λάθος στο “Σχέδιο” όμως. Τι θα γινόταν αν ο Μάρκο έπαιρνε απόφαση να φύγει από το αρχοντικό νωρίτερα από το τρίτο σινιάλο; Δεν το σκέφτηκε ο στρατηγός. Λογικά θα πόνταρε πως ο φόβος θα έκανε τους άεργους αριστοκράτες να μη μπορούν να πάρουν άμεσα απόφαση. Είμαστε αρκετά τυχεροί που δεν έχει πάει κάτι στραβά μέχρι τώρα είναι η αλήθεια.

Μπαίνοντας στο δωμάτιο μου, είδα ακουμπισμένο στο κομοδίνο μου το σακούλι με τα νομίσματα που είχα ξεχάσει να πάρω την προηγούμενη μέρα από το δωμάτιο του Μάρκο. Το πήρα στα χέρια μου και το επεξεργάστηκα λίγο. Δεν ήμασταν τόσο χαζοί ώστε να αφήσουμε τα φάρμακα μέσα στο σακούλι. Από την πρώτη κιόλας στιγμή που βρεθήκαμε στο δωμάτιο μας αυτό, αδειάσαμε τα φάρμακα του Γιόχαν στο πρώτο συρτάρι του. Όταν ο μπάτλερ μας ρώτησε προς τι τα φάρμακα, γιατί τα είχε παρατηρήσει και νωρίτερα λογικά όταν μας εξέταζε για όπλα, του είχαμε εξηγήσει πως είναι για τον αντισυνταγματάρχη, τον τότε “παππού” μου. Η καλή ψυχή που ήταν ο μπάτλερ μας έδινε κάθε βράδυ ένα ποτήρι νερό για να μπορέσει ο “παππούς” μου να καταπιεί το φάρμακο της ημέρας εκείνης πιο εύκολα. Το σακούλι ουσιαστικά ήταν διακοσμητικό και ασήμαντο από εκείνη τη μέρα και μετά. Το έβαζα σε μία τυχαία τσέπη μου μόνο όταν βγαίναμε έξω ώστε να μπορέσω αν χρειαστεί να προσπαθήσω να πληρώσω για κάτι, και λέω “να προσπαθήσω” καθώς με τόσα λίγα νομίσματα δεν νομίζω να μπορούσα να αγοράσω και πολλά πράγματα.

Η υπόλοιπη μέρα κύλισε με εμένα να κοιτάζω έξω από το παράθυρο την πόλη άλλη μία φορά. Νόμιζα πως χθες το βράδυ θα ήταν η τελευταία φορά, αλλά δεν περίμενα οι διαδικασίες του θάνατου του συναδέλφου μου να διαρκέσουν τόσο λίγο. Βλέποντας το φεγγάρι να αστράφτει στον ουρανό, είπα στον εαυτό μου να μη βιαστεί. Έπρεπε να φτάσουν μεσάνυχτα ώστε όλοι οι κάτοικοι του αρχοντικού να κοιμούνται σίγουρα. Αυτές ήταν τουλάχιστον οι οδηγίες του στρατηγού. Η ώρα μέχρι τις έντεκα πέρασε παρατηρώντας τους δείκτες του ρολογιού, που συμβάδιζαν κάθε ώρα για μερικά λεπτά, αλλά η ταχύτητα του μεγάλου δείκτη τους χώριζε μέχρι την επόμενη ώρα. Ο λεπτός δείκτης, από την άλλη, ήταν ο πιο γρήγορος από τους τρεις. Δεν καθόταν για πολύ ώρα μαζί με τους άλλους δύο, αλλά κάθε λεπτό τους υπενθύμιζε την ύπαρξη του περνώντας από δίπλα τους για μία μικρή επίσκεψη. Νιώθω πως αυτό είναι συμβολισμός για κάτι που δεν μου έρχεται τώρα στο μυαλό.

Οι τρεις δείκτες συγχρονίστηκαν για μία τελευταία φορά μπροστά μου. Το ρολόι έδειχνε πλέον δώδεκα. Η τελευταία φάση του “Σχεδίου ” είχε ξεκινήσει. Το να ανοίξω την είσοδο του αρχοντικού στον Ράινερ δεν ήταν και τόσο δύσκολο, ο μπάτλερ πάντα άφηνε το κλειδί σε ένα τραπεζάκι κοντά στην πόρτα, οπότε δεν είχα παρά μόνο να ξεκλειδώσω την είσοδο στον αντισυνταγματάρχη, το οποίο και έκανα.

«Καλώς σας βρήκαμε» μου είπε ο Ράινερ τη στιγμή που άνοιξα την πόρτα και μου πέταξε την καραμπίνα μου. «Λοιπόν, εγώ θα πάρω το ισόγειο, εσύ τον πρώτο όροφο» είπε και όπλισε το πιστόλι του κάνοντας ένα βήμα μέσα στο αρχοντικό.

«Μίσο…» τον κράτησα πριν προχωρήσει παραπέρα. «Δεν υπάρχουν υπνοδωμάτια στο ισόγειο»

«Ωραία» είπε ο αντισυνταγματάρχης και κοίταξε τις σκάλες. «Ο διάδρομος πάνω οδηγεί σε δωμάτια και από τις δύο μεριές;»

«Ν… ναι…»

«Ωραία, θα πάρω εγώ την αριστερή και εσύ την δεξιά» είπε ο αντισυνταγματάρχης και άρχισε να ανεβαίνει τις σκάλες προσπαθώντας να κάνει τον ελάχιστο δυνατό θόρυβο.

Εγώ τον ακολούθησα προφανώς, μιμούμενος τα αργά βήματα του. Βλέποντας τον να στρίβει αριστερά στην κορυφή της σκάλας κατάλαβα κάτι που δεν σκέφτηκα νωρίτερα. Μου είχε ανατεθεί να σκοτώσω τον Μάρκο, τον άνθρωπο με τον οποίο είχα περάσει τις τελευταίες μέρες και ο οποίος ήταν η μόνη πραγματική συντροφιά που είχα τα τελευταία… δεν ξέρω πόσα χρόνια. Φτάνοντας στο δωμάτιο του άνοιξα την πόρτα και τον είδα να στέκεται μπροστά στο αριστερό παράθυρο του δωματίου του, μακριά από το γραφείο του, πάνω στο οποίο είχε ανάψει ένα κεράκι.

«Ξέρεις…» μου είπε «Το όνομα σου δεν έχει νόημα αν δεν απαντάς σε αυτό… Μάρκο»

Σήκωσα το όπλο μου και τον πλησίασα, κλείνοντας ασυναίσθητα την πόρτα πίσω μου.

«Σου είπα πως είχα ένα προαίσθημα ότι θα πεθάνω σύντομα, όμως αυτό…» γύρισε το κεφάλι του ελαφρώς πλάγια ώστε να φανεί το δεξί μέρος του προσώπου του που φωτιζόταν αχνά από το κερί.

«Τι κάνεις τέτοια ώρα ξύπνιος;» τον ρώτησα με το όπλο μου στοχευμένο πάνω του.

«Αποδέχομαι τη μοίρα μου, τι άλλο μπορώ να κάνω; Δεν μπορώ να κατηγορήσω κανέναν για αυτό που γίνεται. Πώς θα μπορούσα να κατηγορήσω τον άνεμο για την αταξία που έφερε, αν εγώ ήμουν αυτός που άνοιξε το παράθυρο;»

«Τι λες;» μπερδεύτηκα με τα λόγια του, δεν χαλάρωσα όμως καθόλου τον στόχο μου.

«Χθες το βράδυ σας άκουσα Μάρκο. Επέστρεψα στο δωμάτιο μου μετά το βραδινό και είδα το σακούλι σου να βρίσκεται ακόμη στο δωμάτιο μου» είπε και μου ξέφυγε μία βαριά εισπνοή καθώς αντιλήφθηκα τι είχε γίνει. «Σας άκουσα Μάρκο που ο παππούς σου, που δεν ξέρω καν αν είναι παππούς σου πλέον, κατάπινε τα φάρμακα ένα ένα χωρίς σταματημό, και εσύ του έδινες ακόμη ένα, και ακόμη ένα, μέχρις ότου να χτυπήσει το σώμα του στο πάτωμα και συ στη συνέχεια να προσπαθείς να τον σηκώσεις για να τον αφήσεις πίσω στο κρεβάτι ώστε να μην υποπτευθούμε τίποτα. Όταν άκουσα τον γέρο σου να λέει “Νομίζω ήρθε η ώρα”, νόμιζα ότι θα μας σκοτώσετε επιτόπου, δεν θα μπορούσα ποτέ να προβλέψω αυτό που θα άκουγα στη συνέχεια»

«Μάρκο…»

«Και όλο το υπόλοιπο βράδυ…» άρχισε να περπατάει πέρα δώθε εκφράζοντας τον εαυτό του περισσότερο με τις κινήσεις του σώματος του. «Το υπόλοιπο βράδυ δεν μπορούσα να κοιμηθώ. Σκεφτόμουν: γιατί να αυτοκτονήσει ο γέρος και εσύ να προσπαθείς να το κρύψεις, κάτι άλλο θα κρύβεται από πίσω. Όλη μέρα προσπαθούσα να σκεφτώ σενάρια, αλλά τίποτα δεν έβγαζε νόημα, μέχρις που σε είδα να κάνεις νεύμα από απόσταση στον μεσήλικα που κληρονόμησε το σπίτι του άντρα που είχε πεθάνει αυτές τις μέρες. Τότε κατάλαβα πως ήταν αργά. Ήμασταν ήδη νεκροί. Για αυτό αποδέχομαι την μοίρα μου»

«Μάρκο, δεν φταις εσύ…»

«Είχα αρκετές ώρες να σκεφτώ Μάρκο. Και έχω να πω πως βρήκα την απάντηση σε όλα αυτά που γίνονται»

«Τι εννοείς;» έφτιαξα τον στόχο μου που είχε πέσει ελάχιστα.

«Να σε ρωτήσω Μάρκο» με πλησίασε αργά σταματώντας δύο μέτρα μπροστά μου. «Πιστεύεις ότι οι άνθρωποι παίρνουν αποφάσεις;»

«Προφανώς» είπα και αυτός γέλασε.

«Όχι… Οι άνθρωποι νομίζουν ότι παίρνουν αποφάσεις. Νομίζουν… Νομίζουν ότι θα στρίψουν δεξιά ή αριστερά… αλλά δεν έφτιαξαν αυτοί τους δρόμους. Οι μεγάλες αποφάσεις έχουν παρθεί για αυτούς εδώ πολύ καιρό» 

Επικράτησε μία στιγμή σιωπής μετά τις λέξεις του Μάρκο.

«Οι μεγάλες αποφάσεις…» συνέχισε «…είναι αυτές που κρίνουν την ιστορία. Μας δίνονται λίγες τέτοιες στη ζωή μας. Εσύ, Μάρκο, είσαι η μία μεγάλη απόφαση που θα έκανα ξανά και ξανά, γιατί για να αποκαλέσω κάποιον φίλο μου, αυτός πρέπει να είναι ίσος μου από κάθε οπτική. Εμείς… δεν ξέρω αν το νιώθεις και συ, αλλά είναι σαν να είμαστε πλευρές του ίδιου νομίσματος»

«Φύγε» του είπα και άφησα ένα δάκρυ που κρατούσα πίσω να κυλήσει.

«Τι;»

«Φύγε από δω…» είπα και άκουσα έναν πυροβολισμό από τα διπλανά δωμάτια. «…σε παρακαλώ…»

«Μην κλαις Μάρκο. Εξάλλου, η προδοσία δεν είναι αυτή που θέτει σε κίνηση ολόκληρη την ιστορία. Ήταν όντως ο Ιούδας κακός ή μήπως ήταν ο άνθρωπος που αναγκάστηκε να παίξει τον πιο άσχημο ρόλο στην πιο σημαντική ιστορία που ειπώθηκε ποτέ;»

«Τι λες πάλι;» τον ρώτησα, προσπαθώντας αμήχανα να κρατήσω τα υπόλοιπα μου δάκρυα πίσω πριν έρθει στο δωμάτιο ο Ράινερ.

«Ο Ιούδας είναι η απόδειξη ότι κάποιοι άνθρωποι είναι πρόθυμοι να αναλάβουν τον ρόλο που κανείς δεν είναι πρόθυμος να αναλάβει, φίλε μου»

«Σε παρακαλώ…» τα δάκρυα μου άρχισαν να κυλούν από τον λαιμό μου, στους ώμους μου, μέχρι το πάτωμα. «Πήδα από το παράθυρο… φύγε…»

«Γιατί να κλαις για το παρελθόν που έχει ήδη περάσει, ενώ μπορείς να χαρείς για το μέλλον μπροστά σου που έρχεται. Θα περάσει ραγδαία, μη τη χάσεις» μου έκλεισε το μάτι και άκουσα την φωνή της κυρίας από τον διάδρομο γεμάτη φόβο πριν ηχήσει ο δεύτερος πυροβολισμός.

«Φύγε» τον σημάδεψα με το όπλο μου φωνάζοντας του για να αποχωρήσει γρήγορα καθώς δεν είχε πολύ χρόνο.

«Μάρκο» άκουσα τη φωνή του Ράινερ από τον διάδρομο.

«Ειρωνικό, ε;» σχολίασε ο συνονόματος μου «…μας ψάχνει»

Μετά από λίγα δευτερόλεπτα, η πόρτα του δωματίου άνοιξε πίσω μου και ο Ράινερ είδε την κατάσταση που επικρατούσε.

«Και νόμιζα ότι εγώ ήμουν αργός που έκανα προσευχή πριν τους σκοτώσω» σχολίασε και πυροβόλησε τον συνονόματο μου στο νεφρό.

«Όχι» φώναξα πετώντας το όπλο μου κάτω και τρέχοντας να πλησιάσω τον Μάρκο.

«Το γεγονός ότι μετανιώνεις για τις πράξεις σου…» είπε καθώς τον σήκωσα στα χέρια μου «…είναι απόδειξη ότι έχεις ωριμάσει» έφτυσε λίγο αίμα. «Το γεγονός ότι εύχεσαι να τα είχες πάει καλύτερα σημαίνει ότι είσαι ήδη κάποιος που θα το έκανε. Και ξέρω κατά βάθος… πως έχεις το ίδιο όνειρο με έμενα… να είσαι ελεύθερος» χαμογέλασε και χαλάρωσε το σώμα του αφήνοντας μου τις τελευταίες του λέξεις.

«Όχι… όχι, όχι, όχι, περίμενε… περίμενε… Μπορώ να το φτιάξω… Μπορώ… άσε… άσε με να προσπαθήσω…» είπα μάταια προσπαθώντας να βρω στο δωμάτιο τρόπο για να σώσω τον συνάνθρωπο μου. «Γιατί φοβάμαι ακόμη πως θα σε χάσω, ενώ έχεις ήδη φύγει;» είπα στον εαυτό μου κοιτώντας τα τώρα κλειστά άψυχα μάτια του Μάρκο. Το μίσος μου για τους αριστοκράτες ήταν αυτό μου με ώθησε να προχωρήσω μέχρι αυτό το σημείο, αλλά κάπου στη διαδρομή αυτό χάθηκε και τώρα είμαι απλά… κενός.

Έψαχνα ασυναίσθητα τόσο καιρό να βρω λάθη στο “Σχέδιο”, αλλά η πρώτη μου σκέψη που είχα κάνει στη συνάντηση μας με τον στρατηγό κατέληξε να είναι και το μεγαλύτερο ελάττωμα. Η υποκρισία μου είχε αναβαθμιστεί σε κάτι παραπάνω. Η ψεύτικη περσόνα μου είχε δεθεί με τον συνομήλικο της τόσο σκληρά που ο δεσμός είχε ξεχειλίσει και στον πραγματικό μου εαυτό. Ήμασταν συμπληρωματικές οπτικές με όμοιο σκεπτικό. Δεν είχα επιλογή, το “Σχέδιο” θα ολοκληρωνόταν είτε με εμένα είτε με τον αντισυνταγματάρχη να παίρνει το τελευταίο αίμα. Ίσως τελικά εμείς ήμασταν οι κακοί από την αρχή.

Άκουσα τον αντισυνταγματάρχη να με πλησιάζει και να σηκώνει το όπλο του, ακουμπώντας το στη ράχη του κεφαλιού μου.

«Αξιολύπητο» ήταν οι τελευταίες του λέξεις πριν τραβήξει για μία τελευταία φορά την σκανδάλη.



\backmatter


%\chapter{Επίλογος}

%Το τέλος.
% ΤΟΔΟ: δεν ξέρω τι να το κάνω αυτό

\end{document}
